\documentclass[11pt]{article}
\usepackage[margin=1in]{geometry}
\usepackage[T1]{fontenc}
\usepackage{lmodern}
\usepackage{amsmath,amssymb,amsthm}
\usepackage{booktabs,array,longtable,tabularx}
\usepackage{microtype,xcolor}
\usepackage[numbers,sort&compress]{natbib}
\usepackage[colorlinks=true,linkcolor=blue!50!black,citecolor=green!40!black,urlcolor=blue!50!black]{hyperref}
\usepackage{url}
\emergencystretch=3em
\setcounter{tocdepth}{2}
\newtheorem{theorem}{Theorem}[section]
\newtheorem{proposition}[theorem]{Proposition}
\newtheorem{lemma}[theorem]{Lemma}
\newtheorem{corollary}[theorem]{Corollary}
\theoremstyle{definition}
\newtheorem{definition}[theorem]{Definition}
\newtheorem{assumption}[theorem]{Assumption}
\newtheorem{example}[theorem]{Example}
\theoremstyle{remark}
\newtheorem{remark}[theorem]{Remark}
\title{Group-Valued Links and Inference in\\Gauge-Covariant Gaussian Attention}
\author{Robert C. Dennis}
\date{Unified theoretical and implementation companion, October 3, 2026}
\providecommand{\GL}{\mathrm{GL}}
\providecommand{\gl}{\mathfrak{gl}}
\providecommand{\tr}{\operatorname{tr}}
\providecommand{\im}{\operatorname{im}}
\providecommand{\ad}{\operatorname{ad}}
\providecommand{\cfg}[1]{{\small\texttt{#1}}}
\providecommand{\KL}{D_{\mathrm{KL}}}
\providecommand{\softmax}{\operatorname{softmax}}
\providecommand{\Tr}{\operatorname{Tr}}
\providecommand{\diag}{\operatorname{diag}}
\providecommand{\E}{\mathbb{E}}
\providecommand{\N}{\mathcal{N}}
\providecommand{\Cov}{\operatorname{Cov}}
\providecommand{\Hent}{\mathcal{H}}
\providecommand{\F}{\mathcal{F}}
\providecommand{\R}{\mathbb{R}}
\providecommand{\C}{\mathbb{C}}
\providecommand{\g}{\mathfrak{g}}
\providecommand{\Ad}{\mathrm{Ad}}
\providecommand{\Log}{\operatorname{Log}}
\providecommand{\Imag}{\operatorname{Im}}
\providecommand{\sym}{\operatorname{sym}}
\providecommand{\asym}{\operatorname{asym}}
\providecommand{\stab}{\mathfrak{stab}}
\providecommand{\Stab}{\mathrm{Stab}}
\providecommand{\spec}{\operatorname{spec}}
\providecommand{\range}{\operatorname{range}}
\providecommand{\vect}{\operatorname{vec}}
\providecommand{\Normal}{\mathcal{N}}
\providecommand{\ip}[2]{\langle #1, #2\rangle}
\providecommand{\Nsc}{N_{\mathrm{sc}}}

\begin{document}
\maketitle
\begin{abstract}
A gauge-covariant attention model must specify how sources are selected, what a communication link carries, how a receiver changes its belief, and how its parameters are learned. This companion unifies six theoretical and implementation notes around exponential group-valued links $\Theta_{ij}=U_i e^{X_{ij}}U_j^{-1}$, the author's primary construction for future integration into the GL(K) attention manuscript. Low-rank exponential identities relate selective channels to the current projected architecture, while inverse-edge triangle products distinguish flat reference identifications from nonflat communication. Gaussian prediction, frozen receiver energies, and learned-template inference provide separate statistical realizations, with their covariance updates and approximation boundaries stated explicitly. Frame optimization is developed through both a Frobenius pullback and a belief-family Fisher metric, including stabilizers, completions, chart conversion, and group retraction. The resulting account preserves detailed derivations, reconciles implementation descriptions with a dated VFE 3.0 snapshot, and explains their relationship to the main attention manuscript and its gradient, supplementary, and nonlinear companions. It supplies an integration map rather than a revision of those documents. Historical measurements and focused CPU checks do not establish a training advantage or a common population evidence bound for the entire composed architecture.
\end{abstract}
\tableofcontents
\clearpage

\section{Scope and continuity with the GL(K) construction}
\label{unified:scope}

This manuscript unifies six developments of gauge-covariant Gaussian attention: frame optimization, the group-valued origin and curvature of communication links, predictive and gradient-based receiver updates, and routing to learned Gaussian templates. Its primary geometric construction is the exponential interaction link
\begin{equation}
\Theta_{ij}=U_i e^{X_{ij}}U_j^{-1}, \qquad X_{ij}\in\mathfrak{gl}(K,\mathbb R).
\label{unified:primary_link}
\end{equation}
This is Option A of the link-origin note, selected by the author as the direction for future integration into \emph{Gauge-Covariant Gaussian Attention through Projected Matrix Links} \cite{unified:attention}. The selection fixes the organization of the theory. It does not assert that the current implementation learns an independent exponential generator on every edge or establish a performance advantage. Compression and the routed head sum remain supporting constructions that identify what the present architecture recovers and what would change.

The six source notes remain separate historical documents. Here their assumptions are made explicit, their derivations are placed in one dependency order, and their implementation descriptions are reconciled with VFE 3.0 as inspected on October 3, 2026. Results for general matrix groups or unrestricted Gaussian covariances retain that scope even when the executable family is narrower. Archived checks and measurements remain dated evidence. This is a standalone companion prepared for later integration; the existing GL(K) manuscripts are not revised by this consolidation.

\subsection{Objects, transformations, and distinct roles}
\label{unified:objects}

At token position $i$, a state $q_i=\mathcal N(\mu_i,\Sigma_i)$ is expressed on a fiber $E_i\cong\mathbb R^K$. An invertible frame $U_i:\mathsf V\to E_i$ identifies a reference space with that fiber. Its reference moments are
\begin{equation}
\xi_i=U_i^{-1}\mu_i, \qquad C_i=U_i^{-1}\Sigma_iU_i^{-\top}.
\label{unified:reference_moments}
\end{equation}
Under a passive local re-expression $(\mu_i,\Sigma_i,U_i)\mapsto(g_i\mu_i,g_i\Sigma_i g_i^\top,g_iU_i)$, the reference moments and reference generator $X_{ij}$ remain fixed. The link obeys $\Theta_{ij}\mapsto g_i\Theta_{ij}g_j^{-1}$. A common reference-basis change is a different action: $U_i\mapsto U_i a$, $\xi_i\mapsto a^{-1}\xi_i$, $C_i\mapsto a^{-1}C_i a^{-\top}$, and $X_{ij}\mapsto a^{-1}X_{ij}a$. Reference metrics and interface maps must transform consistently when that second covariance is claimed.

Geometry asks which directions a link reads, writes, preserves, or composes. Source inference asks which label is selected by a specified record or variational objective. Receiver inference asks whether routed content is averaged predictively or used as evidence in a receiver energy. Parameter learning asks how priors, frames, interfaces, and templates change across training. These roles share variables, but a result about one does not answer the others: endpoint covariance selects neither a likelihood nor an optimizer, and a natural gradient requires a specified metric and objective.

The implemented comparison maps $S_{Q,h},S_{K,h}$ measure Gaussian laws, while $T_h:\mathsf V\to\mathbb R^{d_h}$ sends content and $O_h:\mathbb R^{d_h}\to\mathsf V$ receives it. Their matrix link is $L_{ij}^{(h)}=U_iO_hT_hU_j^{-1}$. With $X_{ij}=sO_hT_h$, this is the derivative at $s=0$ of the exponential link; its exact finite correction uses $\varphi_1(Z)=\sum_{n\ge1}Z^{n-1}/n!$. The next section derives this relation and the assumptions needed for selective aggregation. A real exponential is invertible with positive determinant, but does not parametrize every real invertible matrix.

\subsection{What the main attention manuscript already supplies}
\label{unified:main_bridge}

The main GL(K) manuscript \cite{unified:attention}, in its sections on framed states, the reference-coordinate normal form, and the probabilistic status of routing and content, supplies the common state and interface construction. At fixed input states, nonnegative comparison costs and positive source priors give
\begin{equation}
\beta_{ij}^{(h)}=\frac{\pi_{ij}^{(h)}\exp[-D_{ij}^{(h)}/\tau_h]}{\sum_{k\in S_i}\pi_{ik}^{(h)}\exp[-D_{ik}^{(h)}/\tau_h]}, \qquad \tau_h>0.
\label{unified:source_row}
\end{equation}
This row uniquely minimizes $\sum_j\beta_{ij}^{(h)}D_{ij}^{(h)}+\tau_h\mathrm{KL}(\beta_i^{(h)}\Vert\pi_i^{(h)})$. It is also the exact posterior after observing a declared binary compatibility record with success probability $\exp(-D_{ij}^{(h)}/\tau_h)$, the optimality-record construction of control as inference \cite{unified:levine2018}. Exactness is conditional on that record model. It neither makes the record an observed training datum nor derives the cost from symmetry. At unit temperature a separate fixed-query Gaussian-mixture variational model gives the same row. The probabilistic readings should remain distinct.

The main manuscript already derives the first two moments of a normalized predictive mixture-of-channels kernel, including shared-source, cross-head, and residual self-edge terms. Its evaluated configuration used covariance passthrough. The predictive section below connects that existing covariance derivation to the current moment-matched implementation and states the projection imposed by diagonal storage. The main manuscript's static-covariance normal form continues to describe its reported configuration; a later revision must allow $C_i$ to vary with depth when moment matching or receiver inference updates it.

The main manuscript permits singular individual channels and an invertible routed head sum. Option A introduces a primary invertible interaction object before selective channels are extracted. On causal support $j\le i$, there is no closed directed walk through distinct positions. Invertible links nevertheless permit the inverse-edge triangle $\Theta_{ik}^{-1}\Theta_{ij}\Theta_{jk}$ for $k<j<i$. This observable is compatible with causal message passing and may be nontrivial even for composite links. It supplies neither independent gauge-field degrees of freedom nor an action or continuum limit by itself.

\subsection{Receiver dynamics, templates, and the correlated extension}
\label{unified:companions}

The gradient companion \cite{unified:gradients} distinguishes entropy-regularized routing from entropy-suppressed surrogates and develops Gaussian and frame derivatives. The receiver analysis here specializes that support to projected-gradient attention. It separates ambient alignment of full-dimensional receiving laws from channel likelihoods, identifies their different entropy coefficients, and states the frozen-objective conditions for exact natural-parameter interpolation. A frozen-weight expected-log-likelihood update is a declared variational or generalized-posterior operation. It is not automatically the marginal of a model that first conditions separately on each source label and then averages the resulting posteriors.

Template routing addresses the block MLP that the main manuscript treats as an architectural choice. A mixture prior derives responsibilities and supports a tied posterior update or a separately declared value readout. Reading the incoming Gaussian as an evidence message instead gives a structured mixture posterior. These constructions have different sufficient statistics and covariance updates. Their existence does not show that an arbitrary MLP is a natural-gradient field. A local-coordinate template block also does not acquire full passive GL(K) covariance from the attention sublayer; the reference-coordinate construction supplies the corresponding covariant alternative.

The supplement \cite{unified:supplement} and nonlinear extension \cite{unified:extension} place local mechanisms beside a normalized population model. On a regular observation slice, a fixed joint and a recognition law satisfying the required absolute continuity and finite-expectation conditions obey
\begin{equation}
-\mathbb E_Q\log\frac{p_\theta(z,o)}{q(z)}=\mathrm{KL}\bigl(Q\Vert P_\theta(\cdot\mid o)\bigr)-\log p_\theta(o).
\label{unified:elbo}
\end{equation}
The evidence identity is exact for the declared joint and recognition family \cite{unified:blei2017}. It does not identify arbitrary layer composition with coordinate minimization. The extension supplies a population construction with source labels, latent reliability, correlated Gaussian recognition, and reference-relative information control. Its source weights use expected record residuals, including cross-agent covariance, rather than generically the marginal Gaussian KL of Eq.~\eqref{unified:source_row}. The constructions share a categorical variational mechanism but use different energies.

For a joint state with covariance blocks $\Gamma_{ij}$, the residual $x_i-L_{ij}x_j$ has covariance
\begin{equation}
\Gamma_{ii}+L_{ij}\Gamma_{jj}L_{ij}^\top-\Gamma_{ij}L_{ij}^\top-L_{ij}\Gamma_{ji}.
\label{unified:population_residual}
\end{equation}
The cross terms vanish under a product-state assumption, not from knowledge of the marginals. Regression links below use such cross covariances; current per-token covariance tensors do not provide them. Repeating inference under fixed factors has a frozen-objective interpretation, whereas untied depth-dependent maps need a larger model before one global descent statement is available. A causal mask also does not remove source-side derivatives of a population energy. Prefix-only state construction, shared across candidate next tokens, is what gives the implemented categorical predictor its normalized autoregressive law.

\subsection{Evidence and integration scope}
\label{unified:evidence_scope}

The mathematical statements below are supported by their derivations and declared hypotheses. Implementation descriptions distinguish active settings, reachable alternatives, archived behavior, and proposals. Small-dimensional calculations check identities and execution details; they do not determine the theoretical direction or establish large-network behavior. No GPU workload, training comparison, or large-population experiment is part of this consolidation. The program is to retain Option A as the primary group-valued geometry, specify each extracted channel's statistical role, and connect receiver and template updates to a chosen model without silently changing its conditioning or recognition family.
\section{Exponential interaction links and selective communication}\label{unified:links}
\subsection{Communication links and the geometric construction}
\label{links:origin:sec:setting}

Agent $i$ carries a Gaussian belief $q_i=\mathcal N(\mu_i,\Sigma_i)$ on $\mathbb R^K$ and an invertible frame $U_i$, with reference moments $\xi_i=U_i^{-1}\mu_i$ and $C_i=U_i^{-1}\Sigma_iU_i^{-\top}$. A passive frame change $g_i\in\GL(K)$ acts by $(\mu_i,\Sigma_i,U_i)\mapsto(g_i\mu_i,g_i\Sigma_ig_i^\top,g_iU_i)$ and fixes $(\xi_i,C_i)$. The flat transport is $\Omega_{ij}=U_iU_j^{-1}$. Head $h$ of a layer owns a sending map $T_h:\mathbb R^K\to\mathbb R^d$ and a receiving map $O_h:\mathbb R^d\to\mathbb R^K$, and the code requires $K=Hd$. The matrix link is
\begin{equation}
L_{ij}^{(h)}=U_iO_hT_hU_j^{-1},\qquad L_{ij}^{(h)}\mapsto g_iL_{ij}^{(h)}g_j^{-1},
\label{links:origin:eq:link}
\end{equation}
and the projected predictive mean is $\mu_i^+=r\mu_i+\sum_{h,j}\beta_{ij}^{(h)}L_{ij}^{(h)}\mu_j$ with $r\in\{0,1\}$ set by the attention residual \cite{links:dennis2026projpred}. Stacking the heads gives square matrices $O=[O_1\cdots O_H]$ and $T=[T_1;\dots;T_H]$, both $K\times K$, and the head-diagonal routing matrix $B_{ij}=\bigoplus_h\beta_{ij}^{(h)}I_d$.

A group-valued link is a field $\Theta_{ij}\in\GL(K)$ with the same law $\Theta_{ij}\mapsto g_i\Theta_{ij}g_j^{-1}$, which is the lattice gauge field of Wilson \cite{links:wilson1974}. The primary construction for future integration is the exponential link $\Theta_{ij}=U_i\exp(X_{ij})U_j^{-1}$ developed below. A low-rank content map is its infinitesimal channel term, and on an appropriate logarithm domain its finite channel correction has an exact affine representation. Compression of a primitive link and the presently implemented head sum provide complementary realizations and implementation comparisons. They do not change the choice of the exponential construction as the geometric starting point.

Two facts constrain any answer. First, when $d<K$, $L_{ij}^{(h)}$ has rank at most $d<K$, so for any $M$ the product $L^{(h)}M$ is singular, and no factorization $\Theta=L^{(h)}N$ of a group element exists. The relation between $L$ and $\Theta$ must be additive, through a block of $\Theta$, or infinitesimal, through the Lie algebra. Second, on causal support every source satisfies $j\le i$, so a closed directed walk can only be a self-loop, as the main attention manuscript proves for directed causal walks. The triangle constructions below explicitly complete this graph with inverse reverse edges. They do not claim that a forward causal computation contains a directed cycle. An existing invertible composite edge can be given such a reverse edge without adding an independently learned connection; a separately computed reverse attention map need not equal that inverse. Singular links admit closed matrix products on a declared extended graph, but do not supply inverse reverse edges. This distinction preserves the main manuscript's rank and composition algebra and its separation of $\Omega$ from communication links \cite{unified:attention,links:wilson1974,links:kogut1979}.


The connection discussed here belongs to the declared communication graph over labeled token fibers. It is not the curvature of the smooth contextual bundle in the main manuscript. Its token construction is a constant-context pullback with a common reference trivialization. Passive covariance follows from the endpoint frames and does not constrain the reference maps to be intertwiners. Under a common reference change $a$, every reference link is conjugated by $a^{-1}$; the shared family of reference links therefore retains simultaneous-similarity data. Flatness of the link field is also weaker than alignment with the specific frame transport $\Omega_{ij}$: a distinct coboundary can be flat while changing the communication law.
\subsection{The primary construction: exponential group-valued links}
\label{links:origin:sec:exp}

\begin{definition}[Exponential link]
For an algebra field $X_{ij}\in\gl(K)$ written in reference coordinates, set
\begin{equation}
\Theta_{ij}=U_i\exp(X_{ij})U_j^{-1}.
\label{links:origin:eq:exp_link}
\end{equation}
\end{definition}

The law $\Theta_{ij}\mapsto g_i\Theta_{ij}g_j^{-1}$ holds for every $X_{ij}$, because $X_{ij}$ is expressed in reference coordinates and the frames carry all of the passive action. The identity $\det\exp(X)=e^{\tr X}$ makes $\Theta_{ij}$ invertible for every algebra element, singular ones included. The choice whose first-order channel term reproduces the current links is the rank-at-most-$d$ algebra element $X=O_hT_h$ for one head, or $X=\sum_hO_hT_h=OT$ for all heads with equal weight.

\begin{lemma}[Exponential of a low-rank element]
\label{links:origin:lem:lowrank_exp}
Let $P$ be $K\times r$ and $Q$ be $r\times K$, and let $\varphi_1(Z)=\sum_{n\ge1}Z^{n-1}/n!$, which equals $(e^Z-I)Z^{-1}$ when $Z$ is invertible. Then
\begin{equation}
\exp(PQ)=I_K+P\varphi_1(QP)Q.
\label{links:origin:eq:lowrank_exp}
\end{equation}
\end{lemma}

\begin{proof}
For $n\ge1$, $(PQ)^n=P(QP)^{n-1}Q$ by regrouping the product. Summing the exponential series term by term gives $I_K+P\big(\sum_{n\ge1}(QP)^{n-1}/n!\big)Q$, and the inner series converges for every $r\times r$ matrix.
\end{proof}

\begin{proposition}[Flat transport plus a channel term]
\label{links:origin:prop:split}
With $X_{ij}=O_hT_h$,
\begin{equation}
\Theta_{ij}=\Omega_{ij}+\widetilde L_{ij}^{(h)},\qquad
\widetilde L_{ij}^{(h)}=U_iO_h\varphi_1(T_hO_h)T_hU_j^{-1}.
\label{links:origin:eq:split}
\end{equation}
The following hold.
\begin{enumerate}
\item $\exp(O_hT_h)$ acts as the identity on $\ker T_h$, so on the reference subspace that the head does not read, $\Theta_{ij}$ coincides with the flat transport $\Omega_{ij}$.
\item The image of $\exp(O_hT_h)-I$ lies in $\im O_h$, so $\Theta_{ij}$ differs from $\Omega_{ij}$ only in the directions that the head writes.
\item Under a scaling $X_{ij}=sO_hT_h$, $\Theta_{ij}=\Omega_{ij}+sL_{ij}^{(h)}+O(s^2)$, so the matrix link of Eq.~\eqref{links:origin:eq:link} is the first-order term.
\item $\det\Theta_{ij}=\det(U_iU_j^{-1})e^{\tr(T_hO_h)}$.
\end{enumerate}
\end{proposition}

\begin{proof}
Equation~\eqref{links:origin:eq:split} is Lemma~\ref{links:origin:lem:lowrank_exp} with $P=O_h$ and $Q=T_h$, conjugated by the frames. Items 1 and 2 are read off Eq.~\eqref{links:origin:eq:lowrank_exp}: the correction is $O_h(\cdots)T_h$. Item 3 follows from $\varphi_1(sZ)=I+O(s)$. Item 4 is $\det\exp(X)=e^{\tr X}$ with $\tr(O_hT_h)=\tr(T_hO_h)$.
\end{proof}

Proposition~\ref{links:origin:prop:split} gives the original question an exact answer. The remainder $N$ of the question is the identity on the non-communicating subspace $\ker T_h$, and the communicating part is $\widetilde L_{ij}^{(h)}$. That part equals the matrix link to first order and carries the corrections $\varphi_1(T_hO_h)$ at higher order. Nothing in $\Theta_{ij}$ is hidden: every entry is fixed by the frames and by $(O_h,T_h)$.

\paragraph{The two historical routes as orders of one link.}
Contracting Eq.~\eqref{links:origin:eq:split} against routing weights gives the message
\begin{equation}
\sum_j\beta_{ij}\Theta_{ij}(s)\mu_j
=\underbrace{\sum_j\beta_{ij}\Omega_{ij}\mu_j}_{\text{flat transport}}
+s\underbrace{\sum_j\beta_{ij}L_{ij}^{(h)}\mu_j}_{\text{projected link}}+O(s^2),
\quad \Theta_{ij}(s)=U_i e^{sO_hT_h}U_j^{-1}.
\label{links:origin:eq:two_routes}
\end{equation}
The first term is the transported message of the earlier transported route, and the second is the projected predictive message. At fixed routing and source states, the exponential expansion therefore contains both message forms as its zeroth- and first-order terms. This is an observation about the form of the message, not a claim that either retired objective is recovered.

\paragraph{Routing inside the exponent.}
Setting $X_{ij}=OB_{ij}T=\sum_h\beta_{ij}^{(h)}O_hT_h$ instead dresses the algebra element with the routing weights, and Lemma~\ref{links:origin:lem:lowrank_exp} with $P=O$ and $Q=B_{ij}T$ gives
\begin{equation}
\exp(OB_{ij}T)=I+O\varphi_1(B_{ij}TO)B_{ij}T.
\label{links:origin:eq:exp_routed}
\end{equation}
The first-order term is the multi-head operator of Section~\ref{links:origin:sec:headsum}. Heads interact at second order through the cross-head blocks of $TO$. With routing inside the exponent, the weights no longer normalize a sum over sources, so the aggregation over $j$ must be chosen separately. Section~\ref{links:origin:sec:aggregation} resolves this by keeping the routing outside the exponent.

\paragraph{Loops.}
For a triangle with a reversed edge, the loop is $\Theta_{ij}\Theta_{jk}\Theta_{ik}^{-1}=U_ie^{X_{ij}}e^{X_{jk}}e^{-X_{ik}}U_i^{-1}$, and its conjugacy class, for example its trace, is frame invariant. When the algebra elements commute, which happens for square invertible stacked maps with $TO=I_K$, since then $X_{ij}=OB_{ij}O^{-1}$, the loop equals
\begin{equation}
U_iO\exp\big(B_{ij}+B_{jk}-B_{ik}\big)O^{-1}U_i^{-1},
\label{links:origin:eq:exp_loop}
\end{equation}
so the curvature is the per-head additive defect $\beta_{ij}^{(h)}+\beta_{jk}^{(h)}-\beta_{ik}^{(h)}$. In general the loop follows the Baker-Campbell-Hausdorff series, and the commutators $[X_{ij},X_{jk}]$ make it nonabelian.

\subsubsection{Linear and exponential links are two charts of one group}
\label{links:origin:sec:charts}

The phrase ``first-order term'' in Proposition~\ref{links:origin:prop:split} invites the reading that the matrix link is a weak-field truncation of the group link. There is an exact change of coordinates on a specified logarithm domain; no global equivalence of all linear and exponential parameter values is implied. Fix $O\in\mathbb R^{K\times d}$ and $T\in\mathbb R^{d\times K}$ with $TO$ invertible, so that $\mathbb R^K=W\oplus W'$ with $W=\im O$ and $W'=\ker T$, and let $P=O(TO)^{-1}T$ be the projector onto $W$ along $W'$. For $B\in\GL(d)$ define
\begin{equation}
\iota(B)=(I-P)+OB(TO)^{-1}T,
\label{links:origin:eq:embed}
\end{equation}
the operator that fixes $W'$ pointwise and acts on $W$ as $B$ in the basis formed by the columns of $O$.

\begin{proposition}[Two charts of one group]
\label{links:origin:prop:charts}
\begin{enumerate}
\item $\iota:\GL(d)\to\GL(K)$ is an injective homomorphism with $\det\iota(B)=\det B$. Its image $\mathcal G_W$ is the $\GL(d)$ factor of the block-diagonal Levi subgroup $\GL(d)\times\GL(K-d)$ that acts trivially on $W'$.
\item \emph{Linear links.} For $C\in\mathbb R^{d\times d}$, $I+OCT=\iota(I_d+CTO)$, so $\det(I+OCT)=\det(I_d+CTO)$.
\item \emph{Exponential links.} $\exp(OT)=\iota(e^{TO})$.
\item \emph{Exact conversion.} Independently of the fixed splitting in items 1--3, let $O'\in\mathbb R^{K\times d}$ be such that $I_d+TO'$ has no eigenvalue on the closed negative real axis $(-\infty,0]$. Put $Z=\log(I_d+TO')$ (principal logarithm) and $O=O'\varphi_1(Z)^{-1}$. Then $\varphi_1(Z)$ is invertible, $TO=Z$, and
\begin{equation}
\exp(OT)=I_K+O'T.
\label{links:origin:eq:conversion}
\end{equation}
\end{enumerate}
\end{proposition}

\begin{proof}
For item 1, $\iota(B)Oy=OBy$ because $(I-P)O=0$ and $(TO)^{-1}TO=I$, and $\iota(B)w'=w'$ for $w'\in W'$ because $Pw'=0$. An operator on $W\oplus W'$ is determined by its action on the two summands, so $\iota(B)\iota(B')=\iota(BB')$ and $\det\iota(B)=\det B\cdot1$. For item 2, $(I+OCT)Oy=O(I+CTO)y$ and $(I+OCT)w'=w'$. For item 3, $(OT)^nO=O(TO)^n$ gives $\exp(OT)Oy=Oe^{TO}y$, and $\exp(OT)w'=w'$. For item 4, the eigenvalues of $Z$ have imaginary part in $(-\pi,\pi)$, so $\varphi_1(Z)$, whose eigenvalues $(e^z-1)/z$ vanish only at $z\in2\pi i\mathbb Z\setminus\{0\}$, is invertible. Since $e^Z-I=TO'=Z\varphi_1(Z)$, we get $TO=TO'\varphi_1(Z)^{-1}=Z$, and Lemma~\ref{links:origin:lem:lowrank_exp} gives $\exp(OT)=I+O\varphi_1(TO)T=I+O'\varphi_1(Z)^{-1}\varphi_1(Z)T=I+O'T$.
\end{proof}

Item 4 says that the flat transport plus a linear channel term, $\Omega_{ij}+U_iO'TU_j^{-1}$, is exactly an exponential link for every $O'$ in an open set that contains $O'=0$. The affine parametrization covers the invertible locus of the embedded subgroup; the principal-logarithm parametrization covers the open part whose channel block has spectrum off $(-\infty,0]$. On that common domain they describe the same group elements with $B=I+C TO=e^Z$. The expansion $\Theta=\Omega+L+O(X^2)$ of Proposition~\ref{links:origin:prop:split} is the statement that the charts agree to first order at the identity, and with $\widetilde L=O\varphi_1(TO)T=\exp(OT)-I$ in reference coordinates the series $X=\log(I+\widetilde L)=\widetilde L-\tfrac12\widetilde L^2+\cdots$ converts one chart to the other. Thus the affine correction can represent the exponential exactly after reparametrization, while using the same numerical generator in $I+X$ and $e^X$ remains only a first-order approximation. The same holds for the lattice relation $U=e^{aA}$, where the algebra-valued potential $A$ is the logarithmic coordinate, with the difference that here the frames dress the potential, so that $X$ is frame invariant in reference coordinates.

\begin{corollary}[Tiled heads]
\label{links:origin:cor:tiled_chart}
If the stacked $O$ and $T$ are square and invertible, so that $W=\mathbb R^K$, an operator $M=O'T$, representing the unrestricted stacked content map (or a routed map after absorbing a fixed positive routing matrix into $O'$), is a group element, $M\in\GL(K)$ when $\det M\ne0$, and $M=\exp(X)$ for $X=OT$ exactly when $M$ has a real logarithm, with $O=XT^{-1}$. The identity content initialization $M=I$ corresponds to $X=0$ in the logarithmic chart, although a fixed finite generator $OT=I$ would instead give $eI$. The product of $r$ links along a path with a shared generator is $M^r=e^{rX}$, a one-parameter subgroup.
\end{corollary}

The corollary uses $M^r=e^{rX}$ to restate the path-composition remark of the matrix-link construction: the $r$-edge path $U_{i_0}M^rU_{i_r}^{-1}$ is the group element of a one-parameter subgroup, not a separate object.

\paragraph{What the choice of chart does change.}
First, the ambient affine parameter space includes the singular matrices; only its invertible locus is a group chart. A step in $(O',T)$ can carry $\det(I_K+O'T)=\det(I_d+TO')$ through zero (in the setting of Corollary~\ref{links:origin:cor:tiled_chart}, $\det M$), where at least one channel direction is lost by the linear map. The logarithmic chart cannot reach that boundary, because $\det e^Z=e^{\tr Z}>0$. The same argument shows that the exponential chart covers only matrices with a real logarithm: it excludes negative-determinant components, which an optimizer in the affine chart reaches only by passing through $\det=0$, and a real invertible matrix has a real logarithm exactly when, for each negative real eigenvalue and each Jordan block size, the number of those blocks is even \cite{frame:Culver1966}. Second, the two charts give different optimization geometries, since the code's content M-step takes AdamW updates on the factors $(O',T)$, rather than multiplicative updates of a group element. Such factor updates induce changes in the composite matrix; they are not identical to Euclidean gradient descent on the composite itself. Third, they differ at order one for finite generators: at $X=I$, $e^X=eI$ while $I+X=2I$. Fourth, these statements concern an unrouted operator with a single generator. The routed head sum $\sum_h\beta^{(h)}O_hT_h$ weights heads differently and does not lie in the one-generator family, and arbitrary mixing with the hidden complement remains outside $\mathcal G_W$. Edge-dependent generators may stay within that subgroup or leave it, depending on their specified support.

\subsection{Aggregation, the residual, and the nonlinearity under Option A}
\label{links:origin:sec:aggregation}

The group law by itself does not choose an aggregation kernel. The following construction retains a routing-independent generator, attaches categorical source labels to the sent channel blocks, and preserves the receiver's unsent complement. These are explicit modeling choices connecting the exponential link to the predictive route.

\subsubsection{The complement split}

Stack the heads of a layer into $O$ of size $K\times Hd$ and $T$ of size $Hd\times K$, allow $Hd\le K$, and assume $TO$ is invertible. Let $P=O(TO)^{-1}T$, the projector onto $W=\im O$ along $W'=\ker T$.

\begin{lemma}[Complement split]
\label{links:origin:lem:complement}
\begin{equation}
\exp(OT)=(I-P)+OET,\qquad E=e^{TO}(TO)^{-1}.
\label{links:origin:eq:complement}
\end{equation}
\end{lemma}

\begin{proof}
By Lemma~\ref{links:origin:lem:lowrank_exp}, $\exp(OT)=I+O\varphi_1(TO)T=(I-P)+O\big[(TO)^{-1}+\varphi_1(TO)\big]T$, and $(TO)^{-1}+\varphi_1(TO)=(TO)^{-1}+(e^{TO}-I)(TO)^{-1}=e^{TO}(TO)^{-1}$ because every term is a function of $TO$.
\end{proof}

The first term is the non-communicating part of the question: it is the identity on $W'$ and annihilates $W$. The second term reads the $Hd$ sent coordinates $T\xi$, mixes them through the $Hd\times Hd$ matrix $E$, and writes them into $W$. Equation~\eqref{links:origin:eq:complement} is exact, so the group-valued link separates into a complement transport and a channel transport with no remainder.

\subsubsection{Routing belongs outside the exponent}

A link acts on one source state, whereas a predictive mixture also needs a normalized law over source choices. Placing $B_{ij}$ inside $e^{OB_{ij}T}$ does not itself define that law: neither an unweighted sum of these exponentials nor their product specifies how the different source vectors are selected. Such routed generators are nevertheless admissible if an additional normalized source kernel is declared. For the present construction, categorical routing stays outside the exponential and the shared generator is $X=OT$. This preserves the interpretation of $\beta_{ij}^{(h)}$ as the probability of head $h$ selecting source $j$.

\subsubsection{The aggregation rule}

Following the predictive kernel of the projected predictive route \cite{links:dennis2026projpred}, each head draws its own source label $J^h$ with probability $\beta_{ij}^{(h)}$, independently across heads. In Eq.~\eqref{links:origin:eq:complement} each sent block $T_h\xi$ is read from its head's chosen source, and $E$ mixes the sent blocks into $W$. The complement term needs a carrier of its own, a label $J^0$. The transported reference state is
\begin{equation}
\xi_i'=(I-P)\xi_{J^0}+\sum_hOE_{:,h}T_h\xi_{J^h},
\label{links:origin:eq:transported_state}
\end{equation}
where $E_{:,h}$ is the block column of $E$ that receives head $h$. The expectation over the independent labels is
\begin{equation}
\mathbb E[\xi_i']=\mathbb E[(I-P)\xi_{J^0}]+\sum_h\sum_j\beta_{ij}^{(h)}OE_{:,h}T_h\xi_j.
\label{links:origin:eq:aggregate}
\end{equation}
Each sent block travels with its own head's label, so the cross-head mixing in $E$ does not couple the head routings. Only the complement carrier remains to be chosen. The weights are functions of the states, $\beta_{ij}^{(h)}=\beta_{ij}^{(h)}(\mu,\Sigma)$, through the Gaussian comparison costs of queries and keys. Both projected routes evaluate them at the layer's input states and hold them fixed while the message is formed, which is exact inference in the declared cut model of the projected gradient analysis \cite{links:dennis2026projgrad}. Equation~\eqref{links:origin:eq:aggregate} is therefore the expectation conditional on those states, and the state dependence of $\beta$ does not alter the aggregation rule.

\begin{proposition}[Communication locality with complement persistence]
\label{links:origin:prop:self_complement}
Assume that only sent content may cross between agents and that the update retains one unchanged complement from an existing agent state as in Eq.~\eqref{links:origin:eq:transported_state}. Then its carrier is the receiver, $J^0=i$, and
\begin{equation}
\mu_i^+=U_i(I-P)U_i^{-1}\mu_i+\sum_h\sum_j\beta_{ij}^{(h)}U_iOE_{:,h}T_hU_j^{-1}\mu_j.
\label{links:origin:eq:self_complement}
\end{equation}
\end{proposition}

\begin{proof}
The complement $W'=\ker T$ holds exactly the directions that no head sends. A carrier $J^0=j\ne i$ would place agent $j$'s unsent coordinates in agent $i$'s state, which contradicts the premise, so the receiver is the only admissible carrier within the declared unchanged-complement rule. Without that rule, locality would also permit discarding the complement or applying a receiver-local map; locality alone does not force identity persistence. Equation~\eqref{links:origin:eq:aggregate} with $J^0=i$, conjugated by the frames, gives the formula.
\end{proof}

The alternative, a complement carried by a routed label $J^0\sim\beta^{(0)}$, gives $\sum_j\beta_{ij}^{(0)}(I-P)\xi_j$, which moves the unsent coordinates of other agents. That is a flat-transport message of the kind the earlier transported route used, and it is admissible only if agents are meant to read each other's full state.

\subsubsection{Does the flat term produce a residual?}

Proposition~\ref{links:origin:prop:self_complement} answers the question. Under the communication principle the flat term becomes a residual, but only on the complement: the receiver keeps its own state exactly in the directions that no head writes and takes the transported message in the directions that heads do write. This projected residual $U_i(I-P)U_i^{-1}\mu_i$ differs from the full residual $r\mu_i$ that the code adds under \cfg{attention\_residual}.

\begin{corollary}[Tiled heads leave no residual]
\label{links:origin:cor:tiled}
If the heads tile the fiber, $Hd=K$ with $O$ and $T$ invertible, then $P=I$, the complement is empty, and Eq.~\eqref{links:origin:eq:self_complement} reduces to
\begin{equation}
\mu_i^+=\sum_h\sum_j\beta_{ij}^{(h)}U_iO_h'T_hU_j^{-1}\mu_j,\qquad O_h'=OE_{:,h},
\label{links:origin:eq:tiled}
\end{equation}
the residual-free projected predictive head sum with receiving maps $O_h'$.
\end{corollary}

The code requires $K=Hd$, so with invertible stacked maps its heads tile the fiber, and this declared complement-preserving aggregation has no complement residual. Two consequences follow. First, the $r=1$ residual of the code has no counterpart in Option A under the configuration the code uses, and the receiver's own state reaches $\mu_i^+$ only through the self-weights $\beta_{ii}^{(h)}$. Second, in the tiled case Option A adds no expressive power to the predictive route: its head sum lies in the current function class, with receiving maps of the special form $OE_{:,h}$. We do not claim the converse, that every current set of receiving maps arises this way. Until the generator is made edge dependent, the shared-generator construction changes the group domain and optimization geometry and supplies loop observables, without enlarging the current mean-map family by itself.

A residual of the projected kind requires channels that do not tile, $Hd<K$. Equation~\eqref{links:origin:eq:self_complement} then describes an architecture in which $K-Hd$ coordinates form a residual stream that no head writes and no head reads through $T$. In this reading the residual stream is the complement of the communication channels.

\paragraph{The same answer from Gaussian fusion.}
The projected gradient route reaches the same structure by a different argument \cite{links:dennis2026projgrad}. Its channel-likelihood receiver keeps the prior in directions that no head observes, and in the tiled limit with a flat prior and equal receiving noise it recovers the residual-free predictive head sum. At finite precision, fusion also yields a partial residual in the observed directions. For an isotropic prior of precision $\alpha$ and message evidence of precision $\lambda$ on $W$, the posterior mean is
\begin{equation}
\xi_i^+=(I-P_W)\xi_i+P_W\frac{\alpha\xi_i+\lambda m_i}{\alpha+\lambda},
\label{links:origin:eq:fusion_residual}
\end{equation}
with $P_W$ the orthogonal projector onto $W$ and $m_i$ the message. The residual coefficient is the prior's share of precision, $\alpha/(\alpha+\lambda)$, in the observed directions and one in the unobserved directions, and the orthogonal-complement version of Eq.~\eqref{links:origin:eq:self_complement} is the limit $\lambda/\alpha\to\infty$. The exponential split uses an oblique projector in general, so the two projector formulas coincide only when the chosen splitting is orthogonal in the declared Gaussian precision metric. Both derivations give a residual that is full on unobserved directions and partial or zero on observed ones. The additive residual of the code, $\mu_i+m_i$, is neither: it adds the full own state to the message in every direction without normalization. That form is the architectural residual connection used in residual networks \cite{links:he2016resnet}; it is not the posterior mean of the Gaussian fusion model just specified.

\subsubsection{Is \texorpdfstring{$\varphi_1$}{phi} an activation function?}

Not with the routing-independent generator $X=OT$. There $\varphi_1$ acts on the single-head $d\times d$ round-trip matrix $T_hO_h$, or on the stacked $Hd\times Hd$ round-trip matrix for $X=OT$. These are built from weights, and the link is linear in the transported state, and in the tiled case $\varphi_1$ is absorbed into the receiving maps (Corollary~\ref{links:origin:cor:tiled}). The message is still nonlinear in the states, but only through $\beta(\mu,\Sigma)$, the way standard attention is nonlinear through its softmax.

The answer changes once the generator carries the routing. Because $\beta_{ij}^{(h)}=\beta_{ij}^{(h)}(\mu,\Sigma)$, the routed generator $X_{ij}=OB_{ij}T$ depends on the states, and $\varphi_1(B_{ij}TO)$ in Eq.~\eqref{links:origin:eq:exp_routed} becomes a matrix function of them. On the self-link this is a well-defined single-source transformation, so the objection of the previous subsection to routing inside the exponent does not apply there. For tiled heads with $TO=I$,
\begin{equation}
\exp(OB_{ii}T)=\sum_he^{\beta_{ii}^{(h)}(\mu,\Sigma)}O_hT_h,
\label{links:origin:eq:self_gate}
\end{equation}
so each head's block of the receiver state is multiplied by $e^{\beta_{ii}^{(h)}}\in(1,e]$ on positive self support (and equal to $1$ if self support is absent), a gate whose argument is the self-routing weight $\beta_{ii}^{(h)}=\pi_{ii}^{(h)}e^{-D_{ii}^{(h)}/\tau_h}/Z_i^{(h)}$. Here $D_{ii}$ compares the token's own query and key laws, which differ because $S_Q\ne S_K$, and $Z_i$ compares that cost with every competing source. This is a state-dependent nonlinearity permitted by the chosen routed generator; the group structure alone does not select that generator. It differs from a GELU in three ways. The gate reads the state only through the self-routing weight, which compares the token with its sources rather than reading the token's own coordinates. It acts on whole head blocks rather than coordinatewise. And it lies between $1$ and $e$ inclusive, so it can amplify a block by at most a factor of $e$ but can never suppress it or change its sign.

A general state-dependent generator on the self-link, $\Theta_{ii}=U_i\exp(X(\xi_i))U_i^{-1}$, does produce the gated form of modern activations. For a diagonal generator $X(\xi)=\mathrm{diag}(a(\xi))$ the update is $\xi_k\mapsto\xi_ke^{a_k(\xi)}$, a coordinatewise gate. Every positive gate has this form, since $x\sigma(x)=\exp(\log\sigma(x))x$, so the SiLU $x\sigma(x)$ \cite{links:elfwing2018silu} and the GELU $x\Phi(x)$ \cite{links:hendrycks2016gelu} are exponential self-links with generators $\log\sigma$ and $\log\Phi$. In lattice language the attention links are edge variables and a gated MLP is a state-dependent site term.

This is a rewriting, not a derivation. It places the MLP in the same group-valued language as the links but does not explain why the generator should be $\log\sigma$ or $\log\Phi$. A derivation needs a principle that fixes the generator. Template routing supplies one for a normalized gate, replacing independent two-way competitions by one softmax over Gaussian templates derived from a mixture prior \cite{links:dennis2026template}.

\subsection{Supporting Option B: compression of a primitive link}
\label{links:origin:sec:compression}

Here the group-valued link is taken as primary, $\Theta_{ij}=U_iG_{ij}U_j^{-1}$ with $G_{ij}\in\GL(K)$, and the channel maps act as covariant projectors. When $T_hO_h$ is invertible, define the reference projector $P_h=O_h(T_hO_h)^{-1}T_h$, which projects onto $\im O_h$ along $\ker T_h$, and its frame-dressed version $\Pi_{i,h}=U_iP_hU_i^{-1}$, which transforms by conjugation. Every link then decomposes into four covariant blocks,
\begin{equation}
\Theta_{ij}=\Pi_{i}\Theta_{ij}\Pi_{j}+\Pi_{i}\Theta_{ij}(I-\Pi_{j})+(I-\Pi_{i})\Theta_{ij}\Pi_{j}+(I-\Pi_{i})\Theta_{ij}(I-\Pi_{j}),
\label{links:origin:eq:blocks}
\end{equation}
where the first block is the communicating link, the last is the non-communicating remainder $N$ of the question, and the middle two mix them. More generally, distinct receive and send projectors give the oblique compression $\Pi^R_i\Theta_{ij}\Pi^S_j$, and any rank-$d$ map whose range lies in $\im O_h$ and whose kernel contains $\ker T_h$ arises this way from some invertible $G_{ij}$.

\begin{proposition}[Block factorization]
\label{links:origin:prop:levi}
In a reference basis adapted to $\mathbb R^K=W\oplus W'$ with $W=\im O_h$, write $G=\begin{pmatrix}A&B\\C&D\end{pmatrix}$. If $A$ is invertible, then
\begin{equation}
G=\begin{pmatrix}I&0\\CA^{-1}&I\end{pmatrix}
\begin{pmatrix}A&0\\0&S\end{pmatrix}
\begin{pmatrix}I&A^{-1}B\\0&I\end{pmatrix},\qquad S=D-CA^{-1}B,
\label{links:origin:eq:levi}
\end{equation}
and $\det G=\det A\det S$.
\end{proposition}

The middle factor lies in $\GL(d)\times\GL(K-d)$, the Levi factor of the parabolic subgroup that stabilizes $W$ \cite{links:humphreys1975}. Its first block $A$ is the communicating link, and $S$, the Schur complement \cite{links:zhang2005schur}, is the hidden part with the communicating directions eliminated. The outer unipotent factors carry the mixing. This is the precise multiplicative form of $\Theta=L\times N$: the product holds between the outer mixing factors and the block-diagonal pair $(A,S)$, not between $L$ and $N$ directly.

\begin{remark}[Reduction of structure group]
Assigning the subspace $W_i=U_i\im O_h$ to every agent is a section of the associated Grassmannian bundle, and it reduces the structure group from $\GL(K)$ to the parabolic stabilizer. Reductions of this kind are how a Higgs field breaks a gauge symmetry. The analogy is formal; no potential or dynamics for $W_i$ is proposed here.
\end{remark}

\paragraph{What Option B needs and what it cannot see.}
The current model uses shared $O_h,T_h$ with edge-dependent routing coefficients; it does not store a separate primitive group link or its four blocks. The unrouted composite is constant across edges at a fixed layer, while the routed composite is $OB_{ij}T$. A genuine extension needs a model for $G_{ij}$: a learned relative-position element $G_{i-j}$, a state-dependent element $G(\xi_i,\xi_j)$, or a dynamical link with its own energy. Rotary embeddings are the closest precedent in standard practice, a relative-position orthogonal link inserted between queries and keys \cite{links:su2021roformer}. The code's \cfg{rope\_on\_value=True} also dresses the value link with the rotation, $V_iO_hT_hV_j^{-1}=U_iR_i^\top O_hT_hR_jU_j^{-1}$. The decomposition also states an identifiability limit. If a loss observes $G$ only through the compressed block $\Pi^R G\Pi^S$, changes in the unobserved blocks cannot be identified by that loss. The ambient-alignment and channel-likelihood receiver energies distinguish how a receiver treats directions outside the observed range: the former adds full ambient precision, while the latter leaves unobserved directions at their prior. Neither choice identifies independent hidden blocks $B,C,D$ unless the probabilistic model actually places those blocks in an observed term. A proposed edge-dependent $G_{ij}$ therefore needs both a generative specification and an explicit account of which of its degrees of freedom the loss can see.

\subsection{Supporting Option C: the implemented head sum}
\label{links:origin:sec:headsum}

\begin{proposition}[Head-sum link]
\label{links:origin:prop:headsum}
The edge operator that the projected routes apply at fixed routing,
\begin{equation}
\mathcal C_{ij}=\sum_h\beta_{ij}^{(h)}L_{ij}^{(h)}=U_iOB_{ij}TU_j^{-1},
\label{links:origin:eq:headsum}
\end{equation}
transforms as a link and lies in $\GL(K)$ whenever the stacked $O$ and $T$ are invertible and every $\beta_{ij}^{(h)}>0$. Its determinant is $\det(U_iU_j^{-1})\det(TO)\prod_h(\beta_{ij}^{(h)})^d$.
\end{proposition}

\begin{proof}
The identity $\sum_h\beta^{(h)}O_hT_h=OB T$ is block multiplication with $B=\bigoplus_h\beta^{(h)}I_d$, and the determinant is multiplicative.
\end{proof}

Invertibility of $O$ and $T$ is generic and holds at the code's identity initialization, where $O_h$ and $T_h^\top$ are the coordinate block embeddings. Softmax weights are strictly positive on finite-logit admitted support in exact arithmetic; masked edges and floating-point underflow need not satisfy this condition. Proposition~\ref{links:origin:prop:headsum} is the remark in the attention manuscript that the multi-head operator can be full rank \cite{unified:attention}, made into a statement about group values. The channel factor $B_{ij}$ is head diagonal. The ambient operator $\mathcal C_{ij}$ generally is not block diagonal, because the learned sending and receiving maps mix the full reference state; $TO$ records the cross-head receiving-to-sending transfer.

This is also exactly the value path of standard multi-head attention. With trivial frames, the multi-head output of a transformer is $\sum_hW_O^{(h)}\sum_j\beta_{ij}^{(h)}W_V^{(h)}x_j=\sum_jW_OB_{ij}W_Vx_j$ \cite{links:vaswani2017}, the frame-free case of Eq.~\eqref{links:origin:eq:headsum} with $O=W_O$ and $T=W_V$.

\paragraph{Loops measure a cocycle defect of the logits.}
With square invertible stacked maps and $TO=I_K$ the head blocks commute, and the triangle loop with one reversed edge is
\begin{equation}
\mathcal C_{ij}\mathcal C_{jk}\mathcal C_{ik}^{-1}
=U_iO\Big(\bigoplus_h\frac{\beta_{ij}^{(h)}\beta_{jk}^{(h)}}{\beta_{ik}^{(h)}}I_d\Big)O^{-1}U_i^{-1}.
\label{links:origin:eq:headsum_loop}
\end{equation}
For softmax routing $\beta_{ij}=e^{s_{ij}}/Z_i$, the logarithm of each head factor is
\begin{equation}
\log\frac{\beta_{ij}\beta_{jk}}{\beta_{ik}}=s_{ij}+s_{jk}-s_{ik}-\log Z_j,
\label{links:origin:eq:cocycle}
\end{equation}
because the receiver normalizer $Z_i$ cancels between $\beta_{ij}$ and $\beta_{ik}$. The loop is trivial exactly when the logits form a cocycle up to the intermediate normalizer. In the projected routes the routing is $\beta_{ij}^{(h)}=\pi_{ij}^{(h)}\exp(-D_{ij}^{(h)}/\tau_h)/Z_i^{(h)}$, with $D_{ij}^{(h)}$ the Gaussian KL between head $h$'s query law of $i$ and key law of $j$, so $s_{ij}=\log\pi_{ij}-D_{ij}/\kappa$ and
\begin{equation}
\log\frac{\beta_{ij}\beta_{jk}}{\beta_{ik}}=-\frac{D_{ij}+D_{jk}-D_{ik}}{\tau_h}+\log\frac{\pi_{ij}\pi_{jk}}{\pi_{ik}}-\log Z_j.
\label{links:origin:eq:kl_triangle}
\end{equation}
The leading term is the triangle defect of the routing costs. A metric would make $D_{ij}+D_{jk}-D_{ik}$ nonnegative, but the Kullback-Leibler divergence obeys no triangle inequality, and here it compares a query law with a key law under different maps $S_Q$ and $S_K$, so the defect can take either sign. With uniform priors, Eq.~\eqref{links:origin:eq:kl_triangle} says that the loop of the head-sum link measures how far the routing geometry is from additive along the path $i\to j\to k$, offset by the intermediate token's log normalizer. For dot-product logits $s_{ij}=q_i^\top k_j/\sqrt d$ the defect is instead a bilinear expression in the queries and keys of the three tokens. Reading attention weights as an abelian connection is an interpretation suggested by Eq.~\eqref{links:origin:eq:headsum_loop}; Proposition~\ref{links:origin:prop:headsum} establishes only the algebra.

\paragraph{Limits.}
Group-valuedness in Option C is generic rather than structural. The reference operator has the bound $\kappa(OB_{ij}T)\le\kappa(O)\kappa(B_{ij})\kappa(T)$; the local-coordinate bound also includes $\kappa(U_i)\kappa(U_j)$. The ratio of the largest to the smallest head weight enters this bound, and vanishing weights can produce singularity. The bound is not a lower bound or a monotonicity theorem. A singular square $O$ or $T$ makes the composite singular. If $Hd<K$ its rank is deficient; if $Hd>K$, rectangular factors can still produce a full-rank $K\times K$ composite, although that overcomplete case is outside the implemented $Hd=K$ architecture. In the one trained model measured in Section~\ref{links:origin:sec:measurement} the stacked maps are already ill conditioned, and the equal-head-weight composite is invertible but ill conditioned in the saved numerical measurements.

\subsection{Triangle holonomy on an inverse-completed causal graph}
\label{links:curve:sec:setting}

\begin{definition}[Triangle holonomy]
For $k<j<i$ let
\begin{equation}
F_{ijk}=\Theta_{ik}^{-1}\Theta_{ij}\Theta_{jk}=U_kG_{ik}^{-1}G_{ij}G_{jk}U_k^{-1}.
\label{links:curve:eq:holonomy}
\end{equation}
On the complete causal support, the field is \emph{flat} when $F_{ijk}=I$ for all strict triples. Missing-edge graphs require an explicit set of cycles and do not automatically have the recurrence characterization below.
\end{definition}

Under a frame change $F_{ijk}\mapsto g_kF_{ijk}g_k^{-1}$, so the conjugacy class of $F_{ijk}$, and hence its trace, determinant and spectrum, is gauge invariant. $F_{ijk}$ is the holonomy of the closed loop $k\to j\to i\to k$ based at $k$, in which the last step traverses the long edge backward and uses its inverse; this is the sense in which invertibility of the links is what gives a causal field nontrivial loop observables.

\subsection{Flatness is recurrence}
\label{links:curve:sec:flat}

\begin{proposition}[Flat is recurrent]
\label{links:curve:prop:recurrent}
Let $\Theta_{ij}$, $j<i$, be arbitrary $K\times K$ matrices, and call the field flat if $\Theta_{ij}\Theta_{jk}=\Theta_{ik}$ for all $k<j<i$. Then
\begin{equation}
\Theta_{ij}=\Theta_{i,i-1}\Theta_{i-1,i-2}\cdots\Theta_{j+1,j},
\label{links:curve:eq:recurrent}
\end{equation}
and conversely every choice of adjacent matrices $A_i=\Theta_{i,i-1}$ defines a flat field through the products~\eqref{links:curve:eq:recurrent}.
\end{proposition}

\begin{proof}
For $i-j\ge2$, flatness with the middle index $i-1$ gives $\Theta_{ij}=\Theta_{i,i-1}\Theta_{i-1,j}$, and induction on $i-j$ gives~\eqref{links:curve:eq:recurrent}. Conversely, products of adjacent matrices are associative, so $\Theta_{ij}\Theta_{jk}=\Theta_{ik}$.
\end{proof}

\begin{proposition}[Invertible flat is pure gauge]
\label{links:curve:prop:puregauge}
If every $\Theta_{ij}\in\GL(K)$, the field is flat if and only if there are $g_i\in\GL(K)$ with $\Theta_{ij}=g_ig_j^{-1}$, and then $g_i=\Theta_{i0}$ up to a common right factor. For singular matrices the equivalence fails: the shared projector $\Theta_{ij}=P$ with $P^2=P$ is flat and cannot be written $g_ig_j^{-1}$.
\end{proposition}

\begin{proof}
Set $g_0=I$ and $g_i=\Theta_{i0}=A_i\cdots A_1$. Then $g_ig_j^{-1}=A_i\cdots A_{j+1}=\Theta_{ij}$ by Proposition~\ref{links:curve:prop:recurrent}, and the converse is the cocycle $g_ig_j^{-1}g_jg_k^{-1}=g_ig_k^{-1}$. The projector satisfies $P\cdot P=P$ but is singular.
\end{proof}

Three consequences follow. First, the flat transport $\Omega_{ij}=U_iU_j^{-1}$ built from arbitrary frames, including frames of the form $U_i=\exp(X_{x_i})\exp(Y_i)$ with positional factors, is pure gauge, and so is any reference-coordinate field of the form $G_{ij}=h_ih_j^{-1}$, for example the rotary link $e^{(p_i-p_j)X}$, because it is absorbed into the frames. Second, the frames drop out of the conjugacy class of the holonomy: by Eq.~\eqref{links:curve:eq:holonomy}, $F_{ijk}$ is conjugate to $G_{ik}^{-1}G_{ij}G_{jk}$, so curvature is a property of the reference-coordinate field $G_{ij}$ and no choice of frames, positional or not, can flatten it or create it. A constant $G$ is flat only if $G=I$, whatever the positional dependence of the frames. Third, triangle holonomies detect failure of the cocycle condition on this complete ordered graph. They need not be a complete set of invariants when considered only through separate spectra or traces. With the reference frames also part of the data, the reference links $U_i^{-1}\Theta_{ij}U_j=G_{ij}$ are themselves invariant under passive site changes, up to simultaneous conjugation under a common reference change. Thus graph flatness does not erase all reference-relative content. Inverse-edge completion supplies a triangle observable without contradicting the absence of directed cycles on causal attention support.

\begin{remark}[Recurrent models]
Proposition~\ref{links:curve:prop:recurrent} says that a flat field is generated by its nearest-neighbor links through path composition. The transition products of a linear state-space recurrence $h_i=A_ih_{i-1}+\dots$, including recurrences whose transitions $A_i$ depend on the input, have exactly this form \cite{links:gu2022s4,links:gu2023mamba}, so these transition products form flat link fields. This does not identify every aspect of a recurrent architecture with a link field; input injections, readouts, and nonlinear updates remain additional structure. The triangle defect is the deviation of a long link from composition through an intermediate token. This is a statement about the cocycle structure of the transition products and not about the models' performance.
\end{remark}

\subsection{Curvature of the candidate generators}
\label{links:curve:sec:candidates}

\begin{lemma}[Second-order holonomy]
\label{links:curve:lem:bch}
For $\Theta_{ij}=U_ie^{X_{ij}}U_j^{-1}$, write $A=X_{ij}$, $B=X_{jk}$, $Z=X_{ik}$. Then $F_{ijk}=U_ke^{-Z}e^{A}e^{B}U_k^{-1}$ and
\begin{equation}
\log\big(e^{-Z}e^{A}e^{B}\big)=(A+B-Z)+\tfrac12\big([A,B]-[Z,A]-[Z,B]\big)+O(3).
\label{links:curve:eq:bch}
\end{equation}
The expansion holds near zero, on a domain where the indicated logarithm is analytic. The first-order term $A+B-Z$ is the additive cocycle defect $(dX)_{ijk}$ of the algebra-valued $1$-cochain $X$. On the complete ordered graph its vanishing is equivalent to $X_{ij}=\psi_i-\psi_j$: take $\psi_0=0$, $\psi_i=X_{i0}$, and apply the identity with $k=0$. An additive coboundary does not in general make the exponential links flat when its values fail to commute.
\end{lemma}

\begin{proof}
Apply the Baker-Campbell-Hausdorff series \cite{frame:Hall2015} twice: $\log(e^{A}e^{B})=A+B+\tfrac12[A,B]+\dots$, and then $\log(e^{-Z}e^{C})=C-Z-\tfrac12[Z,C]+\dots$ with $C=A+B+\tfrac12[A,B]$.
\end{proof}

\begin{table}[ht]
\centering
\small
\begin{tabular}{>{\raggedright}p{0.30\linewidth}>{\raggedright}p{0.36\linewidth}>{\raggedright\arraybackslash}p{0.26\linewidth}}
\toprule
Link field & Holonomy $F_{ijk}$ & Flat when \\
\midrule
Positional, $U_ie^{(p_i-p_j)X}U_j^{-1}$ & $I$ & always \\
Constant, $U_ie^{X}U_j^{-1}$ & $U_ke^{X}U_k^{-1}$ & only if $e^{X}=I$ \\
Routed exponent, $X_{ij}=OB_{ij}T$, $TO=I$ & $U_kOe^{B_{ij}+B_{jk}-B_{ik}}TU_k^{-1}$ & $B$ additive: $B_{ij}+B_{jk}=B_{ik}$ \\
Head sum, $U_iOB_{ij}TU_j^{-1}$, $TO=I$ & $U_kO\frac{B_{ij}B_{jk}}{B_{ik}}TU_k^{-1}$ & $B$ multiplicative: $B_{ij}B_{jk}=B_{ik}$ \\
\bottomrule
\end{tabular}
\caption{Triangle holonomy of the candidate edge-dependent link fields. The routing matrices $B_{ij}=\bigoplus_h\beta_{ij}^{(h)}I_d$ act headwise. The last two rows assume square invertible stacked maps and $TO=I_K$; positive head weights are needed for the head-sum inverse.}
\label{links:curve:tab:candidates}
\end{table}

Table~\ref{links:curve:tab:candidates} collects the exact holonomies. The first row is Proposition~\ref{links:curve:prop:puregauge}. The second shows that the constant link of the companion note is nonflat whenever $e^X\ne I$: each triangle has the reference holonomy $e^X$, whose first-order term is $X=(dX)_{ijk}$. No maximality statement follows because no ordered measure of curvature has been specified. For the routed exponent and the head sum with $TO=I$, the head blocks commute and the formulas follow by multiplying block by block; in the routed exponent the generators add, in the head sum the weights multiply.

The multiplicative log defect is Eq.~\eqref{links:origin:eq:kl_triangle}. It includes the source prior and the intermediate normalizer, as well as the triangle defect of the Gaussian comparison costs. Its value therefore cannot be attributed to the KL term alone.

\begin{proposition}[Flat normalized routing]
\label{links:curve:prop:flatrouting}
Let $\beta_{ij}>0$ for $j\le i$ with $\sum_{j\le i}\beta_{ij}=1$. The head-sum link with $TO=I$ is flat on strict triangles $k<j<i$ if and only if $\beta_{ij}=c_i/c_j$ for $0<j<i$ and $\beta_{i0}=c_i$, where $c_i:=\beta_{i0}$ and $c_0=\beta_{00}=1$.
\end{proposition}

\begin{proof}
By Table~\ref{links:curve:tab:candidates}, flatness is $\beta_{ij}\beta_{jk}=\beta_{ik}$ for each head and strict triple. Taking $k=0$ gives $\beta_{ij}\beta_{j0}=\beta_{i0}$, so $\beta_{ij}=c_i/c_j$ with $c_i=\beta_{i0}$, and $c_0=\beta_{00}=1$ because the first row has the single source $0$. Conversely this form satisfies $\beta_{ij}\beta_{jk}=c_ic_j/(c_jc_k)=\beta_{ik}$.
\end{proof}

The form $\beta_{ij}=c_iw_j$ with $w_j=1/c_j$ is a separable attention pattern in which the weight on source $j$ is the reciprocal of $j$'s own first-source weight. The diagonal entries are fixed by normalization as $\beta_{ii}=1-c_i\sum_{j<i}c_j^{-1}$ and must remain positive. For example, $c=(1,0.5,0.2,0.05,0.01)$ satisfies these inequalities, but it is a thin subset: uniform causal attention has $\rho_{ijk}=1/(j+1)$ on every strict triangle and therefore fails as soon as three positions are available (the five-token validation example has largest additive weight defect $1/6$), and generic content-dependent routing has a nonzero multiplicative defect. This characterizes the head-sum field under $TO=I_K$, not general attention or the learned maps after unconstrained training.

\subsection{Conditional regression as a proposed edge-dependent generator source}
\label{links:curve:sec:regression}

The exponential construction ensures invertibility but does not select its generator. Joint second-order statistics give one candidate source of edge dependence. Let $(X_0,\ldots,X_{N-1})$ be a jointly Gaussian reference-coordinate random vector, with means $\xi_i$ and cross-covariances $C_{ab}=\operatorname{Cov}(X_a,X_b)$, and assume $C_{bb}\succ0$. Its centered best linear predictor is
\begin{equation}
W_{ab}=C_{ab}C_{bb}^{-1}
=\arg\min_W\mathbb E\left\|(X_a-\xi_a)-W(X_b-\xi_b)\right\|^2_{C_{aa}^{-1}}.
\label{links:curve:eq:regression}
\end{equation}
The weighting by $C_{aa}^{-1}$ does not change the unrestricted minimizer: differentiating the quadratic risk gives $W C_{bb}=C_{ab}$. The full conditional mean is $\xi_a+W_{ab}(X_b-\xi_b)$ and includes an intercept unless the law is centered. In the present passive convention the reference moments, $C_{ab}$, and $W_{ab}$ are invariant, whereas the local-coordinate link is
\begin{equation}
\Theta^{\mathrm{reg}}_{ab}=U_aW_{ab}U_b^{-1},
\qquad \Theta^{\mathrm{reg}\prime}_{ab}=g_a\Theta^{\mathrm{reg}}_{ab}g_b^{-1}.
\label{links:curve:eq:regression_dressed}
\end{equation}
Equivalently, if the joint law is expressed directly in local coordinates, its cross-covariance transforms as $\Sigma_{ab}\mapsto g_a\Sigma_{ab}g_b^\top$, and $\Sigma_{ab}\Sigma_{bb}^{-1}$ obeys the same endpoint law. A regression matrix can be singular; Gaussianity and positive-definite marginals do not make it a group element.

\begin{proposition}[Regression composition and conditional covariance]
\label{links:curve:prop:markov}
For $k<j<i$, define $C_{ik\mid j}=C_{ik}-C_{ij}C_{jj}^{-1}C_{jk}$. Then
\begin{equation}
W_{ij}W_{jk}=W_{ik}-C_{ik\mid j}C_{kk}^{-1}.
\label{links:curve:eq:partialcov}
\end{equation}
Consequently the regression cocycle is exact on that triple if and only if $C_{ik\mid j}=0$, which for a jointly Gaussian law is $X_i\perp X_k\mid X_j$. If the three links are invertible, the corresponding group holonomy is
\begin{equation}
F_{ijk}=U_k\left(I-W_{ik}^{-1}C_{ik\mid j}C_{kk}^{-1}\right)U_k^{-1}.
\label{links:curve:eq:regression_holonomy}
\end{equation}
\end{proposition}
\begin{proof}
Multiplication gives
$W_{ij}W_{jk}=C_{ij}C_{jj}^{-1}C_{jk}C_{kk}^{-1}=(C_{ik}-C_{ik\mid j})C_{kk}^{-1}$.
Since $C_{kk}$ is invertible, equality to $W_{ik}$ is equivalent to $C_{ik\mid j}=0$. The conditional law of $(X_i,X_k)$ given $X_j$ is Gaussian with that off-diagonal covariance block, and its density or characteristic function factors precisely when the block vanishes \cite{links:lauritzen1996}. In the invertible case, substituting the dressed regression links into the triangle product yields the displayed holonomy.
\end{proof}

The composition defect records conditional linear dependence beyond the middle variable. Its entries or norm are not a mutual information, and its magnitude depends on the chosen tensor norm. For a homogeneous Gaussian vector Markov chain $X_i=A X_{i-1}+\varepsilon_i$ with independent innovations, $C_{ij}=A^{i-j}C_{jj}$ and $W_{ij}=A^{i-j}$. The cocycle is flat even if $A$ is singular; group-valued holonomy additionally requires invertible transitions. Time-dependent transitions give the ordered product $W_{ij}=A_i\cdots A_{j+1}$.

For the complete ordered sequence, the Gaussian cocycle condition on every strict triple also implies the first-order Markov property. Indeed, set $j=i-1$ in Eq.~\eqref{links:curve:eq:partialcov}. Vanishing for each $k<i-1$ makes the entire conditional cross-covariance between $X_i$ and the vector $(X_0,\ldots,X_{i-2})$ zero given $X_{i-1}$; conditional Gaussianity then gives independence of those vectors. The converse follows from the Markov conditional independence. Thus the all-triples statement and the Gaussian Markov property are equivalent at the covariance-cocycle level, rather than merely a one-way example. For non-Gaussian laws, zero partial covariance alone does not imply conditional independence.

\paragraph{The channel as reduced-rank regression.}
Whiten both variables and let
$C_{ii}^{-1/2}C_{ij}C_{jj}^{-1/2}=\sum_s a_su_sv_s^\top$
be a singular value decomposition, with $a_1\ge a_2\ge\cdots$. These singular values are the canonical correlations. The square roots here are the symmetric positive-definite square roots.
\begin{proposition}[Reduced-rank regression]
\label{links:curve:prop:rrr}
Among matrices of rank at most $d$, a minimizer of the centered risk in Eq.~\eqref{links:curve:eq:regression} is
\begin{equation}
W_d=C_{ii}^{1/2}\left(\sum_{s\le d}a_su_sv_s^\top\right)C_{jj}^{-1/2},
\qquad \mathcal R(W_d)-\mathcal R(W_{ij})=\sum_{s>d}a_s^2.
\label{links:curve:eq:rrr}
\end{equation}
The minimizer need not be unique when the truncation falls within a tied singular-value subspace.
\end{proposition}
\begin{proof}
Completing the quadratic risk gives
$\mathcal R(W)-\mathcal R(W_{ij})=\|C_{ii}^{-1/2}(W-W_{ij})C_{jj}^{1/2}\|_F^2$.
The invertible whitening maps preserve rank, so the Eckart--Young singular-value truncation solves the constrained problem and gives the stated discarded-square error \cite{links:izenman1975,links:reinsel1998}.
\end{proof}
A factorization $W_d=OT$ is therefore an optimal width-$d$ channel for this declared risk. Its range is the image under $C_{ii}^{1/2}$ of the retained left canonical directions, and its sending directions are the corresponding right directions after source whitening. Partitioning these directions among several heads is an additional architecture choice; the present learned heads are not trained by this regression objective and need not coincide with its solution.

These regression links are candidates for a correlated extension. Only local token marginals are stored by the implementation. Cross-covariances between distinct random token states vanish under the product-state assumption used for its predictive moment kernel, so regression defined from that product law would be zero. Correlation of deterministic token features across a dataset is a different object and does not by itself supply a correlated variational law. A coupled generative prior or a declared correlated recognition family would provide the required $C_{ij}$. Where a desired regression link has a real logarithm, it can be represented by the primary exponential construction. In general it may be singular or lack a real logarithm, so a choice to use its low-rank factors as a generator, $X_{ij}=O_{ij}T_{ij}$, is a further modeling step; it does not equate regression with its exponential.

\subsection{Fisher geometry of exponential links and the singular boundary}
\label{links:curve:sec:natural}

A group element $g\in\GL(K,\mathbb R)$ sends the fixed nondegenerate Gaussian $q=\mathcal N(\mu,\Sigma)$ to $q_g=\mathcal N(\mu',\Sigma')$, with $\mu'=g\mu$ and $\Sigma'=g\Sigma g^\top$. Write a tangent as $\delta g=Ag$, and let $L'$ be a Cholesky factor of $\Sigma'$.
\begin{proposition}[Gaussian Fisher pullback]
\label{links:curve:prop:fisher}
The Gaussian Fisher--Rao metric pulls back to the quadratic form
\begin{equation}
ds^2=\|Z\eta\|^2+2\left\|\frac{Z+Z^\top}{2}\right\|_F^2,
\qquad Z=L'^{-1}AL',\quad \eta=L'^{-1}\mu'.
\label{links:curve:eq:fisher}
\end{equation}
For $g=e^X$, the right-trivialized differential is
\begin{equation}
A=\int_0^1 e^{sX}\delta X e^{-sX}ds
=\frac{e^{\operatorname{ad}_X}-1}{\operatorname{ad}_X}\delta X,
\label{links:curve:eq:dexp}
\end{equation}
where the quotient denotes its entire power series, including at zero eigenvalues of $\operatorname{ad}_X$.
\end{proposition}
\begin{proof}
The output tangent is $\delta\mu'=A\mu'$ and $\delta\Sigma'=A\Sigma'+\Sigma'A^\top$. Substitution in the Gaussian Fisher form
$\delta\mu'^\top\Sigma'^{-1}\delta\mu'+\tfrac12\operatorname{tr}(\Sigma'^{-1}\delta\Sigma'\Sigma'^{-1}\delta\Sigma')$
and whitening by $L'$ give Eq.~\eqref{links:curve:eq:fisher}. The integral follows by differentiating the exponential series, or the equivalent integral formula for its Fr\'echet derivative, and multiplying by $e^{-X}$ on the right \cite{frame:Hall2015}.
\end{proof}

This is the same belief-dependent pullback used in the frame-chart construction. It is positive semidefinite on $\GL(K)$, not generally a positive-definite metric there: a whitened skew-symmetric $Z$ with $Z\eta=0$ fixes the Gaussian to first order and has zero length. It becomes a metric after quotienting by the Gaussian stabilizer, or after a declared damping or additional term removes the null directions. A routing-weighted sum of source-belief pullbacks is one possible metric for a content field, not the Fisher information of the complete source mixture automatically: observing the source label and marginalizing it lead to different statistical models.

\begin{proposition}[A contraction path reaches the singular boundary only at infinite distance]
\label{links:curve:prop:barrier}
Let $g_t=\operatorname{diag}(e^{-t},1,\ldots,1)$ act on $q_0=\mathcal N(\mu,I)$, with $t\ge0$. Then
\begin{equation}
\mathrm{KL}(q_t\|q_0)
=t+\tfrac12\left(e^{-2t}-1+(1-e^{-t})^2\mu_1^2\right)\sim t,
\label{links:curve:eq:klbarrier}
\end{equation}
the length of this path is $\sqrt{\mu_1^2+2}t$, and its endpoint Fisher--Rao distance obeys $d_{\mathrm{FR}}(q_0,q_t)\ge\sqrt2t$.
\end{proposition}
\begin{proof}
The Gaussian KL formula gives Eq.~\eqref{links:curve:eq:klbarrier}. Along the path, $A=-E_{11}$, $L'=g_t$, $Z=-E_{11}$, and $\eta=\mu$, so Eq.~\eqref{links:curve:eq:fisher} yields $(\mu_1^2+2)dt^2$. The univariate Gaussian Fisher metric is $(dm^2+2ds^2)/s^2$, so every path between standard deviations $1$ and $e^{-t}$ has length at least $\sqrt2t$ by integration of $|d\log s|$. Marginalization onto the first coordinate contracts Fisher length and hence cannot increase endpoint distance \cite{frame:AmariNagaoka2000}.
\end{proof}

The infinite-distance conclusion extends to any finite singular limit of $g$ acting on a fixed nondegenerate Gaussian. To see this, fix a reference covariance $\Sigma_0\succ0$ and an endpoint covariance $\Sigma_g=g\Sigma g^\top$. Choose a scalar projection whose variance ratio is the smallest eigenvalue $\lambda_{\min}$ of $\Sigma_0^{-1/2}\Sigma_g\Sigma_0^{-1/2}$. Marginalization and the same univariate bound give
\begin{equation}
d_{\mathrm{FR}}(q_0,q_g)\ge\frac{1}{\sqrt2}|\log\lambda_{\min}|.
\label{links:curve:eq:general_boundary}
\end{equation}
As $g$ approaches a finite singular matrix, $\lambda_{\min}\to0$. The covariance part of the Gaussian distance therefore diverges even though the pullback has stabilizer null directions. The specific coefficient in $\mathrm{KL}\sim-\log\sigma_{\min}(g)$ was established for the one-direction contraction above. Along a general path, KL also depends on all other singular values and the transformed mean; one universal asymptotic coefficient does not follow from the smallest singular value alone.

\begin{corollary}[Additive and multiplicative updates]
\label{links:curve:cor:updates}
An additive matrix update can cross the singular set in one discrete step. Its spectral-norm distance from an invertible matrix $M$ to that set is $\sigma_{\min}(M)$: if $u,v$ are corresponding unit singular vectors, $M-\sigma_{\min}(M)uv^\top$ is singular. In contrast,
\begin{equation}
g^+=e^{-\eta A}g,
\qquad \det g^+=e^{-\eta\operatorname{tr}A}\det g
\label{links:curve:eq:group_update}
\end{equation}
remains in the same determinant component for every finite real $A$ and $\eta$ in exact arithmetic.
\end{corollary}
\begin{proof}
The stated rank-one perturbation annihilates $v$, and perturbations of norm below $\sigma_{\min}(M)$ cannot make $M$ singular. The exponential determinant identity proves the multiplicative statement.
\end{proof}
A natural-gradient direction does not by itself impose a fixed finite-step Fisher length or a universal iteration count. If an algorithm additionally limits each realized path segment to Fisher length at most $\ell$, reaching distance $D$ requires at least $D/\ell$ segments. That conditional bound applies to Eq.~\eqref{links:curve:eq:general_boundary}; it cannot be inferred solely from the learning-rate parameter. Matrix-exponential overflow, underflow, and numerical singularity also remain numerical concerns despite exact group membership.

The boundary argument concerns an exact Gaussian pushforward with no added noise. An output law with a positive covariance floor or additive observation noise has a different statistical geometry and need not encounter that same degeneracy when its linear mean map becomes singular. It is therefore not a claim that the implemented training objective diverges at a content-map singularity. The code exposes multiplicative frame-table policies through \cfg{m\_phi\_update\_mode="pullback\_group"} and \cfg{"fisher\_group"}; these apply to stored frame tables, with the metric and retraction qualifications in the frame section. Applying such a construction to content links would require a new parameterization and optimizer path. The currently inspected launcher selects \cfg{m\_phi\_update\_mode="adamw"}, and its content maps remain ordinary learned factors.

\subsection{A measurement in a trained model}
\label{links:origin:sec:measurement}

Proposition~\ref{links:origin:prop:charts} predicts which quantities decide whether the affine and logarithmic charts agree for a trained operator: the sign of $\det M$, the smallest singular value, and the position of the eigenvalues relative to the negative real axis. The source note records their evaluation on one finished run of the projected predictive route, \texttt{54.73\_wikitext-103\_K64\_depth-scaled} ($K=64$, $H=2$, $d=32$, four layers, \cfg{relational\_output\_init="depth\_scaled"}, attention residual on, seed 6, 15000 steps, test perplexity $54.73$). The inputs are the final weights and the operator snapshots stored at twelve steps; the stacked operator is $M=\sum_hO_hT_h=OT$ in reference coordinates, which is the message operator when every head routes to one source, and $I+M$ is the residual-augmented self-operator when all heads select the receiver with weight one. The latter condition holds at the first position under a causal prior admitting self-attention. For a general token the self-operator is $I+OB_{ii}T$; the other incoming edges remain $OB_{ij}T$. The whole message is a sum over source states and is not one $K\times K$ group link. The archived scripts and outputs are in \texttt{link\_origin\_checks/run\_54.73/}. The saved outputs were read during consolidation; the original checkpoint measurements were not rerun, and no new model was trained.

\paragraph{The bare operator is close to the singular set.}
The depth-scaled initialization gives $M=2^{-3/2}I$, with all singular values equal to $0.354$ and stable rank $64$. By step $1000$ the stable rank has fallen to $5$--$13$, and from then on $\sigma_{\min}(M)$ lies between $10^{-5}$ and $4\times10^{-3}$, and at the final step $\sigma_{\max}$ is $1.2$ to $1.9$; the minimum over all later snapshots is $9.7\times10^{-6}$. At the final step, $90\%$ of $\|M\|_F^2$ lies in $21$--$30$ of the $64$ singular directions, and the smallest singular values of the stacked $O$ and $T$ lie between $10^{-3}$ and $8\times10^{-2}$, so some directions are barely written and some barely read. The sign of $\det M$ is not stable: counted over consecutive snapshot pairs from step $1000$ on, it changes $5$, $5$, $3$ and $4$ times in layers $0$ to $3$, $17$ times in $40$ pairs. These snapshots show ill-conditioned nonsingular matrices with changing computed determinant signs; they do not show that a sampled matrix lies exactly at $\det=0$, or identify the intervening optimizer path. A factor learning rate is not itself a measured update norm of the composite $M$. At the final step the determinant is negative in layers $1$--$3$, and since $\det(OBT)=\det O\det B\det T$ with $\det B>0$, that sign holds for every positive routing weight. Between $4$ and $20$ eigenvalues of $M$ have negative real part, mostly near zero, with the most negative at $-0.15$ to $-0.38$ at the end and $-0.57$ in layer $0$ at step $1000$.

\paragraph{The sampled residual-augmented operators are well conditioned.}
$I+M$ has positive determinant at every snapshot and layer, because every eigenvalue has positive real part, the smallest being $0.43$ (layer $0$, step $1000$) and $0.54$ or larger from step $1500$ on. The principal logarithm therefore exists, and the hypothesis of item~4 of Proposition~\ref{links:origin:prop:charts} holds throughout. The smallest singular value of $I+M$ over all later snapshots is $0.27$, and at the final step the condition number is $4$--$6$ in every layer (Table~\ref{links:origin:tab:measured}). Over the per-head weight box $(b_1,b_2)\in[0,1]^2$, evaluated on a $41\times41$ grid at the final weights, $I+OBT$ never has $\det\le0$ and has $\sigma_{\min}\ge0.38$ and condition number at most $6.1$.

\begin{table}[ht]
\centering
\small
\begin{tabular}{rrrrrr}
\toprule
Layer & $\operatorname{sign}\det M$ & $\sigma_{\min}(M)$ & $\kappa(M)$ & $\sigma_{\min}(I+M)$ & $\min\operatorname{Re}\lambda(I+M)$ \\
\midrule
0 & $+$ & $8.1\times10^{-4}$ & $2.0\times10^{3}$ & $0.38$ & $0.62$ \\
1 & $-$ & $3.0\times10^{-5}$ & $4.8\times10^{4}$ & $0.56$ & $0.70$ \\
2 & $-$ & $3.8\times10^{-4}$ & $3.2\times10^{3}$ & $0.52$ & $0.74$ \\
3 & $-$ & $1.7\times10^{-3}$ & $1.1\times10^{3}$ & $0.69$ & $0.85$ \\
\bottomrule
\end{tabular}
\caption{Final-step chart quantities for $M=\sum_hO_hT_h$ and $I+M$ in the trained run.}
\label{links:origin:tab:measured}
\end{table}

\paragraph{Interpretation.}
For this run, the sampled $I+M$ lies in the principal real-logarithm domain, while $M$ is substantially closer to the singular set and three final layers have negative determinant. Its invertibility is therefore algebraically sufficient for inverse-edge completion but can be numerically fragile. The final stacked factors have $\kappa(O)$ between $126$ and $2265$ and $\kappa(T)$ between $19$ and $172$. These findings concern the saved reference matrices and the sampled weight grid; they do not establish invertibility of every edge, every possible routing choice, or the whole nonlinear block. Small singular values indicate approximately suppressed directions, not an exact common null complement $\ker T$. The measurement motivates examining explicitly restricted channels and group-constrained links, but gives no performance result for either construction. No residual ablation was run, so the observed conditioning cannot be assigned causally to the residual architecture.

\paragraph{Limits.}
The measurement is one run, one seed, with $K=64$ and two heads. Snapshots are $1000$ to $1500$ steps apart, so crossings between them are not seen. The routing-weight box was evaluated only at the final weights, and the realized per-edge routing was not used; the mean self weights at the end are $0.04$--$0.13$, and at position $0$ the only source is the token itself, so $B=I$ exactly and the unrouted operator is applied there. The source analysis reported no compatible comparison run. It did not measure individual content-optimizer increments or resolve the timing of determinant changes. Its observations do not establish a scaling law, seed stability, or a performance effect.

\subsection{Implemented triangle diagnostics and their scope}
\label{links:sec:diagnostics}

The October 3 implementation snapshot contains a routing-triangle diagnostic in addition to the algebra recorded in the original notes. It measures the current head-sum field; it does not implement a learned exponential edge field. The distinction is visible directly in the runtime path. The launcher selects \cfg{attention\_mode="projected\_predictive"}, \cfg{attention\_residual=True}, and \cfg{relational\_output\_init="identity"}. The configuration derives \cfg{relational\_readout\_mode="predictive\_content"}; model construction passes the residual only to that route, builds equal head widths with $K=Hd$, and initializes $T_h$ and $O_h$ on corresponding coordinate blocks. The \cfg{depth\_scaled} alternative multiplies $O_h$ by $1/\sqrt{2L}$ for $L$ layers. The predictive method contracts $T_h\xi_j$, pools by $\beta_{ij}^{(h)}$, writes by $O_h$, and returns through the value frame. The residual is added afterward. These maps are unrestricted learned rectangular factors; neither stacked invertibility nor $TO=I$ is enforced during training.

When rotary position encoding participates in values, the content frame is $V_i=U_iR_i^\top$ and the link is $V_iO_hT_hV_j^{-1}$. With \cfg{rope\_on\_value=False}, the content path retains $U_i$ instead. In either case the endpoint factors cancel in triangle conjugacy formulas. The currently inspected launcher sets that value flag false. This frame choice does not turn a nontrivial shared processing map into a flat transport.

\begin{proposition}[General tiled-head triangle spectrum]
\label{links:prop:general_head_holonomy}
Assume square invertible stacked $O,T$, positive routing on the three edges, and define
\begin{equation}
G=TO,\qquad R_{ijk}=\bigoplus_h\rho_{ijk}^{(h)}I_d,
\qquad \rho_{ijk}^{(h)}=\frac{\beta_{ij}^{(h)}\beta_{jk}^{(h)}}{\beta_{ik}^{(h)}}.
\label{links:eq:rho_transfer}
\end{equation}
The triangle holonomy of $\mathcal C_{ab}=U_aOB_{ab}TU_b^{-1}$ is similar to $R_{ijk}G$ even when $TO\ne I$.
\end{proposition}
\begin{proof}
Removing the external $U_k$ conjugation gives
\begin{equation}
F^{\mathrm{ref}}_{ijk}
=T^{-1}B_{ik}^{-1}B_{ij}G B_{jk}T.
\label{links:eq:general_head_product}
\end{equation}
Conjugation by $B_{jk}T$ gives $B_{jk}B_{ik}^{-1}B_{ij}G=R_{ijk}G$, because the head-diagonal routing matrices commute with one another. They need not commute with $G$.
\end{proof}

The routing factor and the receiving-to-sending transfer therefore remain distinct. Zero routing log defect gives $R=I$, but the full holonomy is then similar to $G$ and is not generally the identity. Conversely separate eigenvalue summaries of $R$ and $G$ cannot reconstruct the eigenvalues of $RG$ when they fail to commute. This generalization connects the existing diagnostic to the earlier formulas, which simplify to pure routing defects under $TO=I$.

The implementation computes, on admitted strict triangles, the exact path-weighted moments
\begin{equation}
S_0^{(h)}=\sum_{b,k<j<i}w_{bijk}^{(h)},\quad
S_1^{(h)}=\sum_{b,k<j<i}w_{bijk}^{(h)}x_{bijk}^{(h)},\quad
S_2^{(h)}=\sum_{b,k<j<i}w_{bijk}^{(h)}(x_{bijk}^{(h)})^2,
\label{links:eq:diagnostic_moments}
\end{equation}
where $w_{ijk}^{(h)}=\beta_{ij}^{(h)}\beta_{jk}^{(h)}$ and $x_{ijk}^{(h)}=\log\rho_{ijk}^{(h)}$, with batch index $b$. All three edges must be admitted and positive. Matrix products evaluate these sums without materializing an $N^3$ tensor. Training diagnostics report $S_1/S_0$, $\sqrt{S_2/S_0}$, and a row-normalized path mass for both the learned router and its normalized prior in \cfg{diagnostics/routing\_holonomy.csv}. A head with no admitted positive path mass has no defined normalized defect statistic.

The posthoc module also samples strict triangles using the two-step path weight formed from head-mean routing, reads the learned content factors from a checkpoint, and computes the eigenvalues of $RG$ in float64. Snapshot analysis computes the $G$ spectra alone. Its scalar spectral summary is
\begin{equation}
\chi(F)=\left[\frac1K\sum_{a=1}^K
\left\{\log^2|\lambda_a(F)|+\arg^2\lambda_a(F)\right\}\right]^{1/2},
\label{links:eq:spectral_summary}
\end{equation}
using principal complex arguments. The code calls this quantity \cfg{kappa}; the symbol $\chi$ here avoids confusion with routing temperature. It vanishes when all eigenvalues equal one, but cannot detect a nontrivial unipotent Jordan part. It is consequently a spectral summary, not a complete holonomy invariant or a metric on conjugacy classes. The scale-free variant subtracts the mean log modulus before squaring and therefore also removes a positive scalar dilation. A zero eigenvalue produces an infinite logarithm. In the absence of invertibility, the calculated spectrum of $RG$ remains a defined matrix diagnostic but cannot be identified with an inverse-edge holonomy. Very small eigenvalues and nonnormal matrices can make spectral estimates sensitive; the diagnostic does not certify conditioning of the inverse triangle product.

Five existing tiny CPU checks were executed during consolidation, covering exact moments against enumeration, flat and uniform routing, the $RG$ similarity spectrum, training CSV production, and posthoc triangle sampling. They passed using $K=4$, two heads in the model tests, vocabulary size seven, and at most six token positions; these dimensions validate identities and plumbing only. No trained-model triangle statistics or training comparison were generated by this consolidation. The diagnostic files were already present as working-tree changes at initial inspection and were subsequently committed as VFE 3.0 pull request 40 during this consolidation. Their bytes remained unchanged, and this task did not edit them.

\subsection{Consequences for future integration}
\label{links:sec:integration}

The exponential construction is the primary geometric model: it gives structural invertibility for every finite generator, a precise low-rank channel expansion, a complement split, and inverse-edge observables. Its connection to the present predictive route is exact on the specified chart domain after changing coordinates, and its tiled-head aggregation is expressible by transformed receiving maps. A routing-independent generator therefore supplies a geometric parameterization and optimization choice rather than an automatic increase in the represented mean-map family. Edge-dependent generators, a constrained group update, or a declared correlated model would add assumptions beyond that parameterization.

Compression explains how a selected channel can sit inside a larger group element and separates communicating, hidden, and mixing blocks. It is useful for describing what an exponential link exposes to a receiver and what a loss can identify. The current head sum supplies the implementation comparison: it is the ordinary multi-head value contraction in reference coordinates, dressed by the endpoint frames, and it becomes group-valued only under rank and positivity assumptions. The exact exponential map, primitive hidden-block field, correlated regression generator, complement-only residual, and routed exponential self-gate are not current learned content mechanisms merely because their formulas can be written using the present matrices.

The residual, covariance propagation, and receiver objective remain separate decisions. A complement-preserving exponential aggregation yields a projected residual only after its persistence assumption; the current architecture adds a full residual; Gaussian fusion yields precision-weighted persistence; and moment-matched prediction propagates the covariance of its declared stochastic kernel. No one of these choices follows from passive gauge covariance alone. Likewise, positive self-gate representations of SiLU or GELU place existing nonlinearities in exponential notation without deriving their specific generator. Template routing supplies a separately declared probabilistic model for a normalized site nonlinearity.

The outstanding implementation questions are whether group-constrained content optimization or nonconstant generators change training, whether a correlated model supplies useful regression links, and whether a deliberately uncommunicated complement should persist locally. The existing measurements establish none of their performance effects. The triangle diagnostic now provides a way to measure the current routing and transfer factors, but future comparisons must distinguish their spectra from the exactly flat frame coboundary, from direct causal paths, and from a contextual bundle connection.

\subsection{Archived algebra checks for the link-origin construction}
\label{links:origin:sec:checks}

The following values are retained from the original saved float64 output files, which were read during consolidation. These scripts were not rerun. The written derivations above carry the mathematical claims; the small matrices check implementations of the formulas.

The script \texttt{link\_origin\_checks/checks.py}, with its output in \texttt{checks\_output.txt}, evaluates each identity in float64 at $K=6$, $H=3$, $d=2$ with random frames and maps. Lemma~\ref{links:origin:lem:lowrank_exp} holds to a relative error of $5.6\times10^{-16}$, its routed form Eq.~\eqref{links:origin:eq:exp_routed} to $9.1\times10^{-15}$, the determinant identity to $1.5\times10^{-16}$, and the fixed action on $\ker T$ to $5.6\times10^{-16}$. The split of Eq.~\eqref{links:origin:eq:split} holds to $7.3\times10^{-16}$, and the first-order defect $(\Theta-\Omega)/s-L$ scales linearly with $s$ ($2.1\times10^{-2}$ at $s=10^{-2}$, $2.2\times10^{-3}$ at $s=10^{-3}$). The block factorization holds to $2.6\times10^{-11}$, limited by the conditioning of the random block $A$, and its determinant identity to $4.3\times10^{-15}$. The loop formulas Eq.~\eqref{links:origin:eq:headsum_loop} and Eq.~\eqref{links:origin:eq:exp_loop} hold to $2.7\times10^{-16}$ and $7.4\times10^{-16}$, and the loop trace is unchanged under a frame change. The script \texttt{checks\_aggregation.py}, with output in \texttt{checks\_aggregation\_output.txt}, checks Section~\ref{links:origin:sec:aggregation}. The complement split of Lemma~\ref{links:origin:lem:complement} holds to $1.6\times10^{-15}$ at $K=7$ with a rank-4 channel, and $I-P$ fixes $\ker T$ and annihilates $\im O$ to $1.2\times10^{-15}$. Exact enumeration over the head-label tuple reproduces the closed form of Eq.~\eqref{links:origin:eq:aggregate} with a receiver-carried complement to $4.2\times10^{-16}$. With tiled heads $P=I$ to $6.1\times10^{-16}$, and the Option A message equals the head sum with receiving maps $OE_{:,h}$ exactly. The fusion residual of Eq.~\eqref{links:origin:eq:fusion_residual} holds to $5.6\times10^{-16}$, and the SiLU equals the exponential of its log-sigmoid generator to rounding. The script \texttt{checks\_chart.py}, with output in \texttt{checks\_chart\_output.txt}, checks Section~\ref{links:origin:sec:charts} at $K=7$, $d=3$ and, for the tiled case, $K=6$. Item 4 of Proposition~\ref{links:origin:prop:charts} holds to $5.2\times10^{-16}$ and $2.5\times10^{-15}$ in the relative error of Eq.~\eqref{links:origin:eq:conversion} on the two random trials in the stated principal-logarithm domain, with $TO=Z$ to $4.4\times10^{-15}$. The embedding $\iota$ is a homomorphism to $2.1\times10^{-16}$, $\exp(OT)=\iota(e^{TO})$ and $I+OCT=\iota(I+CTO)$ hold to $5.3\times10^{-16}$ and $3.1\times10^{-16}$, and $\det\iota(B)=\det B$ to $2.2\times10^{-16}$. In the tiled case $O'=e^{X}T^{-1}$ gives $O'T=e^{X}$ to $1.8\times10^{-15}$, and the identity initialization has $\log I=0$. At $X=I$ the exponential and affine charts differ by $36\%$ in relative size. These checks ground the algebra at one small size; they say nothing about training behavior.

\subsection{Archived algebra checks for the curvature construction}
\label{links:curve:sec:checks}

These values likewise come from the saved original outputs, rather than a new experiment. Their small sizes were chosen to check matrix identities, without drawing conclusions about large-agent behavior or model performance.

The script \texttt{link\_curvature\_checks/checks.py}, with output in \texttt{checks\_output.txt}, evaluates each identity in float64 at $K=4$ with random invertible matrices. Recurrent links satisfy $F_{ijk}=I$ to $1.8\times10^{-15}$ and the pure-gauge form $g_ig_j^{-1}$ to $1.1\times10^{-15}$; the shared rank-one projector satisfies the cocycle exactly and has rank $1$. The holonomy of the constant link equals $U_ke^{X}U_k^{-1}$ to $2.0\times10^{-15}$, the positional link has $\max|F-I|=3.3\times10^{-13}$, the routed exponent matches its additive formula to $2.2\times10^{-15}$, and the head sum matches its multiplicative formula to $2.1\times10^{-15}$. The Baker-Campbell-Hausdorff expansion of Lemma~\ref{links:curve:lem:bch} has first-order error $3.6\times10^{-2}$, $9.1\times10^{-3}$, $2.3\times10^{-3}$ and second-order error $3.6\times10^{-3}$, $4.6\times10^{-4}$, $5.7\times10^{-5}$ at generator scales $0.1$, $0.05$, $0.025$, scaling as $\varepsilon^2$ and $\varepsilon^3$. The flat routing pattern with $c=(1,0.5,0.2,0.05,0.01)$ has all weights positive, row sums one, and cocycle defect $0$, while uniform causal attention has defect $1/6$. For a generic Gaussian, Eq.~\eqref{links:curve:eq:partialcov} holds to $1.4\times10^{-16}$ and the second identity to $5.3\times10^{-16}$; a generic joint law has $\max|F-I|=3.4$, while an AR(1) chain has $1.2\times10^{-15}$ and $W_{ij}=A^{i-j}$ to $1.1\times10^{-16}$; the covariance law for the local-coordinate regression $W_{ab}\mapsto g_aW_{ab}g_b^{-1}$ holds with distinct per-token frames to $1.3\times10^{-15}$. The reduced-rank excess risk $0.096854$ equals the sum of squared discarded singular values $0.096854$, and no factor perturbation or random rank-$2$ matrix does better. The Fisher expansion $2\mathrm{KL}/(t^2ds^2)$ equals $1.0054$, $1.0027$, $1.0013$ at $t=10^{-2},5\times10^{-3},2.5\times10^{-3}$, converging to one linearly in $t$; the derivative of the exponential map is verified to $8.9\times10^{-9}$ by finite differences; and along the shrinking path the divergence is $2.133,4.499,6.809,9.112$ at $\sigma_{\min}=10^{-1},\dots,10^{-4}$, with increments per decade of $2.37$, $2.31$, $2.30$ approaching $\ln10=2.303$. These checks ground the algebra at one small size and say nothing about trained models.


\section{Gaussian inference and template computation}\label{unified:inference}
\subsection{Projected routing and predictive communication}
\label{inference:p}
These executable projected channels are a present realization for the group-valued construction to reproduce or extend. A rank-deficient head map is not itself an invertible group link.
\subsubsection{States, frames, and encoding}
\label{inference:p:sec:states}

Token position $i$ carries a Gaussian state in local coordinates on a fiber $E_i\cong\mathbb R^K$ and an invertible frame $U_i$ identifying the reference space $\mathsf V=\mathbb R^K$ with $E_i$:
\begin{equation}
q_i=\N(\mu_i,\Sigma_i),\qquad
\xi_i=U_i^{-1}\mu_i,\qquad
C_i=U_i^{-1}\Sigma_iU_i^{-\top}.
\label{inference:p:eq:reference_moments}
\end{equation}
The pair $(\xi_i,C_i)$ is the state in reference coordinates. A passive rechoice of agent $i$'s frame acts by
\begin{equation}
(\mu_i,\Sigma_i,U_i)\mapsto(g_i\mu_i,g_i\Sigma_ig_i^{\top},g_iU_i),\qquad g_i\in\mathrm{GL}(K),
\label{inference:p:eq:passive}
\end{equation}
and leaves $(\xi_i,C_i)$ unchanged. The ideal routing costs and full-covariance messages obey this passive law. A fixed diagonal family and the local-coordinate MLP impose additional restrictions stated below.

\paragraph{Encoding.}
Each token type $c$ has a prior bank entry $(\mu_p^c,\Sigma_p^c,\phi_c)$: a Gaussian embedding in the sense of \citep{inference:vilnis2015} with a Lie-algebra frame coefficient \citep{unified:attention}. Under \cfg{encode\_mode="per\_token"} the input state is the prior of the token, $q_i^{(0)}=p_{x_i}$. The frame is a product of token and positional factors, $U_i=\exp(X_{x_i})\exp(Y_i)$, with $Y_i$ learned when \cfg{pos\_phi="learned"}. Under \cfg{gauge\_group="block\_glk"} both factors are block diagonal with one $d\times d$ block per head, $K=Hd$. The current projected architecture builds one fixed frame per forward pass and shares it across layers. Predictive readout requires \cfg{n\_e\_steps=1}; receiver inference permits additional steps. The former \cfg{e\_phi\_lr} field is retired.

\subsubsection{One layer: comparison and routing}
\label{inference:p:sec:routing}

Head $h$ has full-row-rank comparison maps $S_{Q,h},S_{K,h}:\mathsf V\to\mathbb R^r$, sending map $T_h:\mathsf V\to\mathbb R^d$, and receiving map $O_h:\mathbb R^d\to\mathsf V$. The current implementation fixes $r=d=K/H$ and all maps read or write the full reference state; earlier head-only scopes are retired. Reduced channel rank does not require disjoint content channels \citep{links:vaswani2017}.

\paragraph{Comparison laws.}
The query and key laws are Gaussian pushforwards of the reference states into the comparison space,
\begin{equation}
\begin{aligned}
q_{Q,i}^{h}&=\N\bigl(a_i,V_{Q,i}\bigr),&a_i&=S_{Q,h}\xi_i,&V_{Q,i}&=S_{Q,h}C_iS_{Q,h}^{\top},\\
p_{K,j}^{h}&=\N\bigl(b_j,V_{K,j}\bigr),&b_j&=S_{K,h}\xi_j,&V_{K,j}&=S_{K,h}C_jS_{K,h}^{\top}.
\end{aligned}
\label{inference:p:eq:comparison_laws}
\end{equation}
Both reduced covariances are positive definite because $C_i\succ0$ and the maps have full row rank. The routing cost under \cfg{relational\_routing="gaussian\_kl"} is the exact Gaussian divergence
\begin{equation}
\begin{aligned}
D_{ij}^{h}=\KL\bigl(q_{Q,i}^{h}\|p_{K,j}^{h}\bigr)
=\tfrac12\bigl[&\Tr(P_jV_{Q,i})+(a_i-b_j)^{\top}P_j(a_i-b_j)-r\\
&+\log\det V_{K,j}-\log\det V_{Q,i}\bigr],
\qquad P_j=V_{K,j}^{-1}.
\end{aligned}
\label{inference:p:eq:routing_cost}
\end{equation}
Its mean compatibility contains the bilinear term $a_i^{\top}P_jb_j$, the covariance-weighted analogue of the dot-product score. The terms in $\Tr(P_jV_{Q,i})$ and $\log\det V_{K,j}$ carry covariance information and are pair dependent whenever the key covariances differ across sources. Under \cfg{relational\_routing="noisy\_query"}, available only on the predictive route, the code uses
\begin{equation}
D_{ij}^{h,\mathrm{NQ}}=\tfrac12[(a_i-b_j)^\top W_{ij}^{-1}(a_i-b_j)+\log\det W_{ij}-\log\det V_{Q,i}],\qquad W_{ij}=V_{Q,i}+V_{K,j}.
\label{inference:p:eq:noisy_query}
\end{equation}
Before capping this nonnegative cost differs from the Gaussian negative log likelihood only by a row constant. At unit temperature and inactive cap its row is the noisy-query source posterior; other temperatures and capping modify that declared potential.

\paragraph{Gibbs routing.}
With a source prior $\pi_{ij}^{h}$ and a temperature $\tau_h>0$, the routing weights minimize
$\sum_j\beta_jD_{ij}^{h}+\tau_h\sum_j\beta_j\log(\beta_j/\pi_{ij}^{h})$ on the simplex, with unique solution
\begin{equation}
\beta_{ij}^{h}=\softmax_j\bigl(-D_{ij}^{h}/\tau_h+\log\pi_{ij}^{h}\bigr).
\label{inference:p:eq:gibbs}
\end{equation}
The causal support sets $\pi_{ij}^{h}=0$ for $j>i$. Under \cfg{beta\_attention\_prior="causal\_alibi"} the prior adds a head-dependent linear distance penalty. The implemented temperature is $\tau_h=\kappa_{\beta,h}\sqrt{d_h}$, with equal $d_h=d$ in the admitted architecture. The code bounds the score to $\widehat D=\operatorname{clip}_{[0,D_{\max}]}D$ and maps nonfinite values to configured endpoints before computing $\softmax_j(B_{ij}^h-\widehat D_{ij}^h/\tau_h)$, for unnormalized log-prior $B$. Normalizing $B$ adds only a row constant. The launcher sets $D_{\max}=8K$. Unmodified Gaussian-KL semantics require an inactive cap, while categorical optimality holds for the actual bounded cost. Equation~\eqref{inference:p:eq:gibbs} is an exact posterior at every temperature: attaching to each receiver and head a binary compatibility record with likelihood $\exp(-D_{ij}^{h}/\tau_h)$ makes $\beta_{ij}^{h}$ the posterior of the source label given the record \citep{unified:attention}. This establishes consistency with exact inference; it does not show that attention must be inference.

\subsubsection{One layer: content and the predictive mean}
\label{inference:p:sec:content}

The content maps compose with the frames into the matrix link
\begin{equation}
L_{ij}^{h}=U_iM_hU_j^{-1},\qquad M_h=O_hT_h,
\label{inference:p:eq:link}
\end{equation}
which obeys the endpoint law $L_{ij}^{h}\mapsto g_iL_{ij}^{h}g_j^{-1}$ under Eq.~\eqref{inference:p:eq:passive}. The predictive message and the layer's mean update are
\begin{equation}
m_i=\sum_h\sum_j\beta_{ij}^{h}L_{ij}^{h}\mu_j=U_i\sum_hO_h\sum_j\beta_{ij}^{h}T_h\xi_j,
\qquad
\mu_i^{+}=r\mu_i+m_i,
\label{inference:p:eq:mean}
\end{equation}
with $r=1$ under \cfg{attention\_residual=True} and $r=0$ otherwise. The mean readout is linear in source means at fixed weights; state-dependent routing makes the complete mean map nonlinear. Moment matching adds quadratic dependence on means. At initialization $T_h$ selects head block $h$ and $O_h$ inserts it; only $O_h$ is scaled by $1/\sqrt{2L}$ under \cfg{relational\_output\_init="depth\_scaled"}.

\paragraph{The predictive kernel.}
Equation~\eqref{inference:p:eq:mean} is the mean of a declared kernel. Give each receiver and head an independent source label $J^{h}\sim\beta_i^{h}$ and take
\begin{equation}
k_i'\mid k,J\sim\N\Bigl(rk_i+\sum_hL_{iJ^{h}}^{h}k_{J^{h}},\ Q_i\Bigr),
\label{inference:p:eq:kernel}
\end{equation}
with source states $k_j\sim\N(\mu_j,\Sigma_j)$ independent across $j$ and the self edge included. Averaging over source states and labels gives mean $\mu_i^{+}$. The kernel leaves $T_h$, $O_h$, and $Q_i$ free and declares the unit persistence coefficient, additivity across heads, and one label per head for both routing and content. These are modeling choices recorded as such.

\subsubsection{One layer: the covariance}
\label{inference:p:sec:covariance}

\paragraph{Passthrough}

Under \cfg{predictive\_covariance="passthrough"}, attention returns $\Sigma_i^+=\Sigma_i$. If encoding supplies a token-type covariance and every later seam also passes it through, the stored covariance depends only on token identity. Its reference representation can still depend on position through $U_i$. Historical runs used this contract; the current launcher selects \cfg{moment\_matched}. A run's saved configuration determines its actual contract.

\paragraph{Moment matching}

\begin{proposition}[Moment-matched predictive covariance]
\label{inference:p:prop:covariance}
Under the kernel of Eq.~\eqref{inference:p:eq:kernel}, the covariance of $k_i'$ in reference coordinates is
\begin{equation}
\begin{aligned}
C_i^{\mathrm{att}}={}&\widetilde Q_i
+r\Bigl[C_i+\sum_h\beta_{ii}^{h}\bigl(M_hC_i+C_iM_h^{\top}\bigr)\Bigr]
+\sum_j\sum_{h,h'}c_{ij}^{hh'}M_hC_jM_{h'}^{\top}\\
&+\sum_h\Bigl[\sum_j\beta_{ij}^{h}n_{ij}^{h}n_{ij}^{h\top}-\bar n_i^{h}\bar n_i^{h\top}\Bigr],
\end{aligned}
\label{inference:p:eq:covariance}
\end{equation}
where $\widetilde Q_i=U_i^{-1}Q_iU_i^{-\top}$, $c_{ij}^{hh}=\beta_{ij}^{h}$, $c_{ij}^{hh'}=\beta_{ij}^{h}\beta_{ij}^{h'}$ for $h\ne h'$, $n_{ij}^{h}=M_h\xi_j$, and $\bar n_i^{h}=\sum_j\beta_{ij}^{h}n_{ij}^{h}$. The local covariance is $\Sigma_i^{+}=U_iC_i^{\mathrm{att}}U_i^{\top}$.
\end{proposition}

\begin{proof}
Write $X=rk_i+\sum_hL_{iJ^h}^{h}k_{J^h}+\varepsilon$ with $\varepsilon\sim\N(0,Q_i)$. By the law of total covariance, $\Cov X=\E\Cov(X\mid J)+\Cov\E(X\mid J)$. Given the labels, $X-\varepsilon=\sum_sA_sk_s$ with $A_i=rI+\sum_h\mathbf 1[J^h=i]L_{ii}^{h}$ and $A_s=\sum_h\mathbf 1[J^h=s]L_{is}^{h}$ for $s\ne i$, so $\Cov(X\mid J)=\sum_sA_s\Sigma_sA_s^{\top}+Q_i$ by independence of the sources. Averaging over the independent labels uses $\E(\mathbf 1[J^h=s]\mathbf 1[J^{h'}=s])=c_{is}^{hh'}$ and gives the terms in $C_i$, $\beta_{ii}^{h}$, and $c_{ij}^{hh'}$. For the second term, $\E(X\mid J)=r\mu_i+\sum_hL_{iJ^h}^{h}\mu_{J^h}$ is a sum of independent head contributions, so its covariance is the sum of the per-head spreads. Conjugating by $U_i^{-1}$ removes every frame, since $L_{ij}^{h}\Sigma_jL_{ij}^{h'\top}=U_iM_hC_jM_{h'}^{\top}U_i^{\top}$.
\end{proof}

The proposition is item (2) of the exact-status theorem of the GL(K) manuscript \citep{unified:attention}, in its $r=1$ case. Passthrough retains $C_i$ then; at $r=0$ it is a separate retention contract. Under \cfg{predictive\_covariance="moment\_matched"} the layer returns $\Sigma_i^{+}$ with the reference floor
\begin{equation}
\widetilde Q_i=\epsilon\frac{\Tr C_i}{K}I,\qquad Q_i=\epsilon\frac{\Tr C_i}{K}U_iU_i^{\top},\qquad \epsilon=\cfg{predictive\_covariance\_floor}.
\label{inference:p:eq:floor}
\end{equation}
The trace $\Tr C_i$ is invariant under Eq.~\eqref{inference:p:eq:passive}, so $Q_i$ transforms as a covariance. The kernel leaves $\widetilde Q_i$ free; this choice is made for conditioning. The pointwise lower bound is $\lambda_{\min}(C_i^{\mathrm{att}})\ge\epsilon\Tr(C_i)/K$. This is not a uniform condition-number bound over arbitrary maps or depths. A relative floor $\epsilon C_i$ can compound shrinkage along directions absent from a residual-free content map. The isotropic reference floor fixes a reference Euclidean metric: general reference-basis changes must transform that metric rather than reset it to the identity. Passive site covariance is unaffected.

\paragraph{Properties.}
Four consequences follow directly. First, $C_i^{\mathrm{att}}-\widetilde Q_i$ is a covariance and hence positive semidefinite, so $\Sigma_i^{+}\succeq Q_i\succ0$: the output is positive definite whatever the routing and whatever the rank of $M_h$. Second, $\Sigma_i^{+}$ obeys the passive law $\Sigma_i^{+}\mapsto g_i\Sigma_i^{+}g_i^{\top}$, because $C_i^{\mathrm{att}}$ depends only on reference moments and the invariant weights, and the floor of Eq.~\eqref{inference:p:eq:floor} transforms in the same way. Third, the between-source spread depends on the source means through $\beta$ and $n_{ij}^{h}$, so the covariance becomes context dependent: it grows when a receiver's attended sources carry disagreeing content. Fourth, the cross-head terms $h\ne h'$ are nonzero whenever two heads route to a common source and $T_hC_jT_{h'}^{\top}\ne0$, which occurs for a full covariance or under full input scope. Independent head labels do not make the source states independent when two heads select the same source.

\paragraph{Reductions.}
With the residual off, $C_i^{\mathrm{att}}=\Cov(m_i)+\widetilde Q_i$ in reference coordinates, the predictive reading in which the message replaces the state. With the residual on, the persistence term $C_i$ and the self-edge cross terms appear. There is no monotonicity in depth in general: the cross terms $\beta_{ii}^{h}(M_hC_i+C_iM_h^{\top})$ can be negative. If every $M_h=0$, the output is $r\Sigma_i+Q_i$. If each head routes all weight to one source, the spread term vanishes and the cross-head weights reduce to indicators of a shared source.

\paragraph{Growth with depth under the residual.}
For $\sum_hM_h=I$ and routing only to self, $r=1$ gives $k_i'=2k_i+\varepsilon$ and covariance $4\Sigma_i+Q_i$; $r=0$ gives $\Sigma_i+Q_i$. For $\sum_hM_h=sI$ the respective factors are $(1+s)^2$ and $s^2$. These are exact reductions at specified maps and routing, not predictions for arbitrary trained layers or subsequent MLPs. Predictive moment matching does not apply \cfg{sigma\_max}.

\paragraph{Diagonal families.}
The executable \cfg{gaussian\_full} family stores $\Sigma_i^+$ and \cfg{gaussian\_diagonal} stores its local diagonal. Matching marginal means and variances minimizes $\KL(p\|q)$ over diagonal Gaussian $q$, as follows by separating its cross-entropy by coordinate. This forward-KL projection differs from the reverse-KL receiver optimum. A fixed local-diagonal family is not closed under arbitrary $\mathrm{GL}(K)$ congruences, so the full passive law applies before projection. Earlier \cfg{gaussian\_diagonal\_exact} and \cfg{gaussian\_frame\_diagonal} names are retired.

\paragraph{Value frame.}
Content and covariance use the same frame, the value frame, which coincides with the routing frame except under rotary position encoding with \cfg{rope\_on\_value=False} for the present Gaussian families.

\paragraph{Computation.}
The implementation (\cfg{RelationalOperators.predictive\_covariance\_update}) forms the reference covariances $C_j$, pushes them into the value space per head pair, $Y_j^{hh'}=\hat T_hC_j\hat T_{h'}^{\top}$, pools them with the weights $c_{ij}^{hh'}$, adds the pooled second moment of $v_j^{h}=\hat T_h\xi_j$ minus the outer product of its mean on the head diagonal, lifts by $\hat O_h$, adds the residual terms, and conjugates by the value frame. Here $\hat T_h\in\mathbb R^{d\times K}$ and $\hat O_h\in\mathbb R^{K\times d}$ embed the head-scope maps in the full state space, the full-state maps now used throughout. Before computing spread the code subtracts token 0 from all source means; a common shift preserves the covariance because each routing row sums to one. The pooling is an attention with $d\times d$ values per head pair, so the dominant costs are $O(HdK^2+H^2d^2K)$ per token for the pushforward, $O(N^2H^2d^2)$ per sequence for the pooling, and $O(K^3)$ per token for the frame conjugations. All operations are differentiable, so gradients flow to the source means and covariances, the routing, the content maps, and the frame.

\subsubsection{Between layers}
\label{inference:p:sec:block}

After the attention step, a block applies optional seams in a fixed order: head mixer, cross-type coupling, block normalization, block MLP. The active settings are recorded below. The normalizations act on the mean only. The block MLP (\cfg{block\_mlp\_mode}) is a coordinate GELU MLP by default \citep{links:hendrycks2016gelu}. It passes the covariance through unless \cfg{block\_mlp\_covariance="delta\_full"} applies the delta-method pushforward $J\Sigma J^{\top}$ with a covariant floor. Under \cfg{"gaussian\_template"} it is replaced by template routing, a derived Gibbs readout of learned Gaussian templates that reads the diagonal of the covariance \citep{links:dennis2026template}. Moment-matched attention can supply context-dependent uncertainty to templates, as can a gradient receiver. The stack passes each block's output onward and applies its prior-handoff policy. The predictive readout does not read the prior.

\subsubsection{Decoder and objective}
\label{inference:p:sec:decode}

Under \cfg{decode\_mode="linear"}, the next-token law is $\softmax_v(W_{\mathrm{out}}\mu_i^{(L)}+b+\kappa_{\mathrm{uni}}\log p_{\mathrm{uni}})$, a Gibbs readout of shared-covariance templates that reads the mean only. Under the prior-bank decoders (\cfg{diagonal\_chunked}, \cfg{full\_chunked}, \cfg{family\_chunked}), the readout is
\begin{equation}
\hat p_i(v)\propto\omega_v\exp\bigl[-\KL(q_i^{(L)}\|p_v)/\tau_{\mathrm{dec}}\bigr],
\label{inference:p:eq:decode}
\end{equation}
whose exponent contains $\Tr(\Sigma_{p,v}^{-1}\Sigma_i^{(L)})$. The decoder therefore reads the final covariance. Under passthrough this is the input-token covariance; under moment matching it is the context-dependent covariance of Proposition~\ref{inference:p:prop:covariance}, so an uncertain final belief is scored against each class template with its full spread. Training minimizes next-token cross-entropy through every differentiable component. An optional M-step self-coupling term $\sum_i\alpha_i\KL(q_i^{*}\|p_i)$ is weighted by \cfg{mstep\_self\_coupling\_weight}, which is zero in the active launcher snapshot. No E-step gradient descent runs on this route: the layer update is the readout of Eqs.~\eqref{inference:p:eq:mean} and, optionally,~\eqref{inference:p:eq:covariance}.



\subsection{Receiver inference over the same channels}
\label{inference:g}
\paragraph{The receiver episode.}
A projected gradient layer snapshots the incoming states $(\bar\xi_j,\bar C_j)$, the routing weights $\beta$, and the frames once. It then moves each receiver's state $(\xi,C)$ for $n_{\mathrm e}$ iterations against a receiver energy $F_i(\xi,C)$ with that snapshot held fixed in value, while the outer supervised graph through the snapshot remains attached. The frames do not change within the episode, so the map $(\mu,\Sigma)\mapsto(\xi,C)$ is a fixed linear reparametrization.

\begin{lemma}[Reference coordinates suffice]
\label{inference:g:lem:reference}
Every energy below is a function of reference moments. Its local gradients are $\nabla_\mu F=U_i^{-\top}\nabla_\xi F$ and $\nabla_\Sigma F=U_i^{-\top}(\nabla_CF)U_i^{-1}$. The Gaussian Fisher natural gradient $(\Sigma\nabla_\mu F,2\Sigma\nabla_\Sigma F\Sigma)$ equals $U_i$ applied to $(C\nabla_\xi F,2C\nabla_CF C)$, with $U_i^\top$ on the right of the covariance term. Every Fisher step taken in local coordinates is therefore the same step taken in reference coordinates and mapped back.
\end{lemma}

\begin{proof}
Substitute $\xi=U_i^{-1}\mu$ and $C=U_i^{-1}\Sigma U_i^{-\top}$ into the chain rule. The Fisher metric of the Gaussian family is invariant under affine reparametrization of the sample space \citep{frame:Amari1998}.
\end{proof}

From here on we work in reference coordinates and suppress the receiver index. These unrestricted full-family identities do not make a local-diagonal family reference diagonal. The local quantities used by the code are the pullbacks $\Lambda_{\mathrm{loc}}=U_i^{-\top}\Lambda U_i^{-1}$ and $\eta_{\mathrm{loc}}=U_i^{-\top}\eta$.

\paragraph{One quadratic family.}
Every receiver energy in this section has the form
\begin{equation}
F(\xi,C)=\tfrac12\Tr(\Lambda C)+\tfrac12\xi^\top\Lambda\xi-\eta^\top\xi-\tfrac{c}{2}\log\det C+\text{const},
\qquad \Lambda\succ0,\ c>0.
\label{inference:g:eq:quadratic_family}
\end{equation}

\begin{lemma}[Gradients, minimizer, Fisher step]
\label{inference:g:lem:family}
For Eq.~\eqref{inference:g:eq:quadratic_family}:
\begin{enumerate}
\item The gradients are $\nabla_\xi F=\Lambda\xi-\eta$ and $\nabla_CF=\tfrac12(\Lambda-c C^{-1})$. For a diagonal family $C=\diag(s)$, they are $\partial_{s_k}F=\tfrac12(\Lambda_{kk}-c/s_k)$, with the dense mean force unchanged.
\item The unique minimizer is $\xi^\star=\Lambda^{-1}\eta$ and $C^\star=c\Lambda^{-1}$. In the diagonal family it is $\xi^\star=\Lambda^{-1}\eta$ and $s_k^\star=c/\Lambda_{kk}$.
\item The explicit Fisher mean step $\xi\leftarrow\xi-\varepsilon C(\Lambda\xi-\eta)$ maps the error $e=\xi-\xi^\star$ to $(I-\varepsilon C\Lambda)e$. Since $C\Lambda$ is similar to the positive definite $C^{1/2}\Lambda C^{1/2}$, for fixed $C$ the step contracts in the $C^{-1}$ norm if and only if $\varepsilon\lambda_{\max}(C\Lambda)<2$. At the covariance optimum $C\Lambda=c I$, so the bound is $\varepsilon<2/c$, and $\varepsilon=1/c$ is the exact Newton step.
\end{enumerate}
\end{lemma}

\begin{proof}
Items 1 and 2 follow by differentiation, using $\partial_C\log\det C=C^{-1}$ and strict convexity of $F$ in $(\xi,C)$. For item 3, $e^+=e-\varepsilon C\Lambda e$, and $C\Lambda=C^{1/2}(C^{1/2}\Lambda C^{1/2})C^{-1/2}$ has the same real positive spectrum as $C^{1/2}\Lambda C^{1/2}$.
\end{proof}

The parameter $c$ is the positive coefficient with which the receiver's entropy $\Hent(q)=\tfrac12\log\det C+\text{const}$ enters the energy. The rest of this section is about what $\Lambda$, $\eta$, and $c$ should be.

\subsubsection{Ambient alignment as a weighted divergence energy}
\label{inference:g:sec:ambient}

The selectable ambient receiver uses, with $w_{hj}=\lambda_h\beta_{ij}^h$, $M_h=O_hT_h$, and receiving noise $\nu>0$,
\begin{equation}
\begin{split}
F^{\mathrm A}(\xi,C)&=\alpha\KL\bigl(\N(\xi,C)\Vert\N(\xi_p,C_p)\bigr)
+\sum_{h,j}w_{hj}\KL\bigl(\N(\xi,C)\Vert\N(M_h\bar\xi_j,R^{\mathrm A}_{hj})\bigr),\\
R^{\mathrm A}_{hj}&=M_h\bar C_jM_h^\top+\nu I_K.
\end{split}
\label{inference:g:eq:ambient}
\end{equation}
We call this ambient alignment, because each head's message is lifted to a $K$-dimensional law and the receiver is aligned to it by a divergence. It is the direct replacement of $\Omega_{ij}$ by $L_{ij}^{(h)}$ in the old alignment energy.

\begin{proposition}[Ambient alignment is the quadratic family with $c=\alpha+\sum_h\lambda_h$]
\label{inference:g:prop:ambient}
Eq.~\eqref{inference:g:eq:ambient} has the form of Eq.~\eqref{inference:g:eq:quadratic_family} with
\begin{equation}
\Lambda=\alpha C_p^{-1}+\sum_{h,j}w_{hj}R^{\mathrm A -1}_{hj},\qquad
\eta=\alpha C_p^{-1}\xi_p+\sum_{h,j}w_{hj}R^{\mathrm A -1}_{hj}M_h\bar\xi_j,\qquad
c=\alpha+\sum_h\lambda_h.
\end{equation}
\end{proposition}

\begin{proof}
For any Gaussian target $t$, $\KL(q\Vert t)=-\E_q[\log t]-\Hent(q)$. Summing with weights gives $\sum w_{hj}\KL(q\Vert t_{hj})=\sum w_{hj}(-\E_q\log t_{hj})-(\sum w_{hj})\Hent(q)$. Each $-\E_q\log t_{hj}$ is $\tfrac12\Tr(R^{-1}C)+\tfrac12(\xi-m)^\top R^{-1}(\xi-m)$ plus a constant. On the causal support $\sum_j\beta_{ij}^h=1$, so $\sum_{h,j}w_{hj}=\sum_h\lambda_h$. The prior contributes $\alpha$ to $c$.
\end{proof}

\begin{lemma}[The receiving noise is required]
\label{inference:g:lem:noise_required}
If $d<K$ and $\nu=0$, then $R^{\mathrm A}_{hj}$ has rank at most $d$, and $\KL(\N(\xi,C)\Vert\N(M_h\bar\xi_j,R^{\mathrm A}_{hj}))=+\infty$ for every $C\succ0$.
\end{lemma}

\begin{proof}
$M_h\bar C_jM_h^\top$ has range in $\im O_h$, of dimension at most $d$. A nondegenerate $K$-dimensional Gaussian is not absolutely continuous with respect to a law supported on a proper affine subspace.
\end{proof}

The receiving noise in the implementation therefore exists to make the energy finite, which is how the design report states it. It is not a collapse regularizer, although a positive floor would also play that role.

\begin{proposition}[Complement pinning]
\label{inference:g:prop:pinning}
Let $V_h=\im O_h$ and let $V_h^\perp$ be its Euclidean orthogonal complement. The law $\N(M_h\bar\xi_j,R^{\mathrm A}_{hj})$ has mean component $0$ and covariance $\nu I$ on $V_h^\perp$, independently of the source.
\end{proposition}

\begin{proof}
$M_h\bar\xi_j\in V_h$. The symmetric matrix $M_h\bar C_jM_h^\top$ has range in $V_h$, so it annihilates $V_h^\perp$, and $R^{\mathrm A}_{hj}$ acts on $V_h^\perp$ as $\nu I$.
\end{proof}

At block-identity initialization with unit head weights, $M_h$ is the coordinate projector onto block $h$. Every block $g$ then receives precision $(H-1)/\nu$ toward zero from the other heads and precision about $1/(C_{gg}+\nu)$ toward its routed content. At $C=15I$, $\nu=1$, and $H=2$ these are $1$ and $1/16$, so the frozen minimizer keeps about $6.5\%$ of the incoming mean. This is the earlier source note's $K=20$ illustration, not a new training result. By Lemma~\ref{inference:g:lem:family}, the same numbers give $C\Lambda\approx16.9 I$ at initialization and a stable Euler range $\varepsilon<0.12$, while at the covariance optimum the range is $\varepsilon<2/c=2/3$. At $\varepsilon=0.9$ the unrolled step $I-\varepsilon C\Lambda$ has eigenvalues near $-14$, which reverses and amplifies the outer gradient through each layer.

\begin{remark}
These three properties follow from one design choice: the message is lifted to a $K$-dimensional law before the comparison. Lifting a rank-$d$ message forces a noise floor for finiteness (Lemma~\ref{inference:g:lem:noise_required}). That floor makes every head an opinion about all $K$ coordinates (Proposition~\ref{inference:g:prop:pinning}). Each divergence also brings its own copy of the receiver's entropy (Proposition~\ref{inference:g:prop:ambient}).
\end{remark}

\subsubsection{Two entropies}
\label{inference:g:sec:entropy}

The question whether the gradient should contain an entropy term has two distinct answers, because two entropies are in play.

\paragraph{The receiver entropy.}
A normalized variational free energy has one receiver entropy \citep{unified:blei2017}. Ambient alignment has coefficient $c=\alpha+\sum_h\lambda_h$, but is a valid weighted-divergence objective: after division by $c$ it is KL to the normalized geometric opinion pool
\begin{equation}
p_{\mathrm A}^*(x)\propto p(x)^{\alpha/c}\prod_{h,j}t_{hj}(x)^{w_{hj}/c}.
\label{inference:g:eq:opinion_pool}
\end{equation}
Its covariance $c\Lambda^{-1}$ is not an algebraic error. It differs from a posterior interpreting the same $\Lambda$ as accumulated likelihood precision. The channel model declares that alternative meaning; entropy coefficients must be interpreted with the complete objective and its scaling.

\paragraph{The routing entropy.}
The term $\tau\sum_j\beta_j\log(\beta_j/\pi_j)$ makes the Gibbs weights the minimizer of the headwise routing objective of \citep{unified:attention}. Whether it belongs in the receiver gradient depends on whether routing and content are inferred jointly.

\begin{definition}[Joint routing]
\label{inference:g:def:joint}
For a content cost $\ell_{hj}(q)$ with weight $\gamma_h\ge0$, the joint receiver functional is
\begin{equation}
F^{\mathrm J}(q,\beta)=\alpha\KL(q\Vert p)+\sum_h\Bigl[\sum_j\beta_j^h\bigl(D_{j}^h(q)+\gamma_h\ell_{hj}(q)\bigr)+\tau_h\sum_j\beta_j^h\log\frac{\beta_j^h}{\pi_j^h}\Bigr].
\label{inference:g:eq:joint}
\end{equation}
\end{definition}

\begin{proposition}[Joint routing sees content, and its gradient carries the comparison force]
\label{inference:g:prop:joint}
The minimizing weights of Eq.~\eqref{inference:g:eq:joint} are $\beta_j^{h\star}=\softmax_j(\log\pi_j^h-(D_j^h+\gamma_h\ell_{hj})/\tau_h)$. The reduced functional is
\[
\alpha\KL(q\Vert p)-\sum_h\tau_h\log\sum_j\pi_j^h\exp\bigl(-(D_j^h+\gamma_h\ell_{hj})/\tau_h\bigr).
\]
Its gradient in $q$ is $\alpha\nabla\KL(q\Vert p)+\sum_{h,j}\beta_j^{h\star}\nabla(D_j^h+\gamma_h\ell_{hj})$. The routing entropy contributes no separate term, because its effect is that the derivative through $\beta^\star$ vanishes (envelope theorem). The comparison force is
\begin{equation}
\nabla_\xi D_j^h=S_{Q,h}^\top B_j^{h -1}(a^h-b_j^h),\qquad
\nabla_CD_j^h=\tfrac12S_{Q,h}^\top\bigl(B_j^{h -1}-A^{h -1}\bigr)S_{Q,h},
\label{inference:g:eq:comparison_force}
\end{equation}
where $-\tfrac12\log\det A^h$ inside $D$ is a third entropy, that of the query law.
\end{proposition}

\begin{proof}
Minimize the convex objective in $\beta$ on the simplex with a normalization multiplier. Substitute $\beta^\star$, and differentiate the log-sum-exp. Equation~\eqref{inference:g:eq:comparison_force} differentiates $\KL(\N(a,A)\Vert\N(b,B))$ with $a=S_Q\xi$ and $A=S_QCS_Q^\top$.
\end{proof}

\begin{proposition}[No content-informed mixture functional routes by $D$ alone]
\label{inference:g:prop:no_D_only}
Consider a functional that is affine in $\beta$ except for the routing entropy, with linear coefficients $D_j^h+g_{hj}(q)$. Its minimizing weights equal Eq.~\eqref{inference:p:eq:gibbs} for every $q$ if and only if $g_{hj}(q)$ does not depend on $j$ on the positive-prior support.
\end{proposition}

\begin{proof}
The minimizer is $\softmax_j(\log\pi_j-(D_j+g_{hj})/\tau)$. A softmax is unchanged by $j$-independent shifts of its arguments and is injective modulo them.
\end{proof}

A content term that carries source-specific information therefore changes the routing. Routing by $D$ alone, while letting content move the receiver, is not coordinate descent of any one such functional. It is, however, exact inference in a different declared model.

\begin{definition}[Fixed routing and the cut interpretation]
\label{inference:g:def:cut}
Comparison supplies a source-label posterior $\beta$ whose values remain fixed during the receiver episode. Content does not feed back into that module. This realizes the no-feedback intention of modular or cut inference \citep{inference:plummer2015cuts,inference:jacob2017modules}, but does not alone specify a downstream receiver posterior.
\end{definition}
A literal categorical cut posterior averages normalized conditionals, $p_{\mathrm{cut}}(x)=\sum_J\beta(J)p(x\mid J,y)$, and is generally a mixture. Marginalizing an ordinary likelihood first instead gives $p(x)\sum_J\beta(J)p(y\mid x,J)$ and allows evidence to reweight source labels. The implemented expected-log-likelihood receiver uses a third construction: for $\alpha>0$, its target is the normalized geometric likelihood pool
\begin{equation}
p_{\mathrm{pow}}(x)\propto p(x)\prod_{h,j}p(y_{hj}\mid x)^{w_{hj}/\alpha}.
\label{inference:g:eq:power_posterior}
\end{equation}
Indeed, $\alpha\KL(q\|p)+\sum_{h,j}w_{hj}\E_q[-\log p(y_{hj}\mid x)]=\alpha\KL(q\|p_{\mathrm{pow}})+\text{const}$. This is exact inference for declared likelihood powers when normalizable, and the state-coordinate optimum of a mean-field functional whose label factor is fixed. It is not generally the marginalized categorical cut posterior. Fixing $\beta$ removes both its derivative and the comparison force from the receiver update; the outer supervised graph through its incoming snapshot stays attached. Recomputing routing during the episode requires a separate objective argument. No complete-stack population ELBO follows from these local identities.

\subsubsection{The channel-likelihood energy}
\label{inference:g:sec:channel}

The partial-law reading of Eq.~\eqref{inference:p:eq:comparison_laws} extends from comparison to content. Head $h$ observes the receiver through a receiving measurement $P_h\in\mathbb R^{d\times K}$. The sender's message $T_hx_j$, with $x_j\sim\N(\bar\xi_j,\bar C_j)$ marginalized, predicts that measurement through a channel with noise $Q_h\succeq0$:
\begin{equation}
\begin{split}
&P_hx=T_hx_j+\epsilon,\qquad \epsilon\sim\N(0,Q_h)\\
&\Longrightarrow\quad
p(T_h\bar\xi_j\mid x)=\N\bigl(T_h\bar\xi_j;P_hx,R_{hj}\bigr),\qquad R_{hj}=T_h\bar C_jT_h^\top+Q_h.
\end{split}
\label{inference:g:eq:channel_model}
\end{equation}
With frozen routing weights declared as likelihood powers $w_{hj}=\lambda_h\beta_{ij}^h$, the receiver's variational free energy is
\begin{equation}
F^{\mathrm B}(\xi,C)=\alpha\KL\bigl(\N(\xi,C)\Vert\N(\xi_p,C_p)\bigr)+\sum_{h,j}w_{hj} \E_{\N(\xi,C)}\bigl[-\log\N(T_h\bar\xi_j;P_hx,R_{hj})\bigr].
\label{inference:g:eq:channel}
\end{equation}

\begin{proposition}[Channel likelihood is the quadratic family with $c=\alpha$]
\label{inference:g:prop:channel}
Eq.~\eqref{inference:g:eq:channel} has the form of Eq.~\eqref{inference:g:eq:quadratic_family} with
\begin{equation}
\Lambda=\alpha C_p^{-1}+\sum_hP_h^\top W_hP_h,\qquad
\eta=\alpha C_p^{-1}\xi_p+\sum_hP_h^\top\omega_h,\qquad
c=\alpha,
\label{inference:g:eq:channel_stats}
\end{equation}
where $W_h=\sum_jw_{hj}R_{hj}^{-1}$ and $\omega_h=\sum_jw_{hj}R_{hj}^{-1}T_h\bar\xi_j$ are pooled in $d\times d$ and $d$ dimensions.
\end{proposition}

\begin{proof}
$\E[-\log\N(y;Px,R)]=\tfrac12(y-P\xi)^\top R^{-1}(y-P\xi)+\tfrac12\Tr(P^\top R^{-1}PC)+\text{const}$. The expectation contains no entropy of $q$, which enters only through the prior's divergence.
\end{proof}

\begin{corollary}[Properties]
\label{inference:g:cor:channel}
\begin{enumerate}
\item With $\alpha=1$, the minimizer $\xi^\star=\Lambda^{-1}\eta$, $C^\star=\Lambda^{-1}$ is the exact Gaussian power posterior of the prior under the likelihoods of Eq.~\eqref{inference:g:eq:channel_model}, each raised to the power $w_{hj}$. This is the conjugate linear-Gaussian update \citep[Sec.~2.3.3]{inference:bishop2006}, and the Gaussian variational optimum is exact here \citep{inference:opper2009variational}.
\item The energy is finite with $Q_h=0$ whenever $T_h$ has full row rank, since then $T_h\bar C_jT_h^\top\succ0$. The channel noise is then a declared floor that guards against loss of rank in $T_h\bar C_jT_h^\top$ as $T_h$ is learned. It is not needed for finiteness.
\item The likelihoods are constant along $V_0=\bigcap_h\ker P_h$, the directions no head observes. The posterior law of the $V_0$ component, conditional on the observed component, therefore equals the prior's conditional law. If $C_p$ does not couple $V_0$ to its complement, the posterior marginal on $V_0$ equals the prior marginal. No head imposes a value in $V_0$, so there is no complement pinning.
\item At the covariance optimum $C\Lambda=\alpha I$, so the explicit Fisher step is stable for $\varepsilon<2/\alpha$, and $\varepsilon=1$ is the exact Newton step when $\alpha=1$.
\item The pooled statistics cost $O(BN^2Hd^2)$ and the source factorizations $O(BNHd^3)$. The corresponding ambient costs are $O(BN^2HK^2)$ and $O(BNHK^3)$, so the channel form is a factor $H^2$ cheaper in pooling and $H^3$ in factorization.
\end{enumerate}
\end{corollary}

\begin{proof}
Item 1 follows from Proposition~\ref{inference:g:prop:channel} and Lemma~\ref{inference:g:lem:family} with $c=1$, since the information-form product of Gaussians has precision $\Lambda$ and information vector $\eta$. Item 2 follows from $y^\top T\bar CT^\top y=(T^\top y)^\top\bar C(T^\top y)>0$ for $y\ne0$. For item 3, each likelihood depends on $x$ only through $P_hx$, so it is invariant under $x\mapsto x+v$ for $v\in V_0$. Bayes' rule then leaves the prior's conditional law along $V_0$ unchanged. Items 4 and 5 are Lemma~\ref{inference:g:lem:family} and direct operation counts.
\end{proof}

\paragraph{The role of $O_h$.}
In Eq.~\eqref{inference:g:eq:channel} the write-back is not a free map. The posterior update moves the state along the gain $\Lambda^{-1}P_h^\top W_h=(C^\star/\alpha)P_h^\top W_h$, as in a Kalman update. The receiving interface is the measurement $P_h$. One natural choice ties $P_h=O_h^\top$, which keeps the existing parameter names and the block-identity initialization and gives the predictive limit below. An untied learned $P_h$ is the other choice.

\begin{proposition}[Predictive limit]
\label{inference:g:prop:predictive_limit}
Assume $P_h=O_h^\top$ with $O_h^\top O_g=\delta_{hg}I_d$, so that $\sum_hO_hO_h^\top=I_K$. Assume also $R_{hj}=R_h$ for all $j$, assume $\lambda_h>0$ for every head, and let both $\alpha C_p^{-1}\to0$ and $\alpha C_p^{-1}\xi_p\to0$. Then
\begin{equation}
\xi^\star=\sum_hO_h\sum_j\beta_{ij}^hT_h\bar\xi_j,
\end{equation}
the residual-free predictive head sum of \citep{links:dennis2026projpred} with the tied orthonormal receiving map.
\end{proposition}

\begin{proof}
In the limit, $\Lambda=\sum_hO_h(\lambda_hR_h^{-1})O_h^\top$ and $\eta=\sum_hO_h\lambda_hR_h^{-1}\sum_j\beta_{ij}^hT_h\bar\xi_j$. Both are block diagonal in the orthonormal basis given by the columns of the $O_h$. Solving blockwise gives $O_h^\top\xi^\star=\sum_j\beta_{ij}^hT_h\bar\xi_j$, using $\sum_j\beta_{ij}^h=1$.
\end{proof}

Each relaxed hypothesis has a definite meaning:
\begin{itemize}
\item \emph{Unequal $R_{hj}$.} The head's channel mean becomes the precision-weighted average $(\sum_j\beta_jR_{hj}^{-1})^{-1}\sum_j\beta_jR_{hj}^{-1}T_h\bar\xi_j$, so confident senders count more than their routing weight alone.
\item \emph{A finite prior.} The posterior blends the prior mean and the messages with a gain set by their precisions, rather than adding them. The additive residual $\mu_i+m_i$ of the predictive route has no exact counterpart in a posterior. Reproducing it would require a declared increment semantics for the messages, which this section does not adopt.
\item \emph{Overlapping heads.} The heads' measurements are fused by least squares.
\end{itemize}
The posterior covariance $C^\star=\alpha\Lambda^{-1}$ depends on context. It omits the between-source spread of the moment-matched predictive covariance, because geometric likelihood pooling differs from marginalizing a categorical predictive mixture.

\begin{remark}[Channel alignment]
Replacing the expected likelihood in Eq.~\eqref{inference:g:eq:channel} by the divergence $\KL(\N(P_h\xi,P_hCP_h^\top)\Vert\N(T_h\bar\xi_j,R_{hj}))$ in the value space gives a channel-alignment energy that parallels the comparison divergence $D$. Its mean terms, and therefore $\xi^\star=\Lambda^{-1}\eta$, coincide with the channel likelihood. Each divergence, however, adds $-\tfrac{w_{hj}}{2}\log\det(P_hCP_h^\top)$. That adds one projected entropy per edge, and its covariance optimum has no closed form when heads overlap. The channel likelihood is preferred for its single entropy and its exact posterior.
\end{remark}

\subsubsection{Updates and outer gradients}
\label{inference:g:sec:updates}

Both energies are in the family of Eq.~\eqref{inference:g:eq:quadratic_family}. For frozen covariance the explicit Euler mean step has the stability restriction of Lemma~\ref{inference:g:lem:family}, namely $\varepsilon\lambda_{\max}(C\Lambda)<2$. Two better updates are available.

\begin{proposition}[Exact natural-parameter flow]
\label{inference:g:prop:flow}
The Fisher flow of Eq.~\eqref{inference:g:eq:quadratic_family} in the full family, $\dot\xi=-C(\Lambda\xi-\eta)$ and $\dot C=c C-C\Lambda C$, has the exact solution
\begin{equation}
J(t)=(1-\rho)J(0)+\rho\Lambda/c,\qquad h(t)=(1-\rho)h(0)+\rho\eta/c,\qquad \rho=1-e^{-c t},
\end{equation}
in the natural parameters $J=C^{-1}$ and $h=C^{-1}\xi$. For every $\rho\in[0,1]$ the iterate is SPD and the energy does not increase. At $\rho=1$ it is the minimizer $(\Lambda^{-1}\eta,c\Lambda^{-1})$.
\end{proposition}

\begin{proof}
$\dot J=-J\dot CJ=\Lambda-c J$ and $\dot h=\dot J\xi+J\dot\xi=\eta-c h$, both linear with constant coefficients. A convex combination of SPD matrices is SPD. The energy decreases along the flow, with derivative $-\nabla_\xi F^\top C\nabla_\xi F-2\Tr(C\nabla_CF C\nabla_CF)\le0$. This is the derivation in the design report's variational notes, restated here for general $c$.
\end{proof}

With the context fixed in value, the closed-form minimizer ($\rho=1$) completes the receiver episode in one solve, and further iterations leave it unchanged. Its outer derivatives follow by implicit differentiation:
\begin{equation}
d\xi^\star=\Lambda^{-1}\bigl(d\eta-d\Lambda \xi^\star\bigr),\qquad dC^\star=(dc)\Lambda^{-1}-c\Lambda^{-1}d\Lambda\Lambda^{-1}.
\end{equation}
Automatic differentiation through a Cholesky solve computes both. This removes that repeated Euler overshoot without bounding all outer derivatives of the complete network.

For the diagonal family, the closed form $\xi^\star=\Lambda^{-1}\eta$, $s^\star=c/\diag\Lambda$ is exact. The code uses local $\Lambda_{\mathrm{loc}}$ and $\eta_{\mathrm{loc}}$ for its local-diagonal family; diagonalizing reference precision would solve a different restricted problem. Any partial damping toward it is a declared choice, because interpolating natural parameters does not remain in the diagonal family when $\Lambda$ is dense. An implicit mean step $\xi^+=\xi-\varepsilon(I+\varepsilon C\Lambda)^{-1}C(\Lambda\xi-\eta)$ is unconditionally stable for every $\varepsilon>0$ in both families, and it is the minimal change to the existing step.

\subsubsection{Profiled state-dependent self-coupling}
\label{inference:g:sec:profiled}

The coefficient $\alpha$ on $\KL(q\Vert p)$ has so far been a constant. The self-coupling registry also offers state-dependent forms, in which $\alpha$ is optimized against the regularizer $R(\alpha)=b_0\alpha-c_0\log\alpha$ with $b_0,c_0>0$. In this section $p$ is the layer prior handed to the receiver, $D(q)=\KL(q\Vert p)$, and each receiver energy of Sections~\ref{inference:g:sec:ambient} and \ref{inference:g:sec:channel} is split as
\begin{equation}
F(q;\alpha)=\alpha D(q)+G(q),
\label{inference:g:eq:split}
\end{equation}
where $G$ is the energy at $\alpha=0$. By Propositions~\ref{inference:g:prop:ambient} and \ref{inference:g:prop:channel}, $G$ has the form of Eq.~\eqref{inference:g:eq:quadratic_family} with coefficients $\Lambda_G$ and $\eta_G$ and entropy multiplicity $c_G=\sum_h\lambda_h$ under ambient alignment and $c_G=0$ under the channel likelihood. The divergence $D$ is invariant under the affine change of coordinates of Lemma~\ref{inference:g:lem:reference}, so the per-token statements below hold in reference and local coordinates alike.

\begin{lemma}[Profiling the coefficient]
\label{inference:g:lem:profile}
For $D\ge0$ the function $\alpha\mapsto\alpha D+R(\alpha)$ is strictly convex on $\alpha>0$, with unique minimizer $\alpha^\star(D)=c_0/(b_0+D)$ and minimum value
\begin{equation}
\min_{\alpha>0}\bigl[\alpha D+b_0\alpha-c_0\log\alpha\bigr]=c_0\log(b_0+D)+c_0-c_0\log c_0.
\end{equation}
The profiled receiver energy is therefore
\begin{equation}
E(q)=\min_{\alpha>0}\bigl[F(q;\alpha)+R(\alpha)\bigr]=c_0\log\bigl(b_0+D(q)\bigr)+G(q)+\text{const}.
\label{inference:g:eq:profiled}
\end{equation}
\end{lemma}

\begin{proof}
The second derivative in $\alpha$ is $c_0/\alpha^2>0$, and the first derivative $D+b_0-c_0/\alpha$ vanishes at $\alpha^\star$. Substituting $\alpha^\star$ gives $c_0D/(b_0+D)+b_0c_0/(b_0+D)+c_0\log(b_0+D)-c_0\log c_0$, and the first two terms sum to $c_0$.
\end{proof}

\begin{proposition}[Envelope force]
\label{inference:g:prop:envelope}
The gradient of Eq.~\eqref{inference:g:eq:profiled} is
\begin{equation}
\nabla E(q)=\alpha^\star\bigl(D(q)\bigr)\nabla D(q)+\nabla G(q).
\end{equation}
It equals the gradient at $q$ of the quadratic-family energy $F(\cdot;\alpha)$ with $\alpha$ frozen at $\alpha^\star(D(q))$, whose coefficients are
\begin{equation}
\Lambda=\Lambda_G+\alpha^\star C_p^{-1},\qquad \eta=\eta_G+\alpha^\star C_p^{-1}\xi_p,\qquad c=\alpha^\star+c_G.
\label{inference:g:eq:frozen_alpha}
\end{equation}
\end{proposition}

\begin{proof}
The chain rule gives $\nabla[c_0\log(b_0+D)]=c_0(b_0+D)^{-1}\nabla D=\alpha^\star\nabla D$. Equivalently, the derivative through $\alpha^\star$ in Lemma~\ref{inference:g:lem:profile} is multiplied by $\partial_\alpha[\alpha D+R(\alpha)]$, which vanishes at $\alpha^\star$. Eq.~\eqref{inference:g:eq:frozen_alpha} is Proposition~\ref{inference:g:prop:ambient} or Proposition~\ref{inference:g:prop:channel} with $\alpha=\alpha^\star$.
\end{proof}

The explicit Fisher step on $E$ therefore needs no new kernel. At every iteration the receiver recomputes $\alpha^\star$ from its current state and the layer prior, forms Eq.~\eqref{inference:g:eq:frozen_alpha}, and applies the step of Lemma~\ref{inference:g:lem:family}. The force is exact as a value, but its derivative is not that of the frozen problem. The Hessian of the self term is
\begin{equation}
\nabla^2\bigl[c_0\log(b_0+D)\bigr]=\alpha^\star\nabla^2D-\frac{(\alpha^\star)^2}{c_0}\nabla D\nabla D^\top,
\end{equation}
so the unrolled outer derivative must retain the dependence of $\alpha^\star$ on the state and on the parameters, and the implementation does not detach it. The rank-one term is negative semidefinite. The profiled energy need not be convex, and the stability range of Lemma~\ref{inference:g:lem:family} describes the frozen problem only.

\begin{proposition}[The closed form is a majorize-minimize step]
\label{inference:g:prop:mm}
Let $q_t$ be the current receiver, $D_t=D(q_t)$, and $\alpha_t=\alpha^\star(D_t)$. The function
\begin{equation}
M_t(q)=c_0\log(b_0+D_t)+\alpha_t\bigl(D(q)-D_t\bigr)+G(q)
\end{equation}
satisfies $M_t\ge E$ with equality at $q_t$, and $M_t(q)=F(q;\alpha_t)+\text{const}$. Let $q_{t+1}$ be the minimizer of $F(\cdot;\alpha_t)$ over the full or the diagonal family, given by Lemma~\ref{inference:g:lem:family} with the coefficients of Eq.~\eqref{inference:g:eq:frozen_alpha} at $\alpha^\star=\alpha_t$. Then $E(q_{t+1})\le E(q_t)$. Equality holds only if $q_t$ is a stationary point of $E$ within the family.
\end{proposition}

\begin{proof}
The function $x\mapsto c_0\log(b_0+x)$ is concave, so it lies below its tangent at $D_t$, whose slope is $\alpha_t$. This gives $M_t\ge E$ with equality at $q_t$. Hence $E(q_{t+1})\le M_t(q_{t+1})\le M_t(q_t)=E(q_t)$, where the middle inequality holds because $q_{t+1}$ minimizes $M_t$ over the family. If $E(q_{t+1})=E(q_t)$, then $q_t$ also minimizes the strictly convex $M_t$, so the gradient of $M_t$ vanishes at $q_t$, and that gradient is $\nabla E(q_t)$ by Proposition~\ref{inference:g:prop:envelope}.
\end{proof}

This is the majorize-minimize principle \citep{inference:hunter2004mm}. With $\alpha$ recomputed at every iteration and damping $\rho=1$, the closed-form iterates descend the profiled energy monotonically, and their fixed points are its stationary points.

\paragraph{Damping.}
For $\rho<1$ the argument requires only $F(q_{t+1};\alpha_t)\le F(q_t;\alpha_t)$. In the full family this holds for every $\rho\in(0,1]$. With scalar $c$, Eq.~\eqref{inference:g:eq:quadratic_family} equals $c\KL(q\Vert q^\star)+\text{const}$ with $q^\star=\N(\Lambda^{-1}\eta,c\Lambda^{-1})$. In natural parameters $\theta$ with log-partition $A$, $\KL(q_\theta\Vert q^\star)=A(\theta^\star)-A(\theta)-\langle\theta^\star-\theta,\nabla A(\theta)\rangle$. Along $\theta_\rho=(1-\rho)\theta_0+\rho\theta^\star$, with $\delta=\theta^\star-\theta_0$,
\begin{equation}
\frac{d}{d\rho}\KL(q_{\theta_\rho}\Vert q^\star)=-(1-\rho)\delta^\top\nabla^2A(\theta_\rho)\delta\le0,
\end{equation}
so the damped step of Proposition~\ref{inference:g:prop:flow} descends the frozen energy, and the damped majorize-minimize step descends $E$. In the diagonal family the implemented interpolation $J_k=(1-\rho)/s_k+\rho\Lambda_{kk}/c$ contracts the mean error by the diagonal factor $w_k=(1-\rho)s_k^{-1}/J_k$, which differs between coordinates. With a dense $\Lambda$ a nonuniform diagonal contraction can increase the quadratic term. For $\Lambda=\bigl(\begin{smallmatrix}1&0.99\\0.99&1\end{smallmatrix}\bigr)$, $\eta=0$, $c=1$, $\xi=(5,-5)$, $s=(1/2,2)$ and $\rho=1/2$, the frozen energy without constants rises from $1.500$ to $2.503$, while $\rho=1$ lowers it to $1.000$. The damped diagonal step therefore carries no descent guarantee, for constant $\alpha$ as well as for the profiled energy.

\paragraph{Per-coordinate coefficients.}
The form \cfg{state\_dependent\_per\_coord} assigns $\alpha_k=c_{0,k}/(b_{0,k}+D_k)$ to each local coordinate. On the diagonal family with a diagonal prior, $q=\N(\mu,\diag s)$ and $p=\N(\mu_p,\diag p)$ in local coordinates, and the divergence separates as
\begin{equation}
D=\sum_kD_k,\qquad D_k=\frac12\Bigl[\frac{s_k}{p_k}+\frac{(\mu_k-\mu_{p,k})^2}{p_k}-1-\log\frac{s_k}{p_k}\Bigr].
\end{equation}
Lemma~\ref{inference:g:lem:profile} applies coordinatewise and gives $E=\sum_kc_{0,k}\log(b_{0,k}+D_k)+G+\text{const}$ with $\nabla E=\sum_k\alpha_k^\star\nabla D_k+\nabla G$. Freezing the coefficients gives
\begin{equation}
\sum_k\alpha_kD_k=\tfrac12(\mu-\mu_p)^\top\diag(\alpha/p)(\mu-\mu_p)+\tfrac12\sum_k\frac{\alpha_ks_k}{p_k}-\tfrac12\sum_k\alpha_k\log s_k+\text{const}.
\end{equation}
The prior precision becomes the nonscalar $\diag(\alpha_k/p_k)$, and the entropy weight becomes the vector $c_k=\alpha_k+c_G$. With $\Lambda_{\mathrm{loc}}=\Lambda_{G,\mathrm{loc}}+\diag(\alpha/p)$ and $\eta_{\mathrm{loc}}=\eta_{G,\mathrm{loc}}+\diag(\alpha/p)\mu_p$, the diagonal-family energy is
\begin{equation}
\tfrac12\mu^\top\Lambda_{\mathrm{loc}}\mu-\eta_{\mathrm{loc}}^\top\mu+\tfrac12\sum_k(\Lambda_{\mathrm{loc}})_{kk}s_k-\tfrac12\sum_kc_k\log s_k+\text{const},
\end{equation}
which separates in the variances. Its unique minimizer is $\mu^\star=\Lambda_{\mathrm{loc}}^{-1}\eta_{\mathrm{loc}}$ and $s_k^\star=c_k/(\Lambda_{\mathrm{loc}})_{kk}$. The vector weight therefore leaves the diagonal closed form exact, and Proposition~\ref{inference:g:prop:mm} holds with the sum of the coordinatewise tangent bounds as majorizer. The dense $\Lambda_{G,\mathrm{loc}}$ is the local pullback of the reference statistics, so only the self term is diagonal in local coordinates. In reference coordinates it contributes $U_i^\top\diag(\alpha/p)U_i$, which is not a scalar multiple of $C_p^{-1}$. The per-coordinate form is tied to the local basis and, unlike the per-token form, is not covariant under a change of local frame.

For the full family the code splits $D$ by the ordered chain rule, $D_k=\E_q\KL(q(x_k\mid x_{<k})\Vert p(x_k\mid x_{<k}))$. Write the prior precision as $C_p^{-1}=A\diag(1/v)A^\top$ with $A$ unit upper triangular and $v_k$ the conditional prior variances, and the covariance as $\Sigma=LL^\top$ with $L$ lower triangular. The cross-entropy of coordinate $k$ is $\tfrac1{2v_k}[a_k^\top\Sigma a_k+(a_k^\top(\mu-\mu_p))^2]$ up to a constant, with $a_k$ the $k$th column of $A$, and the conditional entropy is $\log L_{kk}$ up to a constant. The frozen energy is therefore
\begin{equation}
\tfrac12\Tr(\Lambda\Sigma)+\tfrac12\mu^\top\Lambda\mu-\eta^\top\mu-\sum_kc_k\log L_{kk}+\text{const},\qquad
\Lambda=\Lambda_G+A\diag(\alpha/v)A^\top,
\end{equation}
with $\eta=\eta_G+A\diag(\alpha/v)A^\top\mu_p$. It is strictly convex in $(\mu,L)$. Stationarity in $L$ requires the lower triangle of $\Lambda L$ to equal $\diag(c_k/L_{kk})$, so $\Lambda L$ is upper triangular. Writing $\Lambda=\tilde U\diag(d)\tilde U^\top$ with $\tilde U$ unit upper triangular, uniqueness of this factorization gives
\begin{equation}
\Sigma^{\star-1}=\tilde U\diag(d_k/c_k)\tilde U^\top,\qquad \mu^\star=\Lambda^{-1}\eta,
\end{equation}
which reduces to $\Sigma^\star=c\Lambda^{-1}$ when $c$ is uniform. A float64 check at $K=4$ against automatic differentiation of the code's ordered conditional divergence found gradients of $7\times10^{-16}$ in the mean and $2\times10^{-15}$ in the Cholesky factor at this point, while the lower-upper ordering of the factorization gave gradients of order one. This closed form is derived and checked but not implemented. The gradient route rejects \cfg{state\_dependent\_per\_coord} on the full family, and the descent argument for $\rho<1$ above does not apply, because the frozen energy is no longer a scalar multiple of a divergence.

The implementation clips diagonal closed-form variances to $[\cfg{eps},\cfg{sigma\_max}]$. At unit damping this is the exact box-constrained optimum because the variance terms are separable and strictly convex. The majorization proof applies on that fixed box for feasible incoming states; claiming the unconstrained optimum requires inactive bounds. Partial diagonal damping has no general descent guarantee. The full-covariance closed form does not apply this clipping. The implementation selects these forms with \cfg{lambda\_alpha\_mode} and the coefficients \cfg{b0} and \cfg{c0}. It uses the exact, unclipped divergence, since a ceiling on $D$ would zero the force beyond it. The diagnostics report the profiled energy of Eq.~\eqref{inference:g:eq:profiled}, with the constant omitted.



\subsection{Template routing and structured concept posteriors}
\label{inference:t}
A fixed exponential link still gives a linear fixed-weight message in source states. State-dependent inference over alternatives supplies another mechanism for nonlinearity. Template routing connects to feed-forward memory, persistent attention, and associative retrieval \citep{inference:geva2021,inference:sukhbaatar2019,inference:ramsauer2021}.
\subsubsection{Agents on two timescales}

Place a learned template bank in the reference space with the belief $\N(\xi_i,C_i)$. The bank is shared across sequences and changes during training. The executable local-coordinate block is distinguished from this covariant construction below.

\begin{definition}[Concept agents]
A bank of $M$ concept agents is a collection $\{(\pi_a,p_a,w_a)\}_{a=1}^M$, where $\pi_a>0$ with $\sum_a\pi_a=1$ is a prior weight, $p_a=\N(\nu_a,\Lambda_a)$ is a Gaussian template with $\Lambda_a\succ0$, and $w_a\in\mathbb{R}^K$ is a value. Concept agents are not tokens, are shared across all sequences, and hold beliefs that are parameters of the network, updated by training and fixed during a forward pass.
\end{definition}

The distinction between token agents and concept agents is one of timescale. Token beliefs are inferred on the fast timescale of a forward pass; concept beliefs are learned on the slow timescale of training, in the manner of a higher level of a hierarchical generative model whose states change slowly relative to the level below \citep{inference:friston2008hierarchical}. The attention step answers the question \emph{what does this context say about me}. Template routing answers a second question, \emph{which known kind of thing am I, and what does that imply}. The rest of this section makes the second question precise; Section~\ref{inference:t:sec:example} follows one sentence through the construction, and Section~\ref{inference:t:sec:social} examines how far the reading of templates as shared social norms holds.

\subsubsection{The mixture prior and its free energy}

Let agent $i$ place a mixture prior on its latent state $k\in\mathbb{R}^K$,
\begin{equation}
p(k)=\sum_{a=1}^M\pi_a p_a(k),
\label{inference:t:eq:mixture_prior}
\end{equation}
realized by a categorical latent $z\in\{1,\dots,M\}$ with $p(z=a)=\pi_a$ and $p(k\mid z=a)=p_a(k)$. The agent's belief over the joint latent $(k,z)$ is taken in the mean-field form $q_i(k)\rho_i(z)$, with $q_i$ given (it is the belief delivered by the preceding attention step) and $\rho_i$ a distribution on the $M$ labels. The variational free energy of the prior term is
\begin{equation}
\F_i[\rho_i]
=\E_{q_i\rho_i}\bigl[\log q_i(k)+\log\rho_i(z)-\log\pi_z-\log p_z(k)\bigr]
=\sum_a\rho_{ia}\Bigl[\KL(q_i\|p_a)+\log\frac{\rho_{ia}}{\pi_a}\Bigr].
\label{inference:t:eq:free_energy}
\end{equation}

\begin{proposition}[Responsibilities]
\label{inference:t:prop:responsibilities}
The minimizer of Eq.~\eqref{inference:t:eq:free_energy} over the probability simplex is
\begin{equation}
\rho_{ia}^{*}=\frac{\pi_a\exp\bigl(-\KL(q_i\|p_a)\bigr)}{\sum_b\pi_b\exp\bigl(-\KL(q_i\|p_b)\bigr)},
\label{inference:t:eq:responsibilities}
\end{equation}
and the minimum value is the soft minimum
\begin{equation}
\F_i^{*}=-\log\sum_a\pi_a\exp\bigl(-\KL(q_i\|p_a)\bigr).
\label{inference:t:eq:softmin}
\end{equation}
Moreover, $\KL\bigl(q_i\|\sum_a\pi_ap_a\bigr)\le\F_i[\rho_i]$ for every $\rho_i$, so $\F_i^{*}$ is the tightest bound of this form on the divergence from the mixture prior.
\end{proposition}

\begin{proof}
Write $c_a=\KL(q_i\|p_a)-\log\pi_a$. Then $\F_i[\rho_i]=\sum_a\rho_{ia}(c_a+\log\rho_{ia})$. With $Z=\sum_be^{-c_b}$ and $\rho^{*}_{ia}=e^{-c_a}/Z$, this equals $\KL(\rho_i\|\rho_i^{*})-\log Z$, which is minimized exactly at $\rho_i=\rho_i^{*}$ with value $-\log Z$; that is Eqs.~\eqref{inference:t:eq:responsibilities} and~\eqref{inference:t:eq:softmin}. For the bound, Jensen's inequality applied to the concave logarithm gives, pointwise in $k$,
\begin{equation*}
-\log\sum_a\pi_ap_a(k)=-\log\sum_a\rho_{ia}\frac{\pi_ap_a(k)}{\rho_{ia}}
\le-\sum_a\rho_{ia}\log\frac{\pi_ap_a(k)}{\rho_{ia}}.
\end{equation*}
Taking $\E_{q_i}$ and adding $\E_{q_i}[\log q_i]$ to both sides yields $\KL(q_i\|\sum_a\pi_ap_a)\le\F_i[\rho_i]$.
\end{proof}

The entropy of $q_i$ never appears in Eq.~\eqref{inference:t:eq:responsibilities}: in the cross-entropy $\E_{q_i}[-\log p_a]=\KL(q_i\|p_a)+H(q_i)$ it is the same for every template and cancels in the normalization. The routing rule is therefore the same Gibbs form as attention, with divergence to a concept in place of divergence to a peer and unit temperature fixed by the derivation rather than chosen. For a single-component belief, Eq.~\eqref{inference:t:eq:softmin} coincides with the variational approximation to the divergence between a Gaussian and a Gaussian mixture introduced by \citep{inference:hershey2007}. The bound is tight when the posterior over labels, $p(z\mid k)$, is constant over the support of $q_i$, and equality requires this constancy exactly. Dominance or concentration can make the bound approximately tight without proving equality at finite nonzero covariance.

\subsubsection{Explicit logits for diagonal templates}

The implementation uses diagonal templates, $\Lambda_a=\diag(\lambda_{a1},\dots,\lambda_{aK})$, which keeps the parameter count linear in $K$ per template. For a full belief covariance $C_i$, the Gaussian divergence is
\begin{equation}
\KL(q_i\|p_a)=\frac12\Bigl[\Tr(\Lambda_a^{-1}C_i)+(\xi_i-\nu_a)^{\top}\Lambda_a^{-1}(\xi_i-\nu_a)-K+\log\det\Lambda_a-\log\det C_i\Bigr].
\label{inference:t:eq:gaussian_kl}
\end{equation}
The terms $-K$ and $-\log\det C_i$ do not depend on $a$ and cancel in Eq.~\eqref{inference:t:eq:responsibilities}. With a diagonal $\Lambda_a$, the trace reads only the diagonal of $C_i$, so the logits are
\begin{equation}
\ell_{ia}=\log\pi_a-\frac12\sum_{k=1}^K\Bigl[\frac{(C_i)_{kk}+(\xi_{ik}-\nu_{ak})^2}{\lambda_{ak}}+\log\lambda_{ak}\Bigr],
\qquad
\rho^{*}_{ia}=\softmax_a(\ell_{ia}).
\label{inference:t:eq:logits}
\end{equation}
Two consequences follow directly. The logits are exact up to an $a$-independent constant, not an approximation to the divergence, and they require no determinant or inverse of the belief covariance. And a full covariance and its diagonal produce identical responsibilities against diagonal templates, so the construction is indifferent to the storage format of $C_i$. Expanding the square writes the quadratic part as three matrix products,
\begin{equation}
\sum_k\frac{(C_i)_{kk}+\xi_{ik}^2}{\lambda_{ak}}
-2\sum_k\frac{\xi_{ik}\nu_{ak}}{\lambda_{ak}}
+\sum_k\frac{\nu_{ak}^2}{\lambda_{ak}},
\label{inference:t:eq:matmul_form}
\end{equation}
each a contraction of a $K$-vector against an $M\times K$ table, which is how the logits are computed.

\subsubsection{Relation to the MLP}

Consider templates that share one covariance, $\Lambda_a=\Lambda$ for every $a$, with an arbitrary belief covariance $C_i$. Then $\Tr(\Lambda^{-1}C_i)$, $\log\det\Lambda$, and $\xi_i^{\top}\Lambda^{-1}\xi_i$ are all independent of $a$ and cancel in the softmax, leaving
\begin{equation}
\ell_{ia}=\nu_a^{\top}\Lambda^{-1}\xi_i-\frac12\nu_a^{\top}\Lambda^{-1}\nu_a+\log\pi_a+\text{const}_i.
\label{inference:t:eq:affine_logits}
\end{equation}
These are the pre-activations of a linear layer: the rows of the weight matrix are the whitened template means $\Lambda^{-1}\nu_a$, and the bias of each unit is fixed by its template's norm and prior weight. If the update is $\xi_i\mapsto\xi_i+\sum_a\rho^{*}_{ia}w_a$ (Section~\ref{inference:t:sec:update}), the values $w_a$ are the columns of the output matrix. Template routing in this regime is therefore a width-$M$ MLP whose pointwise activation has been replaced by a softmax across units. Table~\ref{inference:t:tab:correspondence} records the correspondence.

\begin{table}[h]
\centering
\begin{tabular}{@{}lll@{}}
\toprule
 & Coordinate MLP, the coordinate GELU MLP & Template routing \\
\midrule
Unit $a$ reads & row $a$ of $W_1$ & template $(\nu_a,\Lambda_a)$ \\
Similarity & $(W_1\xi_i+b_1)_a$ & $-\KL(q_i\|p_a)+\log\pi_a$ \\
Gate & $\operatorname{GELU}$, independent per unit & $\softmax$ across units \\
Unit $a$ writes & column $a$ of $W_2$ & value $w_a$ \\
Depends on $C_i$ & no & yes, through $\diag C_i$ \\
Parameters per layer & $2MK+M+K$ & $3MK+M$ \\
\bottomrule
\end{tabular}
\caption{Correspondence between the coordinate MLP and template routing at width $M$. The two coincide in form, up to the gate, when the templates share a covariance.}
\label{inference:t:tab:correspondence}
\end{table}

The gates differ in normalization and competition. The sigmoid-weighted linear unit satisfies the identity
\begin{equation}
x\sigma(x)=x\cdot\frac{e^{x}}{e^{x}+e^{0}},
\label{inference:t:eq:silu}
\end{equation}
which is softmax routing over two options with logits $(x,0)$ and values $(x,0)$ \citep{links:elfwing2018silu,inference:ramachandran2017}; GELU is close to $x\sigma(1.702x)$ \citep{links:hendrycks2016gelu}. The sigmoid approximation to GELU is thus a two-way competition between ``pass the pre-activation'' and ``pass nothing,'' decided by an exponential weight. Template routing replaces $M$ independent two-way competitions by a single $M$-way competition among concepts.

\begin{remark}[Independent gates]
\label{inference:t:rem:independent}
Independent per-unit gates also have a generative reading. For each concept separately, introduce a binary latent that selects between the template $p_a$ with probability $\pi_a$ and a shared background density $p_0$ otherwise. The same argument as Proposition~\ref{inference:t:prop:responsibilities}, applied to each two-component mixture separately, gives the gate
\begin{equation}
g_{ia}=\sigma\Bigl(\log\frac{\pi_a}{1-\pi_a}-\KL(q_i\|p_a)+\KL(q_i\|p_0)\Bigr),
\end{equation}
which is the two-option competition of Eq.~\eqref{inference:t:eq:silu} with log-odds given by a divergence difference, and the block output is $\sum_ag_{ia}w_a$. The choice between the softmax of Eq.~\eqref{inference:t:eq:responsibilities} and these gates is a modeling statement: one concept explains the agent, or several concepts may explain it at once. The implementation uses the softmax. The independent-gate variant is not implemented.
\end{remark}

\subsubsection{From responsibilities to a mean update}
\label{inference:t:sec:update}

The responsibilities say which concepts the agent recognizes itself as; a mean update must say what that recognition does to the belief. Two constructions are available, and they differ in what the values $w_a$ are allowed to be.

\paragraph{Tied values: descent on the concept free energy.}
Treat the belief mean as free and descend the optimal free energy of Eq.~\eqref{inference:t:eq:softmin}. By the envelope theorem, the gradient of the soft minimum is the responsibility-weighted gradient of the divergences,
\begin{equation}
\frac{\partial\F_i^{*}}{\partial\xi_i}=\sum_a\rho^{*}_{ia}\Lambda_a^{-1}(\xi_i-\nu_a).
\end{equation}
The Fisher metric of the Gaussian location family is $C_i^{-1}$, so the natural-gradient step with rate $\eta$ is
\begin{equation}
\xi_i\mapsto\xi_i-\eta C_i\sum_a\rho^{*}_{ia}\Lambda_a^{-1}(\xi_i-\nu_a).
\label{inference:t:eq:tied_update}
\end{equation}
Each agent is pulled toward the concepts it recognizes, weighted by how strongly it recognizes each one. With $\Lambda_a=\lambda I$ and a Euclidean step of size $\lambda$, the update is $\xi_i\mapsto\sum_a\rho^{*}_{ia}\nu_a$ with $\rho^{*}_{ia}=\softmax_a(\nu_a^{\top}\xi_i/\lambda-\|\nu_a\|^2/2\lambda+\log\pi_a)$, the retrieval step of a modern Hopfield network with inverse temperature $1/\lambda$ and stored patterns $\nu_a$ \citep{inference:ramsauer2021}; the norm term vanishes for patterns of equal norm. In this construction the value written by concept $a$ is tied to its key: recognition moves the belief toward the prototype. The operation cleans up or denoises a state, and it cannot write information that the prototype does not already contain.

\paragraph{Untied values: concepts that predict.}
An MLP does more than clean up a state, because its output columns are independent of its input rows: a recognized feature can add a direction unrelated to the feature itself. For template routing to do the same, a concept must predict something beyond its own location. Let each concept be a joint Gaussian over the agent's state $k$ and an increment $y\in\mathbb{R}^K$,
\begin{equation}
p_a(k,y)=\N\left(\begin{pmatrix}k\\ y\end{pmatrix};\begin{pmatrix}\nu_a\\ w_a\end{pmatrix},\begin{pmatrix}\Lambda_a&\Lambda_aB_a^{\top}\\ B_a\Lambda_a&\Gamma_a\end{pmatrix}\right),
\end{equation}
with a regression block $B_a$ and $\Gamma_a-B_a\Lambda_aB_a^{\top}\succ0$, so that the joint covariance is positive definite. For a point state $k$, conditioning the mixture $\sum_a\pi_ap_a(k,y)$ on $k$ gives the standard Gaussian mixture regression \citep{inference:ghahramani1994,inference:jacobs1991},
\begin{equation}
\E[y\mid k]=\sum_ar_a(k)\bigl(w_a+B_a(k-\nu_a)\bigr),
\qquad
r_a(k)=\frac{\pi_a\N(k;\nu_a,\Lambda_a)}{\sum_b\pi_b\N(k;\nu_b,\Lambda_b)}.
\label{inference:t:eq:gmr}
\end{equation}
Each concept is a local linear expert; the responsibilities decide which experts speak. Setting $B_a=0$ leaves the piecewise-constant readout $\sum_ar_a(k)w_a$, and the implemented block is
\begin{equation}
\xi_i\mapsto\xi_i+\sum_a\rho^{*}_{ia}w_a,
\qquad
C_i\mapsto C_i.
\label{inference:t:eq:implemented}
\end{equation}

Equation~\eqref{inference:t:eq:implemented} uses the variational responsibilities $\rho^{*}_{ia}$ of the belief, not the point responsibilities $r_a(k)$. The two agree in the limit of a concentrated belief: as $C_i\to0$, the $a$-dependent part of $\KL(q_i\|p_a)$ tends to $-\log\N(\xi_i;\nu_a,\Lambda_a)$, so $\rho^{*}_{ia}\to r_a(\xi_i)$. For a belief of finite width, $\sum_a\rho^{*}_{ia}w_a$ is not the exact expectation $\E_{q_i}\E[y\mid k]=\sum_a\E_{q_i}[r_a(k)]w_a$; it is the readout of the optimal mean-field label posterior of Proposition~\ref{inference:t:prop:responsibilities}. The implemented update is therefore a declared variational construction, exact at the point-belief limit, and not an exact conditional expectation for general $C_i$. Regression blocks add input-dependent component values and are not implemented. No model-class inclusion theorem relative to arbitrary GELU MLPs is asserted.

\subsubsection{What the covariance does}

Equation~\eqref{inference:t:eq:logits} contains the belief covariance through $\sum_k(C_i)_{kk}/\lambda_{ak}$, a term that has no counterpart in the coordinate GELU MLP. Its effect is to route uncertain agents toward broad concepts. For an isotropic template $\Lambda_a=\lambda I$, an isotropic belief $C_i=sI$, and a fixed squared distance $d^2=\|\xi_i-\nu_a\|^2$, the divergence is $\frac12[(Ks+d^2)/\lambda+K\log\lambda]$ up to $a$-independent terms, and it is minimized over the template width at
\begin{equation}
\lambda^{*}=s+\frac{d^2}{K}.
\end{equation}
An agent that is confident ($s$ small) prefers a template that is both near and sharp; an agent that is uncertain ($s$ large) pays a large trace penalty $Ks/2\lambda$ against any sharp template and prefers broad ones whose width matches its own. The same token mean is therefore routed differently depending on how certain the agent is, which is behavior a coordinate MLP cannot express because it never reads $C_i$. Template widths control mismatch sensitivity: a common scale $c$ divides quadratic terms by $c$. This is not generally categorical temperature, because relative log determinants and prior odds do not receive that scaling. Unit temperature belongs to the declared mixture-prior objective.

In the implemented block the covariance is read but not written: Eq.~\eqref{inference:t:eq:implemented} passes $C_i$ through unchanged, consistent with the evaluated architecture's treatment of covariance as a routing state. A covariance update from the concept step, for instance by moment matching of the mixture readout, is a possible extension.

\subsubsection{Gauge covariance and the implemented chart}
\begin{proposition}[Covariant template routing]
\label{inference:t:prop:covariant}
Place $p_a=\N(\nu_a,\Lambda_a)$ and $w_a$ in reference coordinates. Site $i$ sees $(U_i)_*p_a=\N(U_i\nu_a,U_i\Lambda_aU_i^\top)$. Responsibilities from $\KL(q_i\|(U_i)_*p_a)$ are invariant and $U_i\sum_a\rho_{ia}^*w_a$ transforms by $g_i$.
\end{proposition}
\begin{proof}
Common invertible pushforward preserves KL, giving $\KL(q_i\|(U_i)_*p_a)=\KL(\N(\xi_i,C_i)\|p_a)$. Reference data are invariant; the final frame supplies the message law.
\end{proof}
The implemented \cfg{gaussian\_template} block reads local $(\mu_i,\Sigma_i)$, adds local values, and ignores its supplied frame context. Its equations follow by replacing $(\xi_i,C_i)$ with $(\mu_i,\Sigma_i)$ and omitting the final frame lift. It is not an intertwiner under arbitrary independent site changes. The reference-bank variant is unimplemented. A reference-basis change must also transform the bank and need not preserve diagonal templates.


\subsubsection{Relation to attention}

Template routing is attention in which the keys are concept templates and the values are concept outputs. Setting the two messages side by side,
\begin{equation}
m_i^{\mathrm{peers}}=\sum_j\beta_{ij}L_{ij}\xi_j,
\qquad
m_i^{\mathrm{concepts}}=\sum_a\rho^{*}_{ia}w_a,
\end{equation}
the peer message transports what another agent believes, while the concept message writes what the corpus associates with a kind of state. Both weights are Gibbs distributions over divergences. The implementation keeps the two normalizations separate and applies the concept step after the attention step and normalization, in the position of the MLP. A single softmax over the union of peers and concepts is the alternative studied by \citep{inference:sukhbaatar2019}, who appended persistent vectors to every attention layer and removed the feed-forward sublayer; it corresponds to one mixture prior over both sources of evidence rather than two sequential inferences.

\subsubsection{A worked example: one sentence, token by token}
\label{inference:t:sec:example}

This section follows two tokens of the sentence ``a man went to the bank to deposit \dots'' through one template step. The numbers are a hand-built illustration in $K=2$ with $M=3$, chosen so that every quantity can be checked by hand; a trained layer has $K$ in the tens, hundreds of templates, and templates without human labels. Read the first coordinate as ``finance'' and the second as ``terrain and water.'' Table~\ref{inference:t:tab:toy_templates} lists the templates, all with uniform prior weight $\pi_a=1/3$.

\begin{table}[h]
\centering
\begin{tabular}{@{}llll@{}}
\toprule
Template & Mean $\nu_a$ & Widths $\lambda_a$ & What its value $w_a$ adds \\
\midrule
$a_1$ ``money institution'' & $(2,0)$ & $(0.5,0.5)$, sharp & toward \emph{money, check, cash, his} \\
$a_2$ ``riverbank'' & $(0,2)$ & $(0.5,0.5)$, sharp & toward \emph{river, water, sat, grass} \\
$a_3$ ``generic place noun'' & $(0.5,0.5)$ & $(3,3)$, broad & small generic push toward \emph{to, and, .} \\
\bottomrule
\end{tabular}
\caption{Toy template bank for the worked example. The labels are for the reader; the model sees only the parameters.}
\label{inference:t:tab:toy_templates}
\end{table}

\paragraph{The token ``bank.''}
Attention is causal, so at position six ``bank'' has seen only ``a man went to the.'' For the hand-built illustration, stipulate an unresolved state between the two senses: $\xi=(1,1)$ and $C=2I$. Equation~\eqref{inference:t:eq:logits} gives the logits and responsibilities in the first row of Table~\ref{inference:t:tab:toy_routing}. The sharp templates are each at squared distance $2$ from the belief, and the trace term $\sum_kC_{kk}/\lambda_{ak}=8$ penalizes them heavily; the broad template pays a trace term of only $4/3$. The uncertain token routes almost entirely to the generic template and receives a generic update, which illustrates mean-field coverage preference without establishing the preferable linguistic representation of ambiguity. The representation of ``bank'' at its own position never sees ``deposit''; it exists to predict the seventh token, and disambiguation happens at later positions.

\paragraph{The token ``deposit.''}
At position eight, attention lets ``deposit'' read its past, including ``bank.'' For the later token, stipulate a confident financial state: $\xi=(1.8,0.2)$ and $C=0.3I$. The second row of Table~\ref{inference:t:tab:toy_routing} shows the result. The token recognizes itself as an instance of $a_1$ with responsibility $0.82$, and the update $\xi\mapsto\xi+0.82w_1+0.18w_3$ moves its belief toward the region the decoder reads as \emph{money}, \emph{his}, or \emph{the check}, illustrating stipulated value meanings rather than a prediction from trained embeddings.

\paragraph{The same mean at lower confidence.}
The third row keeps $\xi=(1.8,0.2)$ but sets $C=2I$, as it would be after a muddled context. The routing reverses toward the hedging template, $0.21$ to $a_1$ and $0.79$ to $a_3$. A coordinate MLP reads only $\xi$ and would return the same output in both rows. Template routing distinguishes ``in financial territory and sure of it'' from ``in financial territory and guessing,'' which is the behavior derived in the preceding sections.

\begin{table}[h]
\centering
\begin{tabular}{@{}lcccccc@{}}
\toprule
Token state & $\ell_{1}$ & $\ell_{2}$ & $\ell_{3}$ & $\rho^{*}_{1}$ & $\rho^{*}_{2}$ & $\rho^{*}_{3}$ \\
\midrule
``bank,'' $\xi=(1,1)$, $C=2I$ & $-5.31$ & $-5.31$ & $-1.85$ & $0.030$ & $0.030$ & $0.941$ \\
``deposit,'' $\xi=(1.8,0.2)$, $C=0.3I$ & $0.01$ & $-6.39$ & $-1.50$ & $0.818$ & $0.001$ & $0.181$ \\
``deposit,'' $\xi=(1.8,0.2)$, $C=2I$ & $-3.39$ & $-9.79$ & $-2.06$ & $0.210$ & $0.000$ & $0.790$ \\
\bottomrule
\end{tabular}
\caption{Logits of Eq.~\eqref{inference:t:eq:logits} (dropping the common $\log\pi_a$) and responsibilities of Eq.~\eqref{inference:t:eq:responsibilities} for the toy bank of Table~\ref{inference:t:tab:toy_templates}, computed exactly and rounded.}
\label{inference:t:tab:toy_routing}
\end{table}

This two-dimensional example grounds the formulas with inspectable numbers. Its coordinates and human labels are stipulated. It does not show that training discovers these concepts, that both blocks are necessary for a prediction, or that every parameter moves monotonically toward frequent situations under supervised learning.

\subsubsection{A structured posterior: multimodal beliefs from the template step}
\label{inference:t:sec:structured}

The construction so far keeps one Gaussian per token and puts all multimodality in the categorical label: the posterior $q_i(k)\rho_i(z)$ is mean-field, so the belief about the state does not depend on which template is responsible. This section develops the structured alternative
\begin{equation}
q_i(k,z)=\rho_i(z)q_i(k\mid z),
\label{inference:t:eq:structured}
\end{equation}
in which each template the token considers carries its own Gaussian refinement of the state, and the token's belief becomes a genuine Gaussian mixture. It is an option recorded for later use; it is not implemented.

\paragraph{Motivation.}
Moment-matched attention already provides context-dependent covariance but still collapses to one Gaussian. A structured template posterior addresses the separate issue of retaining several template-conditioned readings within a block.

\paragraph{Evidence and prior.}
A structured posterior needs something to condition on. Treat the belief delivered by attention as a Gaussian message about the latent state, $\N(k;\xi_i,C_i)$, and the concept bank as the prior $\sum_a\pi_a\N(k;\nu_a,\Lambda_a)$. This is the reading used for conformity in Section~\ref{inference:t:sec:social}, and it is the natural reading when attention outputs are messages rather than posteriors.

\begin{proposition}[Exact structured posterior]
\label{inference:t:prop:structured}
With Gaussian evidence $\N(k;\xi_i,C_i)$ and the mixture prior of Eq.~\eqref{inference:t:eq:mixture_prior}, the exact posterior over $(k,z)$ has the form of Eq.~\eqref{inference:t:eq:structured} with
\begin{align}
\tilde\rho_{ia}&=\frac{\pi_a\N(\xi_i;\nu_a,\Lambda_a+C_i)}{\sum_b\pi_b\N(\xi_i;\nu_b,\Lambda_b+C_i)},
\label{inference:t:eq:bayes_resp}\\
q_i(k\mid z=a)&=\N(k;m_{ia},C_{ia}),\qquad
C_{ia}=\bigl(\Lambda_a^{-1}+C_i^{-1}\bigr)^{-1},\qquad
m_{ia}=\xi_i+C_i(C_i+\Lambda_a)^{-1}(\nu_a-\xi_i).
\label{inference:t:eq:component}
\end{align}
\end{proposition}

\begin{proof}
For each $a$, the product of two Gaussian densities in $k$ factorizes as $\N(k;\nu_a,\Lambda_a)\N(k;\xi_i,C_i)=\N(\xi_i;\nu_a,\Lambda_a+C_i)\N(k;m_{ia},C_{ia})$, with $C_{ia}$ and $m_{ia}$ as stated \citep{inference:bishop2006}; the second form of $m_{ia}$ follows from $C_{ia}C_i^{-1}=\Lambda_a(\Lambda_a+C_i)^{-1}=I-C_i(C_i+\Lambda_a)^{-1}$. The joint density of $(k,z=a)$ given the evidence is proportional to $\pi_a$ times this product. Normalizing over $k$ and $a$ gives Eq.~\eqref{inference:t:eq:bayes_resp} for the label marginal and Eq.~\eqref{inference:t:eq:component} for the conditional.
\end{proof}

Each component mean $m_{ia}$ is the conformity update of Eq.~\eqref{inference:t:eq:conformity}: the token's belief pulled toward template $a$ by the precision-weighted fraction $C_i(C_i+\Lambda_a)^{-1}$. The belief is therefore a mixture of ``what I would believe if template $a$ were the right reading,'' weighted by how well each reading predicts the evidence.

\paragraph{Two routing rules.}
Equation~\eqref{inference:t:eq:bayes_resp} is not the rule of Proposition~\ref{inference:t:prop:responsibilities}. Up to $a$-independent terms,
\begin{align*}
\log\tilde\rho_{ia}&=\log\pi_a-\tfrac12(\xi_i-\nu_a)^{\top}(\Lambda_a+C_i)^{-1}(\xi_i-\nu_a)-\tfrac12\log\det(\Lambda_a+C_i),\\
\log\rho^{*}_{ia}&=\log\pi_a-\tfrac12(\xi_i-\nu_a)^{\top}\Lambda_a^{-1}(\xi_i-\nu_a)-\tfrac12\Tr(\Lambda_a^{-1}C_i)-\tfrac12\log\det\Lambda_a.
\end{align*}
Both reduce to $\log\pi_a+\log\N(\xi_i;\nu_a,\Lambda_a)$ as $C_i\to0$, and they differ for uncertain beliefs. The mean-field rule treats $q_i$ as the posterior itself and asks whether template $a$ can cover the whole belief; the trace term penalizes sharp templates, so uncertain tokens route to broad ones. The Bayesian rule treats $q_i$ as noisy evidence and asks which template most plausibly generated the state; the evidence covariance is added to each template's, the distances shrink, and an uncertain token spreads its weight over the templates in proportion to their predictive density. Table~\ref{inference:t:tab:structured} repeats the worked example of Section~\ref{inference:t:sec:example} under Proposition~\ref{inference:t:prop:structured}. The uncertain ``bank'' now splits $0.37$ and $0.37$ between the financial and riverside readings with $0.26$ on the generic template, where the mean-field rule gave $0.94$ to the generic template, and its component means $m_{i1}=(1.8,0.2)$ and $m_{i2}=(0.2,1.8)$ are the two senses. These numbers illustrate alternative uncertainty semantics without deciding which is appropriate for a learned layer. The choice between the rules is a modeling statement about what attention delivers, a posterior or a message.

\paragraph{Collapsing to one Gaussian.}
Attention, transport, and the gauge results all assume Gaussian beliefs. The structured posterior can be used inside the block and collapsed at its output by moment matching,
\begin{equation}
\xi_i^{+}=\sum_a\tilde\rho_{ia}m_{ia},
\qquad
C_i^{+}=\underbrace{\sum_a\tilde\rho_{ia}C_{ia}}_{\text{within}}
+\underbrace{\sum_a\tilde\rho_{ia}(m_{ia}-\xi_i^{+})(m_{ia}-\xi_i^{+})^{\top}}_{\text{between}},
\label{inference:t:eq:collapse}
\end{equation}
which is the Gaussian closest to the mixture in $\KL(\text{mixture}\|\cdot)$. The between-template term is an ambiguity covariance: it is large exactly when the token splits its weight among templates whose refined means are far apart. It depends on context through $\xi_i$, so it gives the covariance a context-dependent update even on a route where nothing else writes it. In the toy example, $\Tr C^{+}$ has between-template part $0.96$ at ``bank'' and $0.02$ at confident ``deposit,'' a ratio of about $44$, while the incoming traces were $4$ and $0.6$.

\begin{table}[h]
\centering
\begin{tabular}{@{}lccccc@{}}
\toprule
Token state & $\tilde\rho_{1}$ & $\tilde\rho_{2}$ & $\tilde\rho_{3}$ & $\Tr(\text{within})$ & $\Tr(\text{between})$ \\
\midrule
``bank,'' $\xi=(1,1)$, $C=2I$ & $0.369$ & $0.369$ & $0.262$ & $1.22$ & $0.96$ \\
``deposit,'' $\xi=(1.8,0.2)$, $C=0.3I$ & $0.824$ & $0.015$ & $0.160$ & $0.40$ & $0.02$ \\
``deposit,'' $\xi=(1.8,0.2)$, $C=2I$ & $0.587$ & $0.163$ & $0.250$ & $1.20$ & $0.68$ \\
\bottomrule
\end{tabular}
\caption{The toy bank of Table~\ref{inference:t:tab:toy_templates} under the exact structured posterior of Proposition~\ref{inference:t:prop:structured}: Bayesian responsibilities of Eq.~\eqref{inference:t:eq:bayes_resp} and the two parts of the collapsed covariance of Eq.~\eqref{inference:t:eq:collapse}. Compare the mean-field responsibilities in Table~\ref{inference:t:tab:toy_routing}.}
\label{inference:t:tab:structured}
\end{table}

\paragraph{Values become exact.}
With joint templates $p_a(k,y)$ as in Eq.~\eqref{inference:t:eq:gmr} and $y\mid k,a\sim\N(w_a+B_a(k-\nu_a),\cdot)$, the structured posterior gives the readout
\begin{equation}
\E[y\mid\text{evidence}]=\sum_a\tilde\rho_{ia}\bigl(w_a+B_a(m_{ia}-\nu_a)\bigr),
\label{inference:t:eq:structured_readout}
\end{equation}
exactly, because the conditional mean is linear in $k$ within each component and the unobserved $y$ does not change the posterior over $k$. The mean-field readout of Section~\ref{inference:t:sec:update} was exact only as $C_i\to0$; the structured version removes that restriction, and with $B_a=0$ it is the implemented readout with $\rho^{*}$ replaced by $\tilde\rho$.

\paragraph{Gauge covariance and cost.}
In reference coordinates the evidence density and each $\N(\xi_i;\nu_a,\Lambda_a+C_i)$ transform with the same Jacobian factor, which is independent of $a$ and cancels in Eq.~\eqref{inference:t:eq:bayes_resp}, so the argument of Proposition~\ref{inference:t:prop:covariant} carries over to the structured posterior. The cost depends on the covariance used. With diagonal $\Lambda_a$ and $\diag C_i$, every quantity in Eqs.~\eqref{inference:t:eq:bayes_resp} and~\eqref{inference:t:eq:component} is elementwise, $O(MK)$ per token as in the implemented block, and the diagonal of Eq.~\eqref{inference:t:eq:collapse} is also $O(MK)$. A full between-template covariance costs $O(MK^2)$ per token through the low-rank form $D^{\top}\diag(\tilde\rho_i)D$ with $D$ the matrix of centered component means. With a full $C_i$, each template needs its own $(\Lambda_a+C_i)^{-1}$, which is $O(MK^3)$ per token and not practical at the widths used here.

\paragraph{Carrying mixtures through attention.}
The alternative to collapsing is to pass the mixture to the next layer. Attention would then compare mixtures, whose divergence has no closed form; the soft minimum of Proposition~\ref{inference:t:prop:responsibilities} and its generalization by \citep{inference:hershey2007} supply a variational surrogate, at $C^2$ component pairs per token pair for $C$-component beliefs, while routing would require a specified mixture cost. Total covariance remains exact for a declared mixture kernel; cost approximation and component truncation require separate error statements. The collapsed version keeps every existing result and adds multimodality only where the model already has a mixture prior. It retains the single-Gaussian interface.

\paragraph{Diagnostics.}
The structured posterior makes ambiguity measurable per token and per context: the entropy of $\tilde\rho_i$, the between-template trace of Eq.~\eqref{inference:t:eq:collapse}, and their correlation with next-token cross-entropy. These are the context-dependent counterparts of calibration of whichever covariance contract a saved run actually uses.

\subsubsection{Templates as a shared cultural repertoire}
\label{inference:t:sec:social}

The agent reading invites a social analogy: an individual reasoning within a conversation also consults a stock of shared expectations, norms, and scripts that no single conversation created. This section states what that analogy captures exactly, what it captures only in a modified variant, and where it fails. The short answer is that the templates behave like a shared cultural repertoire of norms and scripts that an individual consults, and not like a meta-agent in the sense of the multi-agent program.

\paragraph{What matches.}
Four features of the template bank have direct social counterparts. The bank is shared by every token in every sequence, as norms are shared by the members of a population. It changes on a slower timescale than the beliefs that consult it, as norms change across generations of cultural transmission rather than within a single exchange \citep{inference:boyd1985}. It is consulted and not renegotiated: a token cannot alter a template during a forward pass, just as an individual does not rewrite a convention in the course of following it. And the quantities in the free energy of Eq.~\eqref{inference:t:eq:free_energy} have natural readings. The divergence $\KL(q_i\|p_a)$ is the cost of an individual's belief deviating from norm $a$, the weight $\pi_a$ is the prevalence of the norm, and the width $\Lambda_a$ is its permissiveness: a broad template tolerates a wide range of beliefs, and a sharp one tolerates little.

\paragraph{Conformity and uncertainty.}
The tied-value construction of Section~\ref{inference:t:sec:update} is the analog of conformity. Suppose agent $i$ fuses its belief with norm $a$ by precision weighting, treating $p_a$ as a Gaussian prior and its current belief as the evidence. The resulting conjugate Gaussian has mean
\begin{equation}
\xi_i^{+}=\xi_i+C_i(C_i+\Lambda_a)^{-1}(\nu_a-\xi_i),
\label{inference:t:eq:conformity}
\end{equation}
of which the natural-gradient step Eq.~\eqref{inference:t:eq:tied_update} is the first-order form when $C_i\ll\Lambda_a$. For isotropic $C_i=sI$ and $\Lambda_a=\lambda I$, the agent moves the fraction $s/(s+\lambda)$ of the way to the norm. Uncertain agents conform more, confident agents conform less, and strict norms pull harder than permissive ones. This is the pattern of informational social influence, in which others' judgments are taken as evidence \citep{inference:deutsch1955}, and it is consistent with the classical observation that shared norms form readily when the stimulus is ambiguous and offers no external frame of reference \citep{inference:sherif1936}. Combined with the routing result of the worked example, the mean-field model makes two predictions: an uncertain agent defers to broad, generic norms, and a confident one consults sharp, specific norms but moves little toward them. Under the Bayesian routing of Section~\ref{inference:t:sec:structured} the first prediction changes: an uncertain agent hedges among the specific norms that could apply, in proportion to how well each explains its evidence, rather than retreating to a generic one. The implemented block uses untied values, so its update magnitude does not scale with $C_i$ in this way; only its routing does. Equation~\eqref{inference:t:eq:conformity} describes the tied variant.

\paragraph{Norms versus scripts.}
Untied values correspond to a different social object. A value $w_a$ that is independent of the template mean does not say ``be more like the norm''; it says ``in a situation of kind $a$, this is what tends to come next.'' That is a script in the sense of \citep{inference:schank1977}: a stored expectation about how a recognized kind of situation unfolds. A language model is trained to predict rather than to prescribe, so its trained templates should be read as scripts and schemas more than as norms of conduct. The two readings coexist in the joint-template construction of Eq.~\eqref{inference:t:eq:gmr}, where $\nu_a$ says what the situation looks like and $w_a$ says what follows it.

\paragraph{How norms would form: the population average.}
The analogy suggests asking how a template would be set if the agents themselves determined it, rather than an external optimizer. Sum the free energy of Eq.~\eqref{inference:t:eq:free_energy} over a population of beliefs $\{q_i\}$ with fixed responsibilities and minimize over the template. For a full-covariance template, the stationary point is
\begin{equation}
\pi_a=\frac1N\sum_i\rho^{*}_{ia},
\qquad
\nu_a=\frac{\sum_i\rho^{*}_{ia}\xi_i}{\sum_i\rho^{*}_{ia}},
\qquad
\Lambda_a=\frac{\sum_i\rho^{*}_{ia}\bigl[C_i+(\xi_i-\nu_a)(\xi_i-\nu_a)^{\top}\bigr]}{\sum_i\rho^{*}_{ia}},
\label{inference:t:eq:mstep}
\end{equation}
the M-step of expectation--maximization for a Gaussian mixture fitted to beliefs rather than points \citep{inference:dempster1977}. The diagonal case takes the diagonal of $\Lambda_a$. In social terms, a norm settles at the average belief of its adherents, weighted by how strongly each identifies with it; its prevalence is the fraction who adhere; and its permissiveness is the spread of their beliefs plus their own uncertainty, so a norm is strict exactly when its adherents are homogeneous and certain. This mathematical population-fitting rule has a social interpretation; it is not a theorem about human norm formation. Prevalence enters the routing only through $\log\pi_a$, so an agent's probability of adopting a norm, all else equal, is proportional to the norm's prevalence. That is unbiased transmission. Conformist transmission, in which common variants are adopted more than in proportion to their frequency \citep{inference:henrich1998}, would correspond to replacing $\log\pi_a$ by $\kappa\log\pi_a$ with $\kappa>1$; the derivation fixes $\kappa=1$. The implemented model does not use Eq.~\eqref{inference:t:eq:mstep}. It learns $(\nu_a,\Lambda_a,\pi_a,w_a)$ by gradient descent on next-token prediction, so trained templates are the norms that best serve prediction, which need not be the averages of their adherents.

\paragraph{Where the analogy fails.}
The templates are not meta-agents in the sense of the multi-agent program. There, a higher-level agent is formed from a particular coalition of constituent agents once they reach coherence, its state is built from theirs, and it acts back on them. A template is built from no particular tokens, exists before any sequence is read, and is not changed by the tokens that consult it during inference. A coherent token cluster could motivate a coalition construction, but ordinary attention alone does not prove meta-agent formation or its criteria. The distinction is between a culture's repertoire, which templates resemble, and a particular group, which a meta-agent is. A second failure is that the model has informational influence but no normative influence. \citep{inference:deutsch1955} separated conformity driven by evidence from conformity driven by the wish for approval, and the second requires sanctions that nothing in the model supplies; the only pressure toward a template is the prior cost. A third concerns the gate. Human beings hold several roles and norms at once, which fits the independent gates of Remark~\ref{inference:t:rem:independent} better than the softmax, which asserts that one template explains the agent. The softmax reading is closer to identifying the single situation one is in than to the set of norms one obeys.

\subsection{Current implementation of the inference routes}
\label{inference:implementation}
This source mapping describes the VFE\_3.0 working tree read on October 3, 2026, not the saved configuration of a historical training run. The executable attention modes are \cfg{projected\_predictive} and \cfg{projected\_gradient}; their comparison and content maps have full input and output width. Admitted families are \cfg{gaussian\_diagonal} and \cfg{gaussian\_full}, with at least two equal-width heads and the configured \cfg{block\_glk} or equal-irrep \cfg{so\_n} frames. These operators do not install independent group-valued edge variables.

The active launcher selects predictive attention, Gaussian-KL routing, \cfg{moment\_matched} covariance with floor $10^{-4}$, an additive attention residual, local-diagonal storage, a coordinate GELU block with covariance passthrough, and a linear decoder. It uses a causal source prior and $\kappa_\beta=1$. Receiver-law controls in that same dictionary are inactive on the predictive route. Its state-dependent self-coupling selectors have no consumer while the optional M-step self-coupling weight is zero. These are dated active settings, not evidence for predictive performance.

\paragraph{Predictive execution.}
\cfg{vfe\_block} dispatches to the predictive branch of \cfg{e\_step\_iteration}, then \cfg{relational\_predictive\_readout}. This computes \cfg{relational\_score}, applies \cfg{safe\_kl\_clamp}, obtains the row with \cfg{attention\_weights}, and calls \cfg{predictive\_content}. Moment matching uses those same weights and the value frame in \cfg{predictive\_covariance\_update}, including shared-source cross-head and residual self-edge terms. The active diagonal route stores the local diagonal of those kernel moments. The linear decoder reads only the final mean; the evolved covariance can still affect subsequent routing.

\paragraph{Receiver execution.}
\cfg{build\_full\_link\_context} snapshots sources and routing with their outer graphs attached. \cfg{\_source\_channel} builds either $M_h\bar C_jM_h^\top+\nu I_K$ or $T_h\bar C_jT_h^\top+\nu I_d$. The channel law pools $d$-dimensional statistics and lifts with $O_h$, implementing $P_h=O_h^\top$. Both \cfg{ambient\_alignment} and \cfg{channel\_likelihood} are implemented, as are \cfg{euler} and \cfg{closed\_form}. New defaults are channel likelihood and closed form; serialized configurations predating these selectors retain historical ambient/Euler semantics. Code requires strictly positive receiver noise: zero noise with full-row-rank $T_h$ is a theoretical case, not a current toggle.

\cfg{full\_link\_context\_at} recomputes profiled $\alpha^\star$ at each current receiver using unclamped self KL. \cfg{full\_link\_gradients} supplies the quadratic-family force and \cfg{full\_link\_minimizer} the full-family interpolation or diagonal update. Full-family per-coordinate coupling is rejected. Closed form reads \cfg{relational\_receiver\_damping}; Euler rates and trust controls are inactive there. If covariance updates are disabled, the covariance stays fixed and the mean interpolates directly toward its minimizer. Constant coupling and unit damping need only one exact solve; profiled coupling can require further iterations even at unit damping because the majorizer is recomputed. Neither joint content-informed routing nor a structured categorical cut posterior is implemented.

\paragraph{Template execution.}
\cfg{GaussianTemplateBlockMLP} is selected by \cfg{use\_block\_mlp=True} and \cfg{block\_mlp\_mode="gaussian\_template"}. Each layer owns $M=K\cdot\cfg{block\_mlp\_expansion}$ concepts with \cfg{nu}, \cfg{log\_lambda}, \cfg{log\_pi}, and \cfg{values}, totaling $3MK+M$ parameters. Widths start at one and priors uniform; means and values use the corresponding coordinate-MLP initialization scales. Logits read the local covariance diagonal and are evaluated outside autocast at at least float32 before being returned at the mean dtype. Configured dropout acts on the value sum before its residual addition. The deterministic readout formula therefore applies in evaluation or with zero dropout. Covariance passthrough is the only admitted template contract; activation and covariance-floor controls are inactive. Leading work is approximately $3MK$ multiply-adds per token versus $2MK$ for the width-$M$ two-layer MLP, excluding lower-order operations. Equal expansion does not match parameter count.

Reference-coordinate templates, regression blocks, independent gates, structured posterior mixtures, template covariance updates, and the population-fitted M-step remain mathematical extensions. The structured posterior is exact under its separately declared Gaussian-message likelihood. Treating an existing attention posterior as new likelihood evidence can count earlier evidence twice unless a larger model specifies the reuse. Source tracing and finite checks do not establish improved loss, calibrated uncertainty, or behavior at the target scale; those require matched training studies.

\section{Statistical geometry of stored frames}\label{frame:sec:overview}
The Gaussian orbit geometry also supplies a metric for the learned frame factors. A group-valued link and a frame act on a Gaussian through the same linear pushforward, but they are different parameters: the implemented optimizer described here updates the stored token and positional frames, whereas applying an analogous metric to independently parameterized exponential links requires specifying their statistical family and optimizer. No Fisher update of the content maps is inferred from the frame implementation.

Two metrics on stored frame coordinates must be distinguished. The ambient Frobenius metric pulled back through the exponential chart gives projected Euclidean descent on the frame matrix. The Fisher information of a reference Gaussian prior gives a positive semidefinite metric whose kernel consists of transformations that preserve that Gaussian; a specified completion makes the implemented row solves nonsingular. Both constructions yield a tangent direction, and both require a finite-step retraction. Their exact identities, branch assumptions, symmetry limits, and computational realizations are developed below.

These results extend the distinction between the extrinsic exponential pullback and the left-invariant group metric in the \emph{GL(K) gradients} companion. That companion's Gaussian natural gradient concerns belief moments; the frame Fisher below instead pulls that metric back through a separately declared family indexed by a frame. The right-trivialized differential called \(\mathrm{dexp}\) in the companion is \(\Psi_R(\operatorname{ad}_X)\) here. We use reference means \(\xi_i\), reference covariances \(C_i\), and frames \(U_i\) consistently with the main attention manuscript; the step generator is denoted \(\mathsf Y\).

The mathematical statements apply under their displayed assumptions. Numerical examples retained from the September 28 derivation are historical CPU calculations at small block sizes, used to check formulas and finite-precision behavior. They are not training comparisons or measurements of current GPU performance. The current implementation distinctions are stated explicitly in Section~\ref{frame:sec:implementation}.

\subsection{Setting and notation}\label{frame:sec:setting}

\subsubsection{Frames, algebras and stored tables}

Let $\g\subset\gl(K)$ be a real Lie subalgebra of dimension $n$ with basis $T_1,\dots,T_n$ and Gram matrix $B_{ab}=\ip{T_a}{T_b}$, where $\ip{A}{C}=\tr(A^\top C)$ is the Frobenius pairing. The coordinate map is $\iota:\R^n\to\g$, $\iota(\varphi)=\sum_a\varphi_aT_a$. A stored chart row $\varphi$ defines $X=\iota(\varphi)$ and the frame $U=e^X$. Let $G\subset\GL(K)$ be the connected subgroup with Lie algebra $\g$ and $\Pi_\g$ the Frobenius-orthogonal projector onto $\g$.

For the untied \texttt{block\_glk} group, $K=Hd$ with $H$ heads of dimension $d$, $\g=\bigoplus_{h=1}^{H}\gl(d)$, the basis is the row-major elementary basis, $B=I$, and $X$ is the row $\varphi$ reshaped to $(H,d,d)$. The launcher inspected on October 3, 2026 uses $K=20$, $H=2$, and $d=10$; the results are stated for general dimensions. The \texttt{glk} group is the case $H=1$, $d=K$. For these groups $\iota$ is a Frobenius isometry, and we identify a row with its matrix whenever no confusion arises.

The token-type table stores, for each type $v$, a chart row $\varphi_v$ with $E_v=e^{X_v}$, $X_v=\iota(\varphi_v)$, together with the local prior $q_v=\Normal(\mu_v,S_v)$, $S_v=\operatorname{diag}(s_v)$. The learned positional table stores $\theta_p$ with $F_p=e^{\Theta_p}$, $\Theta_p=\iota(\theta_p)$. The realized frame of occurrence $i$, a token of type $v_i$ at position $p_i$, is the exact group product
\begin{equation}\label{frame:eq:realized}
U_i=E_{v_i}F_{p_i},\qquad U_i^{-1}=F_{p_i}^{-1}E_{v_i}^{-1}.
\end{equation}
A fixed right positional rotation can be appended as \(U_i=E_{v_i}F_{p_i}R_{p_i}^{\top}\). The current projected routes use this right insertion when RoPE is enabled. Because \(R_{p_i}\) is fixed during a frame update, the type-metric position-independence argument below is unchanged; it is a parameter-independent bijection of the reference sample space. When \texttt{pos\_phi="none"}, \(F_p=I\) and there is no learned positional row to optimize.

For $A\in\GL(K)$ write $A\cdot\Normal(m,\Sigma)=\Normal(Am,A\Sigma A^\top)$ for the linear pushforward. At layer $0$ the belief of occurrence $i$ is the prior, $q_i=q_{v_i}$. The attention read path forms the reference belief
\begin{equation}\label{frame:eq:reference}
r_i=U_i^{-1}\cdot q_i=\Normal(\xi_i,C_i),\qquad \xi_i=U_i^{-1}\mu_i,\qquad C_i=U_i^{-1}\Sigma_iU_i^{-\top},
\end{equation}
with full covariance blocks $C_i$ even when $\Sigma_i$ is diagonal.

\subsubsection{The three group actions}

Three actions organize the symmetry statements. The \emph{left gauge action} at occurrence $i$ is $(\mu_i,\Sigma_i,U_i)\mapsto(g_i\mu_i,g_i\Sigma_ig_i^\top,g_iU_i)$; it fixes every reference belief $r_i$. The \emph{right reference change} is $U_i\mapsto U_ih$ for one $h$ common to all occurrences; it sends $r_i\mapsto h^{-1}\cdot r_i$, and the learned maps downstream of the reference co-transform. The \emph{split redundancy} of the stored tables is $E_v\mapsto E_vh$ for every type and $F_p\mapsto h^{-1}F_p$ for every position; by \eqref{frame:eq:realized} it leaves every $U_i$, and therefore the loss, unchanged. Since the positional factor sits on the right of $E_v$, only per-type left actions $E_v\mapsto g_vE_v$ are symmetries of the table parametrization itself.

The full general-linear transformation statements below concern unrestricted invertible frames and full Gaussian moments. For a block or other restricted subgroup, the admissible changes of reference basis must preserve that subgroup and the declared family. A transformation can preserve the underlying frame matrices while leaving a single exponential chart, so chart-domain and finite-step claims require their additional hypotheses.

\subsubsection{The exponential chart}

Write $\ad_XY=[X,Y]$ and define the entire functions $\Psi_L$, $\Psi_R$ and their meromorphic reciprocal $h$ by
\begin{equation}\label{frame:eq:psi}
\Psi_L(z)=\frac{1-e^{-z}}{z},\qquad \Psi_R(z)=\frac{e^{z}-1}{z},\qquad h(z)=\frac{1}{\Psi_L(z)}=\frac{z}{1-e^{-z}}.
\end{equation}
The principal strip is $\mathcal{S}=\{z\in\C:|\Imag z|<\pi\}$.

\begin{lemma}[derivative of the exponential]\label{frame:lem:dexp}
For $X,Y\in\gl(K)$, the standard exponential derivative \cite[Ch.~5]{frame:Hall2015} is
\begin{equation}\label{frame:eq:dexp}
L_{\exp}(X,Y):=\frac{d}{dt}e^{X+tY}\Big|_{t=0}=e^X\Psi_L(\ad_X)(Y)=\Psi_R(\ad_X)(Y)e^X.
\end{equation}
\end{lemma}

Consequently, a chart velocity $\delta X$ produces the body generator $Y=U^{-1}\delta U=\Psi_L(\ad_X)\delta X$ and the spatial generator $W=(\delta U)U^{-1}=\Psi_R(\ad_X)\delta X=\Ad_UY$. The implementation's \texttt{xi} in \texttt{PullbackGroupDirectionResult} is a body generator, because the staging compares $\exp(\varphi-\eta v)$ with $\exp(\varphi)\exp(-\eta\mathsf{Y})$.

\begin{lemma}[adjoints]\label{frame:lem:adjoint}
With respect to the Frobenius pairing, $\ad_X^*=\ad_{X^\top}$, and $f(\ad_X)^*=f(\ad_{X^\top})$ for every power series $f$ with real coefficients. In particular
\begin{equation}\label{frame:eq:dexp-adjoint}
L_{\exp}(X,\cdot)^*(M)=\Psi_L(\ad_{X^\top})(U^\top M)=L_{\exp}(X^\top,M).
\end{equation}
\end{lemma}
\begin{proof}
$\ip{A}{[X,Y]}=\tr(A^\top XY)-\tr(A^\top YX)=\ip{X^\top A}{Y}-\ip{AX^\top}{Y}=\ip{[X^\top,A]}{Y}$. Adjoints of real power series in $\ad_X$ follow term by term. For \eqref{frame:eq:dexp-adjoint}, $\ip{M}{U\Psi_L(\ad_X)Y}=\ip{\Psi_L(\ad_{X^\top})(U^\top M)}{Y}$. Left multiplication by $U^\top=e^{X^\top}$ commutes with $\ad_{X^\top}$ because $e^{X^\top}$ commutes with $X^\top$, so $\Psi_L(\ad_{X^\top})(U^\top M)=U^\top\Psi_L(\ad_{X^\top})M=L_{\exp}(X^\top,M)$.
\end{proof}

For a loss $\ell$ and a real-coefficient chart with $B=I$, Lemma~\ref{frame:lem:adjoint} converts covectors: with $d\ell=\ip{G_Y}{Y}=\ip{G_W}{W}$,
\begin{equation}\label{frame:eq:covectors}
\nabla_\varphi\ell=\Psi_L(\ad_{X^\top})(G_Y)=\Psi_R(\ad_{X^\top})(G_W),\qquad G_W=U^{-\top}G_YU^{\top}.
\end{equation}

\begin{lemma}[regularity of the chart]\label{frame:lem:regular}
$L_{\exp}(X,\cdot)$ restricted to $\g$ is injective if and only if no eigenvalue of $\ad_X|_\g$ lies in $2\pi i(\mathbb{Z}\setminus\{0\})$. This holds whenever $\spec(X)\subset\mathcal{S}$.
\end{lemma}
\begin{proof}
$U$ is invertible and $\ad_X$ preserves $\g$, so injectivity is invertibility of $\Psi_L(\ad_X|_\g)$, whose eigenvalues are $\Psi_L(\lambda)$ for $\lambda\in\spec(\ad_X|_\g)$. $\Psi_L$ vanishes exactly at $2\pi ik$, $k\neq0$. The eigenvalues of $\ad_X$ on $\gl(K)$ are the differences $\lambda_i-\lambda_j$ of eigenvalues of $X$, and those on $\g$ form a subset. If $\spec(X)\subset\mathcal{S}$ then $|\mathrm{Im}(\lambda_i-\lambda_j)|<2\pi$.
\end{proof}

The chart therefore has full rank everywhere inside the principal strip. On $\gl(K)$ it loses rank whenever two eigenvalues of $X$ differ by $2\pi ik$, $k\neq0$. A path that leaves the strip continuously loses rank at the moment it leaves, when a conjugate pair of eigenvalues reaches $a\pm i\pi$; off the strip rank is also lost for eigenvalues that are not a conjugate pair, for example $\pm i\pi/2$ together with $\pm5i\pi/2$.

\begin{assumption}\label{frame:ass:main}
The algebra \(\g\) is closed under transposition. This holds for real general linear, block general linear, tied block general linear, orthogonal, special linear, and symplectic algebras and their direct sums. For the interpretation as descent on the frame matrix, the loss has the form \(\ell(\varphi)=L(e^{\iota(\varphi)})\), with \(L\) twice continuously differentiable near the group. A coordinate penalty or an active clamp on the exponential argument must be treated separately. The exact projection identities below use the undamped pullback matrix; a ridge changes the direction. Floating-point gradient error is an additional numerical input, not an assumption with a universal error magnitude.
\end{assumption}

\subsection{The Frobenius-pullback step is projected Euclidean descent on \texorpdfstring{$U$}{U}}\label{frame:sec:pullback}

Let $D=L_{\exp}(X,\cdot)\circ\iota:\R^n\to\gl(K)$ be the differential of the chart, and $G_0=D^\top D$ the pullback of the Frobenius metric. The undamped \texttt{pullback\_group} direction is $v=G_0^{-1}\nabla_\varphi\ell$.

\begin{theorem}\label{frame:thm:projected}
Assume $D$ is injective. Let $\Gamma=\nabla_UL$ be the ambient gradient, $\mathcal{T}=\range(D)$ and $P_\mathcal{T}$ the Frobenius-orthogonal projector onto $\mathcal{T}$. Then
\begin{equation}\label{frame:eq:projected}
v=D^+\Gamma=D^+P_\mathcal{T}\Gamma,\qquad Dv=P_\mathcal{T}\Gamma,\qquad e^{\iota(\varphi-\eta v)}=U-\eta P_\mathcal{T}\Gamma+O(\eta^2),\qquad \mathcal{T}=U\g=\g U.
\end{equation}
\end{theorem}
\begin{proof}
By the chain rule $\nabla_\varphi\ell=D^\top\Gamma$. Since $G_0=D^\top D$ is invertible, $v=(D^\top D)^{-1}D^\top\Gamma=D^+\Gamma$. The operator $DD^+$ is the orthogonal projector onto $\range(D)$, and $D^+$ annihilates $\range(D)^\perp=\ker D^\top$, so $D^+=D^+P_\mathcal{T}$ and $Dv=DD^+\Gamma=P_\mathcal{T}\Gamma$. A first-order Taylor expansion of $e^{X-\eta\iota(v)}$ gives $U-\eta Dv+O(\eta^2)$. Finally $\mathcal{T}=U\Psi_L(\ad_X)(\g)=U\g$, because $\Psi_L(\ad_X)$ is a bijection of $\g$ by Lemma~\ref{frame:lem:regular}, and $U\g U^{-1}=e^{\ad_X}\g=\g$ gives $U\g=\g U$.
\end{proof}

Geometrically, $F:\varphi\mapsto e^{\iota(\varphi)}$ is a local diffeomorphism onto $G$ and $G_0$ is by construction $F^*$ of the Frobenius metric restricted to $G$. So $F$ is a local isometry, and it carries the $G_0$-gradient of $\ell$ to the gradient of $L|_G$ for the induced metric, which is the orthogonal projection of the ambient gradient onto the tangent space \cite[Sec.~3.6.1]{frame:Absil2008}. For $\g=\gl(K)$ the projector is the identity and the step is literally Euclidean gradient descent on $U$; for block $\gl$ it keeps the diagonal blocks.

\begin{proposition}[left-trivialized normal equations]\label{frame:prop:normal}
$P_\mathcal{T}\Gamma=U\mathsf{Y}$, where $\mathsf{Y}\in\g$ is the unique solution of
\begin{equation}\label{frame:eq:normal}
A_U\mathsf{Y}=c,\qquad A_U:=\Pi_\g(U^\top U(\cdot))|_\g,\qquad c:=\Pi_\g(U^\top\Gamma).
\end{equation}
$A_U$ is symmetric positive definite on $\g$, since $\ip{Y}{A_UY}=\|UY\|_F^2$.
\end{proposition}
\begin{proof}
By Theorem~\ref{frame:thm:projected}, $P_\mathcal{T}\Gamma=U\mathsf{Y}$ where $\mathsf{Y}$ minimizes $\|UY-\Gamma\|_F^2$ over $Y\in\g$. Stationarity is $\ip{UZ}{UY-\Gamma}=0$ for all $Z\in\g$, which is $\Pi_\g(U^\top UY)=\Pi_\g(U^\top\Gamma)$. Positive definiteness follows from invertibility of $U$.
\end{proof}

\begin{remark}[the pullback metric is not a group translate]\label{frame:rem:translate}
For the ambient Frobenius metric, the body-coordinate quadratic form is \(\langle Y_1,U^\top UY_2\rangle_F\), and the spatial-coordinate form is \(\langle W_1,W_2UU^\top\rangle_F\). These are not invariant under general left or right multiplication. On an orthogonal subgroup both reduce to the same bi-invariant Frobenius form. The naive translated direction \(\Pi_\g(U^{-1}\Gamma)\) solves the normal equations when left multiplication by \(U^\top U\) preserves \(\g\); without that condition it need not equal the projected-descent direction. Near the identity its discrepancy is generally first order in \(X\).
\end{remark}

\begin{proposition}[closed forms of $A_U^{-1}$]\label{frame:prop:closed}
Let $\mathcal{S}_U=U^\top U$.
\begin{enumerate}
\item[(i)] If $\mathcal{S}_U\g\subset\g$ (true for $\gl$, block $\gl$, tied block $\gl$), then $\mathsf{Y}=\mathcal{S}_U^{-1}c$. For $\gl$ this is $\mathsf{Y}=U^{-1}\Gamma$; for block $\gl$ it is $\mathsf{Y}=U^{-1}\Pi_\g(\Gamma)$, computed blockwise.
\item[(ii)] If $U^\top U=I$ (any $\g\subset\mathfrak{o}(K)$), then $\mathsf{Y}=c$; for $\mathfrak{so}(K)$, $\mathsf{Y}=\operatorname{skew}(U^\top\Gamma)$.
\item[(iii)] For $\mathfrak{sl}(K)$, $\mathsf{Y}=\mathcal{S}_U^{-1}(c+\kappa I)$ with $\kappa=-\tr(\mathcal{S}_U^{-1}c)/\tr(\mathcal{S}_U^{-1})$.
\item[(iv)] For $\mathfrak{sp}(2m)$ with $J=\begin{pmatrix}0&I\\-I&0\end{pmatrix}$, write $\mathcal{S}_U=V\operatorname{diag}(\sigma)V^\top$ and $\hat C=Jc$. Then $\mathsf{Y}=J^\top V\tilde\Sigma V^\top$ with $\tilde\Sigma_{jk}=2(V^\top\hat CV)_{jk}\sigma_j\sigma_k/(\sigma_j+\sigma_k)$.
\item[(v)] For a general transpose-closed $\g$, solve the $n\times n$ system $Ny=\rho$ with $N_{ab}=\ip{UT_a}{UT_b}$ and $\rho_a=\ip{T_a}{c}$, or run conjugate gradients with the matrix-vector product $Y\mapsto\Pi_\g(\mathcal{S}_UY)$ at cost $O(K^3)$ per product.
\end{enumerate}
\end{proposition}
\begin{proof}
(i) The map $Y\mapsto\mathcal{S}_UY$ is injective on $\g$ and maps $\g$ into $\g$, hence bijective, so $\mathcal{S}_U^{-1}\g=\g$ and $\Pi_\g(\mathcal{S}_U\mathsf{Y})=\mathcal{S}_U\mathsf{Y}$. For block-diagonal $U$, $\Pi_\g(U^\top\Gamma)=U^\top\Pi_\g(\Gamma)$. (ii) is immediate. (iii) $\Pi_{\mathfrak{sl}}(M)=M-(\tr M/K)I$, so the equation reads $\mathcal{S}_U\mathsf{Y}=c+\kappa I$ for some scalar $\kappa$, and $\tr\mathsf{Y}=0$ determines $\kappa$ as stated. (iv) $Y\in\mathfrak{sp}(2m)$ if and only if $JY$ is symmetric, and $\Sigma\mapsto J^\top\Sigma$ is a Frobenius isometry from symmetric matrices onto $\mathfrak{sp}(2m)$, so $\Pi_{\mathfrak{sp}}(M)=J^\top\sym(JM)$. $\mathcal{S}_U$ is symmetric positive definite and symplectic, so $\mathcal{S}_UJ\mathcal{S}_U=J$ and $J\mathcal{S}_UJ^\top=\mathcal{S}_U^{-1}$. Hence $A_U(J^\top\Sigma)=J^\top\sym(J\mathcal{S}_UJ^\top\Sigma)=J^\top(\mathcal{S}_U^{-1}\Sigma+\Sigma\mathcal{S}_U^{-1})/2$, and the equation $A_U(J^\top\Sigma)=c=J^\top\hat C$ is the Lyapunov equation $(\mathcal{S}_U^{-1}\Sigma+\Sigma\mathcal{S}_U^{-1})/2=\hat C$, which is diagonal in the eigenbasis of $\mathcal{S}_U$ and gives $\tilde\Sigma$. (v) is \eqref{frame:eq:normal} in coordinates.
\end{proof}

\begin{proposition}[recovering the ambient gradient from the chart gradient]\label{frame:prop:recovery}
Let $m=\iota(B^{-1}\nabla_\varphi\ell)\in\g$; for \texttt{block\_glk} this is the gradient row reshaped. Under the transpose-closure assumption, $m=\Psi_L(\ad_{X^\top})(c)$. If $\spec(X)\subset\mathcal{S}$, then
\begin{equation}\label{frame:eq:recovery}
c=\Psi_L(\ad_{X^\top})^{-1}(m)=U^\top L_{\log}(U^\top,m),\qquad L_{\log}(U^\top,m)=\left[\Log\begin{pmatrix}U^\top&m\\0&U^\top\end{pmatrix}\right]_{12},
\end{equation}
where $L_{\log}$ is the Fr\'echet derivative of the principal logarithm. For $\g=\gl(K)$ this gives $\Gamma=L_{\log}(U^\top,m)$ exactly; for block $\gl$ it gives $\Pi_\g(\Gamma)=L_{\log}(U^\top,m)$ exactly, which is all that Proposition~\ref{frame:prop:closed}(i) uses.
\end{proposition}
\begin{proof}
First, $(\nabla_\varphi\ell)_a=\ip{U\Psi_L(\ad_X)T_a}{\Gamma}=\ip{T_a}{\Psi_L(\ad_{X^\top})(U^\top\Gamma)}$ by Lemma~\ref{frame:lem:adjoint}. By transpose closure, $\ad_{X^\top}$ preserves $\g$, and it also preserves $\g^\perp$: for $Y\perp\g$ and $Z\in\g$, $\ip{\ad_{X^\top}Y}{Z}=\ip{Y}{\ad_XZ}=0$. Hence $\ad_{X^\top}$ commutes with $\Pi_\g$, and $m=\Pi_\g\Psi_L(\ad_{X^\top})(U^\top\Gamma)=\Psi_L(\ad_{X^\top})(c)$.

Next, on the open set of matrices $A$ with $\spec(A)\subset\mathcal{S}$ we have $\Log\circ\exp=\mathrm{id}$ \cite[Thm.~1.31]{frame:Higham2008}. Differentiating and using Lemma~\ref{frame:lem:dexp}, $L_{\log}(e^A,e^A\Psi_L(\ad_A)H)=H$, so $\Psi_L(\ad_A)^{-1}(M)=L_{\log}(e^A,e^AM)$. For a primary matrix function $f$ and diagonalizable $B$, the Daleckii--Krein representation $L_f(B,E)_{jk}=f[\beta_j,\beta_k]E_{jk}$ in the eigenbasis of $B$ shows $L_f(B,BE)=BL_f(B,E)$, and the identity extends to all $B$ in the domain by continuity. Taking $A=X^\top$, $e^A=U^\top$, gives $c=U^\top L_{\log}(U^\top,m)$. The block formula is the block-triangular representation of the Fr\'echet derivative \cite{frame:Mathias1996}, \cite[Thm.~3.6]{frame:Higham2008}. For $\gl$, $c=U^\top\Gamma$; for block $\gl$ with block-diagonal $U$, $c=U^\top\Pi_\g(\Gamma)$.
\end{proof}

The strip condition fixes the logarithm branch required for recovery. Outside it, the principal logarithm, when defined, need not equal \(X^\top\), and its Fr\'echet derivative inverts the exponential at a different chart representative. A diagonalization-based inverse using \(X\) itself is an alternative only when \(X\) is diagonalizable and its exponential differential is invertible; it is not the current Fisher fallback.

\begin{corollary}[no series in $\ad_X$]\label{frame:cor:nodexp}
For the groups of Proposition~\ref{frame:prop:closed}(i)--(iv), $\mathsf{Y}$ is computed from one principal-logarithm Fr\'echet derivative (square roots and $d\times d$ solves) and one closed-form $A_U^{-1}$. With the group retraction, the chart velocity $v=\iota^{-1}h(\ad_X)\mathsf{Y}=\iota^{-1}L_{\log}(U,U\mathsf{Y})$ is needed only for diagnostics. The additive chart retraction does require it.
\end{corollary}

\subsection{Retraction, logarithm branch, and trust region}\label{frame:sec:retraction}

Three local matrix updates share the first derivative $-U\mathsf{Y}=-P_\mathcal{T}\Gamma$: the additive $R_{\mathrm{add}}(\eta)=U(I-\eta\mathsf{Y})$ followed by $\Log$ where defined; the exponential $R_{\exp}(\eta)=Ue^{-\eta\mathsf{Y}}$ followed by $\Log$; and the implementation's additive chart step $R_{\mathrm{chart}}(\eta)=e^{X-\eta V}$ with $V=h(\ad_X)\mathsf{Y}$. We use
\begin{equation}\label{frame:eq:grouplog}
\varphi_{\mathrm{new}}=\iota^{-1}\Pi_\g\Log\left(Ue^{-\eta s\mathsf{Y}}\right),
\end{equation}
with the trust factor $s$ of Lemma~\ref{frame:lem:trust}. Two reasons decide the choice. $Ue^{-\eta\mathsf{Y}}\in G$ for every subalgebra, whereas $U(I-\eta\mathsf{Y})\notin G$ for $\mathfrak{so}$, $\mathfrak{sl}$, $\mathfrak{sp}$. And $Ue^{-\eta\mathsf{Y}}$ is exactly the reference element $\exp(\varphi)\exp(\delta)$ against which the implementation certifies its candidate, so the certificate residual reduces to the accuracy of the logarithm plus storage rounding. We call \eqref{frame:eq:grouplog} the \texttt{group\_log} retraction.

\begin{proposition}[second-order comparison]\label{frame:prop:second}
Assume $\spec(X)\subset\mathcal{S}$, and let $W(\eta)$ be the solution of $W'=-h(\ad_W)\mathsf{Y}$, $W(0)=X$, on an interval on which $\spec(W)\subset\mathcal{S}$. Then $e^{W(\eta)}=Ue^{-\eta\mathsf{Y}}$, and
\begin{align}
W''(0)&=\frac{d}{dt}h(\ad_{X+tV})\mathsf{Y}\Big|_{t=0}=\tfrac16[[X,\mathsf{Y}],\mathsf{Y}]+O(\|X\|^2\|\mathsf{Y}\|^2),\label{frame:eq:W2}\\
\varphi_{\log}-\varphi_{\mathrm{chart}}&=\tfrac{\eta^2}{12}\iota^{-1}[[X,\mathsf{Y}],\mathsf{Y}]+O(\eta^2\|X\|^2\|\mathsf{Y}\|^2,\eta^3),\label{frame:eq:gap}\\
e^{X-\eta V}-Ue^{-\eta\mathsf{Y}}&=-\tfrac{\eta^2}{2}L_{\exp}(X,W''(0))+O(\eta^3),\\
R_{\mathrm{add}}(\eta)-R_{\exp}(\eta)&=-\tfrac{\eta^2}{2}U\mathsf{Y}^2+O(\eta^3).
\end{align}
At $X=0$ the chart step coincides with $R_{\exp}$ exactly.
\end{proposition}
\begin{proof}
On the strip $\ad_W$ has no eigenvalue in $2\pi i(\mathbb{Z}\setminus\{0\})$, so $h(\ad_W)=\Psi_L(\ad_W)^{-1}$ is defined. By Lemma~\ref{frame:lem:dexp}, $M=e^{W}$ satisfies $M'=M\Psi_L(\ad_W)(W')=-M\mathsf{Y}$ with $M(0)=U$, and uniqueness for this linear equation gives $M(\eta)=Ue^{-\eta\mathsf{Y}}$. Also $W'(0)=-h(\ad_X)\mathsf{Y}=-V$. Differentiating again, $W''(0)=-\frac{d}{d\eta}h(\ad_{W(\eta)})\mathsf{Y}|_0$, which equals the stated derivative along $X+tV$. With $h(z)=1+z/2+z^2/12+O(z^4)$ and $V=\mathsf{Y}+\frac12[X,\mathsf{Y}]+O(\|X\|^2\|\mathsf{Y}\|)$, the linear term contributes $\frac12[V,\mathsf{Y}]=\frac14[[X,\mathsf{Y}],\mathsf{Y}]+O(\|X\|^2\|\mathsf{Y}\|^2)$ and the quadratic term contributes $\frac1{12}([V,[X,\mathsf{Y}]]+[X,[V,\mathsf{Y}]])=-\frac1{12}[[X,\mathsf{Y}],\mathsf{Y}]+O(\|X\|^2\|\mathsf{Y}\|^2)$; the sum is $\frac16[[X,\mathsf{Y}],\mathsf{Y}]$. The chart step is linear in $\eta$ in the chart and matches $W$ at order $\eta$, which gives \eqref{frame:eq:gap} and the third line. The last line is the expansion $e^{-\eta\mathsf{Y}}=I-\eta\mathsf{Y}+\frac{\eta^2}{2}\mathsf{Y}^2+O(\eta^3)$. At $X=0$, $V=h(0)\mathsf{Y}=\mathsf{Y}$ and $e^{-\eta V}=Ue^{-\eta\mathsf{Y}}$ with $U=I$.
\end{proof}

The chart retraction's product-residual certificate compares two updates that agree to first order, so its leading discrepancy is approximately \(\eta^2\|U[[X,\mathsf{Y}],\mathsf{Y}]\|/(12\|U\|)\). A fixed residual tolerance can therefore reduce a step already bounded by a tangent trust radius. The group-logarithm update removes this truncation discrepancy, but its candidate must still pass logarithm reconstruction, storage-rounding, chart-bound, and descent-pairing checks. It does not guarantee acceptance or decrease of the nonlinear training loss.

\begin{proposition}[logarithm branch and existence]\label{frame:prop:branch}
\begin{enumerate}
\item[(L1)] $\Log(e^X)=X$ if and only if $\spec(X)\subset\mathcal{S}$.
\item[(L2)] With $W$ the solution of Proposition~\ref{frame:prop:second}, if $\spec(W(t))\subset\mathcal{S}$ for all $t\in[0,\eta]$, then the principal logarithm satisfies $\Log(Ue^{-\eta\mathsf{Y}})=W(\eta)$.
\item[(L3)] If the continuation leaves $\mathcal{S}$ and the product has no eigenvalue on the closed negative real axis, $\Log$ returns a re-charted logarithm $X'$ with $e^{X'}=Ue^{-\eta\mathsf{Y}}$ exactly; for diagonalizable $W$, $X'=W(\eta)-Q\operatorname{diag}(2\pi ik_j)Q^{-1}$, and $\|X'-X\|_2\ge2\pi-O(\eta)$. The realized frame, and hence the loss when the loss depends only on the realized frame, is unchanged.
\item[(L4)] Existence can fail: $Ue^{-\eta\mathsf{Y}}$ can leave $\exp(\g)$. A real matrix has a real logarithm if and only if it is nonsingular and every Jordan block of each negative eigenvalue occurs an even number of times \cite{frame:Culver1966}; two distinct negative eigenvalues therefore admit none.
\end{enumerate}
\end{proposition}
\begin{proof}
(L1) is uniqueness of the logarithm with spectrum in $\mathcal{S}$ \cite[Thm.~1.31]{frame:Higham2008}. (L2) By Proposition~\ref{frame:prop:second}, $Ue^{-\eta\mathsf{Y}}=e^{W(\eta)}$, and (L1) applied to $W(\eta)$ gives $\Log(e^{W(\eta)})=W(\eta)$. (L3) $X'$ and $W(\eta)$ are two logarithms of the same matrix; for diagonalizable $W$ they share eigenvectors and their eigenvalues differ by integer multiples of $2\pi i$, with at least one nonzero multiple when $W(\eta)$ has left $\mathcal{S}$. (L4) is Culver's theorem.
\end{proof}

A candidate logarithm must reconstruct the target group element: the implementation checks \(\|e^{X_{\mathrm{new}}}-Ue^{-\eta\mathsf{Y}}\|/\|Ue^{-\eta\mathsf{Y}}\|\). Inverse scaling and squaring uses Denman--Beavers square roots \cite{frame:DenmanBeavers1976}; reconstruction is needed because convergence of an iteration alone does not certify the intended logarithm. Near the strip boundary, chart displacements can be amplified by \(\|h(\ad_X)\|\). For a normal real matrix whose extremal imaginary eigenvalues approach \(\pm i\pi\) with margin \(\delta'\), a corresponding adjoint resonance gives amplification of order \(\pi/\delta'\).

\begin{proposition}[strip certificate]\label{frame:prop:bendixson}
Every eigenvalue $\lambda$ of a real matrix $X$ satisfies $|\Imag \lambda|\le\|\operatorname{skew}X\|_2\le\|\operatorname{skew}X\|_F$, where $\operatorname{skew}X=(X-X^\top)/2$. Hence $\|\operatorname{skew}X_h\|_F<\pi-\delta$ certifies $\spec(X_h)\subset\mathcal{S}$ with margin $\delta$.
\end{proposition}
\begin{proof}
Let $Xv=\lambda v$ with $v\in\C^K$, $v^*v=1$. Then $\lambda=v^*Xv$ and, $X$ being real, $\bar\lambda=v^*X^\top v$. So $\Imag \lambda=(\lambda-\bar\lambda)/(2i)=v^*(\operatorname{skew}X)v/i$, and $|v^*(\operatorname{skew}X)v|\le\|\operatorname{skew}X\|_2\le\|\operatorname{skew}X\|_F$.
\end{proof}

\begin{lemma}[trust region]\label{frame:lem:trust}
With $s=\min(1,R/\|\eta\mathsf{Y}\|_F)$, the embedded norm taken over all $H$ blocks, and $Z=-\eta s\mathsf{Y}$, the relative frame displacement satisfies $\|U^{-1}\Delta U\|_F=\|e^Z-I\|_F\le e^{R}-1\le Re^R$.
\end{lemma}
\begin{proof}
$\|e^Z-I\|_F\le\sum_{k\ge1}\|Z\|_F^k/k!=e^{\|Z\|_F}-1$ by submultiplicativity, and $\|Z\|_F\le R$.
\end{proof}

The implemented trust rule uses this Frobenius norm of the body step. With \texttt{group\_log}, backtracking can still be required by a logarithm failure, a chart-bound violation, an inaccurate reconstruction after rounding, or a positive covector-displacement pairing. The last check is a first-order descent test, not a reevaluation of the loss. The closed-form Frobenius direction is used only when every block passes the strip margin and \(\kappa_2(U_h)\le10^3\); other rows use the damped series solver, whose condition gate is \(10^6\). For normal \(X_h\) with real spectrum, the singular values of its exponential differential are exponential divided differences and \(\kappa(D)\le\kappa_2(U_h)\), so this frame bound implies \(\kappa(G_0)\le10^6\). Outside that case it is a routing threshold, not an equivalent chart-condition certificate.

\begin{remark}[rejected rows must keep their chart]\label{frame:rem:rejected}
A rejected row must keep its original coordinates. Setting its step generator to zero is insufficient if a logarithm is still evaluated: off the principal strip, \(\Log(e^X)\ne X\), even though the realized matrix is unchanged. The stager therefore excludes rejected rows from retraction and retains their stored chart exactly.
\end{remark}

\paragraph{Positional factor.} The free coordinates are the two tables, so the loss is $\ell(\varphi_{\mathrm{tok}},\theta_{\mathrm{pos}})$ and the existing \texttt{pullback\_group} metric is the row-wise sum $\sum_v\|\delta E_v\|_F^2+\sum_p\|\delta F_p\|_F^2$ pulled back to table rows. Each row gets its own partial covector, recovered from its own row gradient by Proposition~\ref{frame:prop:recovery}, and its own step. All candidates are staged from the pre-step state, so the first-order composite motion is $\delta U_i=-\eta(\Gamma_{E_v}F_{p}+E_v\Gamma_{F_p})$. This sum differs from assigning an independent Euclidean update to every occurrence frame. Shared factors constrain which collections of occurrence velocities can be represented by one simultaneous table update.

\subsection{Gauge behavior of the pullback step}\label{frame:sec:gauge-pullback}

\begin{proposition}[orthogonal-only equivariance]\label{frame:prop:orth}
Suppose $L$ is invariant under $E\mapsto gE$ with the belief tables co-transformed, so that $\Gamma(gE)=g^{-\top}\Gamma(E)$. For $\g=\gl$ let $\Phi(E,\Gamma)=E\exp(-\eta E^{-1}\Gamma)$ be the first-order pullback map. Then $\Phi(gE,g^{-\top}\Gamma)=g\Phi(E,\Gamma)$ for all $\Gamma$ if and only if $g^\top g=I$, and $\Phi(Eh,\Gamma h^{-\top})=\Phi(E,\Gamma)h$ for all $\Gamma$ if and only if $h^\top h=I$.
\end{proposition}
\begin{proof}
Invariance gives $\ip{\Gamma(gE)}{g\delta E}=\ip{\Gamma(E)}{\delta E}$, so $\Gamma(gE)=g^{-\top}\Gamma(E)$. At $gE$ the direction is $(gE)^{-1}g^{-\top}\Gamma=E^{-1}(g^\top g)^{-1}\Gamma$, which equals $\mathsf{Y}=E^{-1}\Gamma$ for every $\Gamma$ if and only if $g^\top g=I$. At $Eh$ it is $h^{-1}E^{-1}\Gamma h^{-\top}=h^{-1}\mathsf{Y} h^{-\top}$, while equivariance requires $Ehe^{-\eta h^{-1}\mathsf{Y} h}=Ee^{-\eta\mathsf{Y}}h$, i.e. the direction $h^{-1}\mathsf{Y} h$; the two agree for every $\mathsf{Y}$ if and only if $h^{-\top}=h$.
\end{proof}

\begin{proposition}[left-invariant alternative]\label{frame:prop:leftinv}
The left-invariant metric $\|E^{-1}\delta E\|_F^2$ gives the direction $\mathsf{Y}_L=E^\top\Gamma=c$ with no solve. It is left-equivariant under all of $\GL$ and right-equivariant only under orthogonal $h$. It differs from the pullback direction at first order: its $U$-velocity is $EE^\top\Gamma$.
\end{proposition}

\begin{theorem}[no bi-invariant metric]\label{frame:thm:milnor}
For $d\ge2$ no Riemannian metric on $\GL^+(d)$ is both left- and right-invariant. Consequently no algebra-only metric on the frame charts is equivariant under both the left gauge action and the right reference change by all of $\GL$.
\end{theorem}
\begin{proof}
A bi-invariant metric restricts at the identity to an inner product on $\gl(d)$ invariant under $\Ad_g$ for $g$ in the identity component. Differentiating, every $\ad_Y$ is skew-adjoint for that inner product, so every $\ad_Y$ has purely imaginary spectrum. Take $Y=\operatorname{diag}(1,-1,0,\dots,0)$: then $\ad_YE_{12}=2E_{12}$, a real nonzero eigenvalue, which is a contradiction. The general criterion is that a connected Lie group admits a positive-definite bi-invariant metric precisely when its adjoint image has compact closure \cite{frame:Milnor1976}; the direct contradiction above suffices for $\GL^+(d)$.
\end{proof}

On $G\subset\mathrm{O}(K)$ the Frobenius, left-invariant and right-invariant metrics coincide and are bi-invariant. The diagonal covariance family is not closed under arbitrary general linear congruences, and a decoder acting directly on local means can break the passive-site symmetry. The invariance statements here concern the corresponding full-covariance statistical family and a loss with the stated symmetry. Section~\ref{frame:sec:symmetries} shows how a belief-dependent Fisher metric differs from an algebra-only group metric.

\subsection{Cost of the \texorpdfstring{$U$}{U}-space step}\label{frame:sec:cost-pullback}

Table~\ref{frame:tab:cost-pullback} gives operation counts per row-block in float64 at $d=10$. The materialized-series route analyzed in the September 28 sizing example materializes a $100\times100$ series for $\Psi$ at the illustrative order $16$ and needs twelve $100\times100$ workspaces per row-block, which gave chunks of $69$ rows under that workspace budget. The $U$-space route needs about thirty $d\times d$ matrices (24 KB) and $2$ to $6$ square roots at $4$ to $8$ Denman--Beavers iterations each.

\begin{table}[h]
\centering
\small
\begin{tabular}{lll}
\toprule
Route & Dominant work & Flops per row-block\\
\midrule
Materialized series route & $15$ products of $100\times100$; assembly; eigensolve & $\approx3.8\times10^7$\\
$U$-space with recovery & $\exp$, pair Fr\'echet $\Log$, plain $\Log$, solve & $3\times10^5$ to $10^6$\\
$U$-space with gradient hook & as above without the Fr\'echet $\Log$ & $1.2\times10^5$ to $3\times10^5$\\
\bottomrule
\end{tabular}
\caption{Operation counts per row-block at $d=10$ (arithmetic estimates, not measured).}\label{frame:tab:cost-pullback}
\end{table}

\subsection{The belief Fisher: modeling choices}\label{frame:sec:em}

\subsubsection{The EM reading: \texorpdfstring{$\varphi$}{phi} stays a slow M-step parameter}

In a variational-EM interpretation \cite{inference:dempster1977,frame:NealHinton1998}, belief revision plays the role of an approximate E-step and the outer parameter update a generalized M-step. The frame tables belong to the outer parameters. Changing their metric does not turn the predictive forward route into exact coordinate ascent on one evidence lower bound, nor does it make the frames inference variables inside each layer. The construction here specifies only the preconditioner and retraction of the outer frame gradient.

Amari's natural gradient for the parameter of a statistical model uses the Fisher information of the model family indexed by that parameter \cite{frame:Amari1998,frame:AmariNagaoka2000}. The frame parameter enters the generative model through the reference prior of each occurrence,
\begin{equation}\label{frame:eq:prior-family}
\Normal\left(U^{-1}\mu_v,U^{-1}\Sigma_vU^{-\top}\right)=U^{-1}\cdot q_v,
\end{equation}
where $q_v=\Normal(\mu_v,\Sigma_v)$ is the prior-bank entry of the type. The metric of this note is the Fisher information of the family \eqref{frame:eq:prior-family} with respect to the frame, evaluated at the prior-bank entry, that is at layer $0$. It is a function of the parameters alone. For positional rows it also depends on which types occupy a column of the batch; these play the role of covariates, as the design matrix does in the Fisher information of a regression model, and the metric does not depend on the responses.

Two limitations of this choice should be stated plainly. First, the family \eqref{frame:eq:prior-family} is the prior factor of the complete-data model, not the marginal likelihood of the tokens. The frame also enters the likelihood through the attention transports and the write-back, and the Fisher information of the training objective itself is the output Fisher of Section~\ref{frame:sec:output}. The prior-family Fisher is therefore a metric on the frame parameters that is intrinsic to the generative model and deterministic; it is not claimed to approximate the Hessian of the loss. Second, the E-step posteriors of deeper layers are deliberately not used. They are inference-state quantities: a metric built from them would change within an E-step, depend on context, and read the previous write-back, which itself depends on $U$. Using them would also destroy the position independence of the type metric (Proposition~\ref{frame:prop:position}) and require a capture in the read path and a dense solve per touched type row. Pooling deeper-layer reads is a possible extension that would let data inform the weight of the completion in Section~\ref{frame:sec:completion}; it is not part of the route analyzed here.

\subsubsection{The Gaussian Fisher metric and the orbit model}

On the full Gaussian manifold $\Normal_d$ with coordinates $(m,S)$ the Fisher--Rao metric \cite{frame:Skovgaard1984} is
\begin{equation}\label{frame:eq:gaussian-fisher}
g_{\mathrm{FR}}\big((\delta m,\delta S),(\delta m,\delta S)\big)=\delta m^\top S^{-1}\delta m+\tfrac12\tr\left(S^{-1}(\delta S)S^{-1}\delta S\right).
\end{equation}
This is the metric inverted by the implementation's \texttt{retraction.natural\_gradient} for beliefs. On the diagonal family it restricts to $\sum_k\delta m_k^2/s_k+\frac12\sum_k(\delta s_k/s_k)^2$.

\begin{lemma}[Fisher information of a submodel]\label{frame:lem:submodel}
If $\vartheta\mapsto p_{f(\vartheta)}$ is a smooth submodel of a regular family $\{p_\omega\}$, its Fisher information is the pullback $f^*g$ of the Fisher information $g$ of the ambient family.
\end{lemma}
\begin{proof}
The score of the submodel is $d\log p_{f(\vartheta)}=d_\omega\log p\circ df$, and the Fisher information is the covariance of the score.
\end{proof}

The diagonal family models the local prior, while its frame orbit \(\mathcal M_q=\{U^{-1}\cdot q:U\in G\}\) lies in the full Gaussian family: \(C=U^{-1}\Sigma U^{-\top}\) is generally full. The frame metric therefore pulls back the full Gaussian Fisher even when the stored local covariance is diagonal. Replacing \(C\) by its diagonal before computing the Fisher would define a different statistical family and would discard directions visible to the comparison maps. A later projection of predictive covariance back to a diagonal local family is a separate operation and does not justify that replacement.

\subsection{The frame Fisher metric}\label{frame:sec:fisher}

Throughout this section fix one occurrence, drop the index $i$, and write $r=\Normal(\xi,C)$ for its reference belief. Perturb the frame by a body generator, $U\mapsto Ue^{tY}$.

\begin{theorem}[frame Fisher]\label{frame:thm:fisher}
The pullback of \eqref{frame:eq:gaussian-fisher} under $Y\mapsto(Ue^{Y})^{-1}\cdot q$ at $Y=0$ is
\begin{equation}\label{frame:eq:gU}
g_U(Y,Y)=\xi^\top Y^\top C^{-1}Y\xi+\tfrac12\tr\left(C^{-1}(YC+CY^\top)C^{-1}(YC+CY^\top)\right).
\end{equation}
Its polarization is $g_U(Y_1,Y_2)=\ip{Y_1}{\mathcal{F}(Y_2)}$ with the Frobenius-self-adjoint operator
\begin{equation}\label{frame:eq:F}
\mathcal{F}(Y)=C^{-1}Y(C+\xi\xi^\top)+Y^\top.
\end{equation}
\end{theorem}
\begin{proof}
$(Ue^{tY})^{-1}=e^{-tY}U^{-1}$, so the perturbed moments are $e^{-tY}\xi$ and $e^{-tY}Ce^{-tY^\top}$, with derivatives $\delta\xi=-Y\xi$ and $\delta C=-(YC+CY^\top)$ at $t=0$. Substituting into \eqref{frame:eq:gaussian-fisher} gives \eqref{frame:eq:gU}. For the operator, the mean term is $\tr(Y^\top C^{-1}Y\xi\xi^\top)=\ip{Y}{C^{-1}Y\xi\xi^\top}$. Expanding the trace term into four products, $\tr(C^{-1}YCC^{-1}YC)=\tr(Y^2)$, $\tr(C^{-1}CY^\top C^{-1}CY^\top)=\tr(Y^2)$, and the two cross terms each equal $\tr(Y^\top C^{-1}YC)=\ip{Y}{C^{-1}YC}$. Half the sum is $\tr(Y^2)+\ip{Y}{C^{-1}YC}=\ip{Y}{Y^\top}+\ip{Y}{C^{-1}YC}$. Adding the mean term gives $\ip{Y}{\mathcal{F}(Y)}$. Self-adjointness: $\ip{Y_1}{C^{-1}Y_2M}=\tr(Y_1^\top C^{-1}Y_2M)=\ip{C^{-1}Y_1M}{Y_2}$ for symmetric $C$, $M$, and $\ip{Y_1}{Y_2^\top}=\ip{Y_1^\top}{Y_2}$.
\end{proof}

\begin{proposition}[whitened form]\label{frame:prop:whitened}
For any factor $LL^\top=C$ set $Z=L^{-1}YL$ and $\omega=L^{-1}\xi$. Then
\begin{equation}\label{frame:eq:whitened}
g_U(Y,Y)=\|Z\omega\|^2+2\|\sym Z\|_F^2=\ip{Z}{\mathcal{L}(Z)},\qquad \mathcal{L}(Z)=Z(I+\omega\omega^\top)+Z^\top,
\end{equation}
and $\ip{Y}{C^{-1}YC}=\|Z\|_F^2$. Covectors whiten as $R=L^\top G_YL^{-\top}$. The result is unchanged under $L\mapsto LQ$ with $Q$ orthogonal, and $a:=|\omega|^2=\xi^\top C^{-1}\xi=\mu^\top\Sigma^{-1}\mu$ does not depend on the frame.
\end{proposition}
\begin{proof}
$YC+CY^\top=LZL^{-1}LL^\top+LL^\top L^{-\top}Z^\top L^\top=L(Z+Z^\top)L^\top$, so $C^{-1}(YC+CY^\top)=L^{-\top}(Z+Z^\top)L^\top$, whose squared trace is $\|Z+Z^\top\|_F^2=4\|\sym Z\|_F^2$. Also $Y\xi=LZ\omega$ and $\xi^\top Y^\top C^{-1}Y\xi=\|Z\omega\|^2$. The identity $2\|\sym Z\|_F^2=\ip{Z}{Z}+\ip{Z}{Z^\top}$ gives $\mathcal{L}$. With $Y=LZL^{-1}$, $\tr(Y^\top C^{-1}YC)=\tr(L^{-\top}Z^\top ZL^\top)=\|Z\|_F^2$. For covectors, $\ip{G_Y}{LZL^{-1}}=\ip{L^\top G_YL^{-\top}}{Z}$. Under $L\mapsto LQ$, $Z\mapsto Q^\top ZQ$ and $\omega\mapsto Q^\top\omega$, which preserve both terms. Finally $\xi^\top C^{-1}\xi=\mu^\top U^{-\top}U^\top\Sigma^{-1}UU^{-1}\mu$.
\end{proof}

Write $\mathcal{P}(Y)=C^{-1}YC$ for the whitened Frobenius operator, $\ip{Y}{\mathcal{P}Y}=\|Z\|_F^2$.

\begin{proposition}[positivity and kernel]\label{frame:prop:kernel}
$g_U\ge0$, and $g_U(Y,Y)=0$ if and only if $\sym Z=0$ and $Z\omega=0$. Hence
\begin{equation}\label{frame:eq:kernel}
\ker g_U=\stab(\xi,C)=\{Y\in\gl(d):Y\xi=0,\ YC+CY^\top=0\},
\end{equation}
the Lie algebra of the stabilizer of $r$ under the linear action. Its dimension is $(d-1)(d-2)/2$ if $\xi\neq0$ and $d(d-1)/2$ if $\xi=0$; at $d=10$ these are $36$ and $45$ of $100$ per block.
\end{proposition}
\begin{proof}
Both terms of \eqref{frame:eq:whitened} are nonnegative and vanish exactly under the stated conditions; transforming back, $\sym Z=0$ is $YC+CY^\top=0$ and $Z\omega=0$ is $Y\xi=0$. These are the linearizations of $A\xi=\xi$, $ACA^\top=C$. In whitened coordinates the kernel consists of skew $Z$ annihilating $\omega$, which is $\mathfrak{so}(d-1)$ when $\omega\neq0$ and $\mathfrak{so}(d)$ when $\omega=0$.
\end{proof}

\begin{theorem}[generalized spectrum]\label{frame:thm:spectrum}
Relative to $\mathcal{P}$, the generalized eigenvalues of $\mathcal{F}$ are
\begin{equation}\label{frame:eq:spectrum}
\{2+a\}\cup\left\{\tfrac12\left(2+a+\sqrt{a^2+4}\right)\right\}^{d-1}\cup\left\{\tfrac12\left(2+a-\sqrt{a^2+4}\right)\right\}^{d-1}\cup\{2\}^{d(d-1)/2}\cup\{0\}^{(d-1)(d-2)/2},
\end{equation}
with multiplicities as exponents.
\end{theorem}
\begin{proof}
By Proposition~\ref{frame:prop:whitened}, $Y\mapsto Z$ is an isometry from $(\gl(d),\ip{\cdot}{\mathcal{P}\cdot})$ to $(\gl(d),\ip{\cdot}{\cdot})$ carrying $\mathcal{F}$ to $\mathcal{L}$, so the generalized spectrum is the spectrum of $\mathcal{L}$. If $\omega\neq0$, choose an orthonormal basis with $e_1=\omega/|\omega|$ (allowed by the $Q$-invariance); then $\mathcal{L}(Z)=Z+aZe_1e_1^\top+Z^\top$. The entry $z_{11}$ is mapped to $(2+a)z_{11}$. For $j\neq1$, the pair $(z_{1j},z_{j1})$ is mapped by $\begin{pmatrix}1&1\\1&1+a\end{pmatrix}$, with trace $2+a$, determinant $a$ and eigenvalues $\frac12(2+a\pm\sqrt{a^2+4})$. On the block $Z_{\perp\perp}$ of indices $j,k\neq1$, $\mathcal{L}(Z)=Z+Z^\top$, with eigenvalue $2$ on its $d(d-1)/2$ symmetric and $0$ on its $(d-1)(d-2)/2$ skew directions. The counts add to $1+2(d-1)+(d-1)^2=d^2$. If $\omega=0$ then $a=0$ and the formula still holds, with the $d-1$ lower pair eigenvalues equal to $0$.
\end{proof}

The near-kernel values \(\frac12(2+a-\sqrt{a^2+4})\sim a/2\) as \(a\to0\) identify weakly visible directions in addition to the exact stabilizer. Their size depends on the local Mahalanobis magnitude \(a\), which can change during training.

\begin{proposition}[restriction to subalgebras]\label{frame:prop:subalg}
For a subalgebra $\g$ the metric is the restriction of $g_U$, with kernel $\g\cap\stab(\xi,C)$. For $\mathfrak{so}(d)$ the kernel is $\{Y:[Y,C]=0,\ Y\xi=0\}$; it is trivial for generic $C$, and of dimension $(d-1)(d-2)/2$ for $C=cI$ and $\xi\neq0$. For $\mathfrak{sp}(2m)$ with distinct symplectic eigenvalues, $\mathfrak{sp}\cap\mathfrak{o}(C)$ is the Lie algebra of an $m$-torus, so the kernel is $m$-dimensional for $\xi=0$ and trivial for generic $\xi$. For \texttt{block\_glk} with block-diagonal local covariance, $C$ is block-diagonal and the metric is exactly block-diagonal across heads, each block being the Fisher of the marginal $(\xi_h,C_h)$.
\end{proposition}
\begin{proof}
For skew $Y$, $YC+CY^\top=YC-CY=[Y,C]$. A matrix commuting with a symmetric $C$ of distinct eigenvalues is diagonal in its eigenbasis, and a skew diagonal matrix vanishes. For $C=cI$ the condition reduces to $Y\xi=0$. The symplectic statement follows from Williamson's normal form \cite{frame:Williamson1936}: a symplectic congruence brings $C$ to $\operatorname{diag}(\Lambda,\Lambda)$, whose symplectic-orthogonal stabilizer is a product of $m$ circle groups when the entries of $\Lambda$ are distinct. For the block statement, with block-diagonal $Y$ and $C$, $\tr(Y^\top C^{-1}YM)$ reads only the diagonal blocks of $M=C+\xi\xi^\top$, and $\tr(Y^2)$ splits by blocks.
\end{proof}

\begin{proposition}[domination of the read path]\label{frame:prop:dominate}
Every regular Gaussian submodel obtained from $(\xi,C)$ by a fixed linear map, such as the projected moments $(S\xi,SCS^\top)$ for query and key, the content $T\xi$, or $M_hCM_h^\top$, has a frame Fisher dominated by $g_U$ in the Loewner order.
\end{proposition}
\begin{proof}
For a parameter-independent observation map, the score of the observed family is the conditional expectation of the full score given the observation. Conditional Jensen gives a covariance matrix no larger in Loewner order. Applying this data-processing statement to each regular linear Gaussian observation gives the result \cite{frame:Cencov1982,frame:AmariNagaoka2000}. A singular image can instead be treated on its fixed support when that support yields a regular family.
\end{proof}

Thus $g_U$ bounds from above the Fisher information of each such fixed read map. This comparison uses the full reference law and does not identify a Fisher metric of the complete predictor.

\subsection{Symmetries of the frame Fisher}\label{frame:sec:symmetries}

\begin{proposition}[left gauge invariance]\label{frame:prop:left}
Under the left gauge action $(\mu,\Sigma,U)\mapsto(g\mu,g\Sigma g^\top,gU)$ the metric is invariant; for an invariant loss, the direction obtained with a covariant positive-definite completion is invariant as well.
\end{proposition}
\begin{proof}
The action fixes $r=U^{-1}\cdot q$ and the body generator: $(gU)^{-1}\delta(gU)=U^{-1}\delta U$. The metric \eqref{frame:eq:gU} depends only on $r$ and $Y$. For a gauge-invariant loss the body covector $G_Y$ is invariant as well, and so is the solution of the completed metric equation. Without completion, a quotient direction requires a covector annihilating the stabilizer.
\end{proof}

This is a property of the formula as a function of $(\mu,\Sigma,U)$. It is not a symmetry of the diagonal table parametrization unless $g$ is monomial.

\begin{proposition}[right equivariance]\label{frame:prop:right}
Under $U\mapsto Uh$, $g_{Uh}(h^{-1}Yh,h^{-1}Yh)=g_U(Y,Y)$. Covectors transform coadjointly, $G_Y\mapsto h^\top G_Yh^{-\top}$, and a covariantly completed natural direction transforms as $\mathsf{Y}\mapsto h^{-1}\mathsf{Y} h$.
\end{proposition}
\begin{proof}
$Uhe^{t h^{-1}Yh}=Ue^{tY}h$, so the perturbed reference beliefs at $Uh$ are $h^{-1}\cdot\big((Ue^{tY})^{-1}\cdot q\big)$, the pushforward of the family at $U$ by the fixed bijection $h^{-1}$. Fisher information is invariant under a parameter-free bijective change of the sample space. The covector law follows from $\ip{G'}{h^{-1}Yh}=\ip{G}{Y}$ for all $Y$, and the direction law from the two transformation laws.
\end{proof}

\begin{proposition}[right invariance for a fixed belief]\label{frame:prop:rightinv}
Let $h_q$ denote \eqref{frame:eq:gU} evaluated at $U=I$, i.e. with $(\xi,C)=(\mu,\Sigma)$. Then
\begin{equation}\label{frame:eq:rightinv}
g_U(Y,Y)=h_q(\Ad_UY,\Ad_UY)=h_q((\delta U)U^{-1},(\delta U)U^{-1}).
\end{equation}
For fixed $q$ the frame metric is therefore the right-invariant positive semidefinite metric on $G$ with identity value $h_q$. Its kernel at $U$ is $\{WU:W\in\stab(q)\}$; it descends to $\Stab_q\backslash G$, and $\Stab_qU\mapsto U^{-1}\cdot q$ is an isometric embedding onto the orbit in $(\Normal_d,g_{\mathrm{FR}})$.
\end{proposition}
\begin{proof}
$Ue^{tY}=e^{t\Ad_UY}U$, so $(Ue^{tY})^{-1}\cdot q=U^{-1}\cdot(e^{-t\Ad_UY}\cdot q)$, the pushforward by the fixed $U^{-1}$ of the family whose Fisher at $t=0$ is $h_q(\Ad_UY)$. Right invariance is the statement that $(\delta U)U^{-1}$ is unchanged under $U\mapsto Uk$. The orbit map is constant on cosets $\Stab_qU$, and by Lemma~\ref{frame:lem:submodel} its pullback of $g_{\mathrm{FR}}$ is $g_U$, whose kernel is exactly the tangent space to the coset.
\end{proof}

For $\GL(d)$ the orbits are the level sets $\{(m,S):m^\top S^{-1}m=a\}$; transitivity follows by whitening two beliefs and rotating one whitened mean onto the other. They have dimension $\frac12(d^2+3d)-1$ and are hypersurfaces for $a>0$. For $a=0$ the orbit is the cone of positive definite matrices with $g=\frac12\tr((C^{-1}\delta C)^2)$, one half of the affine-invariant metric \cite{frame:Skovgaard1984}. This is the frame's own geometry: moving the frame moves the reference belief along its Mahalanobis orbit, measured in Fisher--Rao distance. The belief metric escapes Theorem~\ref{frame:thm:milnor} because it is not a metric on $G$ alone: it depends on $q$, which moves with $U$ under the left action.

\begin{remark}[relation to Amari's natural gradient on matrices]\label{frame:rem:amari-bss}
Amari's natural gradient on the space of nonsingular matrices for blind source separation uses a right-invariant metric in $(\delta U)U^{-1}$ with identity value the Frobenius norm \cite{frame:Amari1998}. Proposition~\ref{frame:prop:rightinv} has the same invariance with identity value $h_q$. With the completion of Section~\ref{frame:sec:completion} at its default, and for a standard local belief $q=\Normal(0,I)$, the metric is $2\|(\delta U)U^{-1}\|_F^2$, twice Amari's; a general $q$ whitens and weights it. This follows by substituting the standard Gaussian moments into the completed metric.
\end{remark}

\begin{proposition}[covariance under the split redundancy]\label{frame:prop:split}
Use for each type row and each positional row the corresponding diagonal block of the batch Fisher of Section~\ref{frame:sec:batch}, completed as in Section~\ref{frame:sec:completion}. Then the realized-frame velocity $W_i^{\mathrm{tot}}=(\delta U_i)U_i^{-1}$ produced by the natural step does not depend on the representative $(E_vh,h^{-1}F_p)$ of the split redundancy. If the steps are retracted by the group exponential $E_v\leftarrow e^{-\eta W_v}E_v$, $F_p\leftarrow e^{-\eta W'_p}F_p$ with unit trust factor, the updated realized frames are also representative-independent.
\end{proposition}
\begin{proof}
The per-occurrence metric (Fisher plus completion) is a function of $q_{v_i}$ and of $W_i^{\mathrm{tot}}$ only: by \eqref{frame:eq:rightinv} the Fisher part is $h_{q_{v_i}}(W_i^{\mathrm{tot}})$, and the completion of Section~\ref{frame:sec:completion} is defined through $Z_i=L_i^{-1}U_i^{-1}W_i^{\mathrm{tot}}U_iL_i$ with $L_i=U_i^{-1}\Sigma_{v_i}^{1/2}$, which equals $\Sigma_{v_i}^{-1/2}W_i^{\mathrm{tot}}\Sigma_{v_i}^{1/2}$. The type block is the restriction of this metric and of the loss differential to the subspace of realized perturbations $\{e^{tW}E_vF_p\}$, and the positional block to $\{E_ve^{tW'}F_p\}$. Under the change of representative these subspaces are the same sets of realized perturbations: $e^{tW}(E_vh)(h^{-1}F_p)=e^{tW}E_vF_p$ and $(E_vh)e^{tW''}(h^{-1}F_p)=E_ve^{t\Ad_hW''}F_p$. A natural gradient restricted to a fixed subspace, with a metric and a covector defined on that subspace, is intrinsic, so the realized velocity in each block is independent of the coordinates used for it. The finite-step statement follows from $e^{-\eta W}(E_vh)e^{-\eta W''}h^{-1}F_p=e^{-\eta W}E_ve^{-\eta\Ad_hW''}F_p$ and $\Ad_hW''_{\mathrm{nat}}=W'_{\mathrm{nat}}$.
\end{proof}

Under AdamW, and under the pullback, the realized trajectory therefore depends on the arbitrary split of each frame into type and positional factors, while under the Fisher route it does not. One caveat applies to the finite step: the implementation's trust rule measures the Frobenius norm of the body generator, $\|\eta\mathsf{Y}_{\mathrm{body}}\|_F$ with $\mathsf{Y}_{\mathrm{body}}=E_v^{-1}W_vE_v$, and that norm changes under $E_v\mapsto E_vh$. When the trust region binds, the finite step is representative-dependent unless the trust norm is replaced by the metric norm. The velocity statement is unaffected.

\subsection{Which beliefs define the metric: the batch Fisher}\label{frame:sec:batch}

\begin{proposition}[position independence of the type metric]\label{frame:prop:position}
Perturb the stored type factor, $E_v\mapsto E_ve^{Y}$. The Fisher of every occurrence $i$ of type $v$ equals the Fisher $g_{E_v}(Y,Y)$ of the middle-fiber belief $r_v^0=E_v^{-1}\cdot q_v=\Normal(\xi^0_v,C^0_v)$, independently of the position. In the spatial generator $W$ of $E_v\mapsto e^{W}E_v$ the metric is $h_{q_v}(W,W)$, a function of the stored row $(\mu_v,s_v)$ only.
\end{proposition}
\begin{proof}
$E_ve^YF_p=U_ie^{\Ad_{F_p^{-1}}Y}$. By Proposition~\ref{frame:prop:right} with $h=F_p^{-1}$, $g_{U_i}(\Ad_{F_p^{-1}}Y)=g_{U_iF_p^{-1}}(Y)=g_{E_v}(Y)$. The spatial statement is \eqref{frame:eq:rightinv} with $U=E_v$ and $W=\Ad_{E_v}Y$.
\end{proof}

Hence the per-occurrence encoded beliefs and the per-type prior give the same metric, $\mathbb{E}_{i\in v}g_i=g_{r^0_v}$ exactly, and no pooling or forward statistics are needed for type rows.

\paragraph{The batch Fisher.} Perturb all factors spatially, $E_v\mapsto e^{W_v}E_v$ and $F_p\mapsto e^{W'_p}F_p$. Then $W_i^{\mathrm{tot}}=W_{v_i}+\Ad_{E_{v_i}}W'_{p_i}$ and the occurrence Fisher is $h_{q_{v_i}}(W_i^{\mathrm{tot}})$. Treat the reference priors attached to input occurrences as independent, so that information adds, and normalize by the number $\Nsc$ of scored tokens used in the cross-entropy denominator. Input occurrences can include context positions whose targets are not scored; the implemented metric counts those occurrences as design information:
\begin{equation}\label{frame:eq:batch}
\mathcal{G}=\frac{1}{\Nsc}\sum_ih_{q_{v_i}}\left(W_{v_i}+\Ad_{E_{v_i}}W'_{p_i}\right).
\end{equation}
Its diagonal blocks are $\frac{n_v}{\Nsc}h_{q_v}(W_v)$ for type row $v$, where $n_v$ is the number of input occurrences of $v$, and $\frac{B}{\Nsc}\cdot\frac1B\sum_bg_{r^0_{v_b}}(W'_p)$ for positional row $p$, using $h_q(\Ad_EW')=g_{E^{-1}\cdot q}(W')$ and $v_b$ the type at $(b,p)$. The route uses these blocks, a block-diagonal approximation of $\mathcal{G}$. The natural steps are therefore count-normalized: the type direction is $(\Nsc/n_v)\bar h_v^{-1}G_{W_v}$, the natural normalization for this declared independent-prior batch model. Using the occurrence mean $\bar h_v$ alone would give count-proportional, SGD-like steps and is not the Fisher of the batch model. The dropped cross blocks couple type row $v_i$ with positional row $p_i$, and $\mathcal{G}$ is degenerate along the split redundancy ($W_v=\Ad_{E_v}Z$, $W'_p=-Z$ gives $W_i^{\mathrm{tot}}=0$), a direction that moves no realized frame and that the block-diagonal approximation does not see.

Count normalization gives each touched type a direction on the scale of its mean per-occurrence gradient. A type seen once receives a single-occurrence estimate; whether that scaling helps optimization depends on the loss, data distribution, and learning rate. The metric is a chosen independent-prior model, not the output Fisher of the autoregressive predictor.

\paragraph{Positional rows are pooled.} Column $p$ carries the types $v_b$, so its metric is a sum of Fishers at different Gaussians. Its kernel is $\bigcap_b\stab(r^0_{v_b})$, which is trivial already for two generic Gaussians: after whitening one of them, a skew generator that also preserves a generic second covariance commutes with a matrix of distinct eigenvalues and vanishes. With the completion of Section~\ref{frame:sec:completion} each summand is positive definite, so the pooled operator is positive definite for every batch.

\subsection{Completion on the stabilizer}\label{frame:sec:completion}

A loss can depend on the frame through write-back, local covariance projection, coordinate normalization, a multilayer perceptron, or the decoder in addition to the reference-prior family. Its covector need not annihilate that family's stabilizer. In that case the uncompleted Fisher equation has no solution for the full covector. A small isotropic ridge regularizes the linear system but assigns a potentially large step to those directions; choosing a completion states the additional geometry explicitly.

\begin{proposition}[covariant completion family]\label{frame:prop:completion}
A completion that depends only on $C$ and is invariant under $\mathrm{O}(d)$ conjugation of $Z$ adds a combination of $\|\sym Z\|_F^2$, $\|\asym Z\|_F^2$ and $(\tr Z)^2$. The family below uses the first two, with $\lambda_s\ge0$ and $\lambda_a>0$, making it positive definite:
\begin{equation}\label{frame:eq:completion}
\mathcal{M}_{\lambda_s,\lambda_a}(Y,Y)=g_U(Y,Y)+\lambda_s\|\sym Z\|_F^2+\lambda_a\|\asym Z\|_F^2.
\end{equation}
has the operator
\begin{equation}\label{frame:eq:M}
\mathcal{M}(Y)=C^{-1}Y(c_1C+\xi\xi^\top)+c_2Y^\top,\qquad c_1=1+\tfrac12(\lambda_s+\lambda_a),\qquad c_2=1+\tfrac12(\lambda_s-\lambda_a),
\end{equation}
and generalized spectrum relative to $\mathcal{P}$
\begin{equation}\label{frame:eq:M-spectrum}
\{2+\lambda_s+a\}\cup\left\{\tfrac12\left(2c_1+a\pm\sqrt{a^2+4c_2^2}\right)\right\}^{d-1}\cup\{2+\lambda_s\}^{d(d-1)/2}\cup\{\lambda_a\}^{(d-1)(d-2)/2}.
\end{equation}
Every member with these parameter restrictions has the symmetries of Section~\ref{frame:sec:symmetries}.
\end{proposition}
\begin{proof}
The $\mathrm{O}(d)$-conjugation-invariant quadratic forms on $\gl(d)$ are combinations of $\|\sym Z\|^2$, $\|\asym Z\|^2$ and $(\tr Z)^2$; the family uses the first two. From $\|\sym Z\|^2=\frac12(\|Z\|^2+\tr Z^2)$, $\|\asym Z\|^2=\frac12(\|Z\|^2-\tr Z^2)$, $\|Z\|^2=\ip{Y}{\mathcal{P}Y}$ and $\tr Z^2=\tr Y^2=\ip{Y}{Y^\top}$, the added operator is $\frac{\lambda_s+\lambda_a}{2}C^{-1}YC+\frac{\lambda_s-\lambda_a}{2}Y^\top$; adding \eqref{frame:eq:F} gives \eqref{frame:eq:M}. In whitened form $\mathcal{M}$ becomes $Z\mapsto Z(c_1I+\omega\omega^\top)+c_2Z^\top$, and the proof of Theorem~\ref{frame:thm:spectrum} applies with $z_{11}\mapsto(c_1+c_2+a)z_{11}$, pair matrix $\begin{pmatrix}c_1&c_2\\c_2&c_1+a\end{pmatrix}$, and $c_1\pm c_2$ on the symmetric and skew parts of $Z_{\perp\perp}$. The added terms are functions of $Z$ that are invariant under $Z\mapsto Q^\top ZQ$, and $Z$ is determined by $(C,Y)$ covariantly under the three actions, so the proofs of Section~\ref{frame:sec:symmetries} carry over.
\end{proof}

The condition number of $\mathcal{M}$ relative to $\mathcal{P}$ is the ratio of the extremes of \eqref{frame:eq:M-spectrum}. The smallest eigenvalue is $\min\{\lambda_a,\frac12(2c_1+a-\sqrt{a^2+4c_2^2})\}$ when $\lambda_a\le2+\lambda_s$, and $2+\lambda_s$ when $\lambda_a>2+\lambda_s$ (for $d\ge3$; the member $\lambda_a$ is absent at $d=2$). Indeed, since $2c_1-2(2+\lambda_s)=-2c_2$, the lower pair eigenvalue is at most $2+\lambda_s$ if and only if $a-2c_2\le\sqrt{a^2+4c_2^2}$, which for $a>0$ holds if and only if $c_2\ge0$, that is $\lambda_a\le2+\lambda_s$; for $c_2<0$ and $a>0$ both the lower pair eigenvalue and $\lambda_a$ exceed $2+\lambda_s$. The largest is $2+\lambda_s+a$ when $\lambda_a\le2$ (with $\lambda_s=0$), but for $\lambda_a>2$ it is $\frac12(2c_1+a+\sqrt{a^2+4c_2^2})$, which then exceeds $2+a$; a condition diagnostic must use $\max\{2+\lambda_s+a,\frac12(2c_1+a+\sqrt{a^2+4c_2^2})\}$ as numerator.

\begin{corollary}[the default completion is a Kronecker product]\label{frame:cor:kron}
The member $(\lambda_s,\lambda_a)=(0,2)$ leaves the Fisher unchanged on the symmetric (covariance-visible) directions and on $z_{11}$, adds the Fisher's own shear weight $2$ on every skew direction, and makes the transpose term vanish ($c_2=0$). The skew directions include the $d-1$ skew components $z_{1j}-z_{j1}$ of the pairs, which the mean sees ($Z\omega\neq0$). On each pair the spectrum changes from $\frac12(2+a\pm\sqrt{a^2+4})$ to $\{2,2+a\}$, so the lower pair eigenvalue $\approx a/2$ is raised to $2$: the completion does change the metric on directions the belief sees. Then
\begin{equation}\label{frame:eq:kron}
\mathcal{M}_{0,2}(Y)=C^{-1}Y(2C+\xi\xi^\top),\qquad \ip{Y}{\mathcal{M}_{0,2}Y}=2\|Z\|_F^2+\|Z\omega\|^2,
\end{equation}
which in row-major vectorization is $C^{-1}\otimes(2C+\xi\xi^\top)$. Its spectrum relative to $\mathcal{P}$ is $\{2\}^{d^2-d}\cup\{2+a\}^{d}$, with condition number $(2+a)/2$ for $d\ge2$; for $d=1$ the sole eigenvalue is $2+a$ and the condition number is $1$.
\end{corollary}
\begin{proof}
Set $c_1=2$, $c_2=0$ in \eqref{frame:eq:M} and \eqref{frame:eq:M-spectrum}: $z_{11}$ gives $2+a$, each pair gives $\frac12(4+a\pm a)\in\{2+a,2\}$, and the complement gives $2$. The multiplicities are $1+(d-1)=d$ for $2+a$ and $(d-1)+(d-1)^2=d^2-d$ for $2$. The vectorization uses $\vect(AYB)=(A\otimes B^\top)\vect Y$ in row-major order.
\end{proof}

The implementation fixes $\lambda_s=0$ and exposes $\lambda_a>0$ as \texttt{m\_phi\_fisher\_rotation\_weight}, with default $2$. The belief does not determine this weight, and optimality for the training loss is not established.

\subsection{Solvers}\label{frame:sec:solvers}

\begin{proposition}[single Gaussian, general completion]\label{frame:prop:woodbury}
The whitened equation $c_1Z+Z\omega\omega^\top+c_2Z^\top=R$ has the solution
\begin{equation}\label{frame:eq:woodbury}
Z=\mathcal{L}_0^{-1}(R)-\mathcal{L}_0^{-1}(w\omega^\top),\qquad w=\frac{b-\beta\omega(\omega^\top b)/(1+a(\alpha+\beta))}{1+a\alpha},\qquad b=\mathcal{L}_0^{-1}(R)\omega,
\end{equation}
where $\mathcal{L}_0^{-1}(R)=\sym R/(2+\lambda_s)+\asym R/\lambda_a$, $\alpha=\frac12(\frac1{2+\lambda_s}+\frac1{\lambda_a})$ and $\beta=\frac12(\frac1{2+\lambda_s}-\frac1{\lambda_a})$. It costs $O(d^2)$, uses no basis adapted to $\omega$, and is valid at $\omega=0$.
\end{proposition}
\begin{proof}
$\mathcal{L}_0(Z)=c_1Z+c_2Z^\top$ acts as $c_1+c_2=2+\lambda_s$ on symmetric and $c_1-c_2=\lambda_a$ on skew matrices, which gives $\mathcal{L}_0^{-1}$. Write $w=Z\omega$; then $Z=\mathcal{L}_0^{-1}(R-w\omega^\top)$. Since $\mathcal{L}_0^{-1}(w\omega^\top)=\alpha w\omega^\top+\beta\omega w^\top$, multiplying by $\omega$ gives $w=b-\alpha aw-\beta\omega(\omega^\top w)$. Contracting with $\omega^\top$ yields $\omega^\top w=\omega^\top b/(1+a(\alpha+\beta))$, and substituting back gives $w$.
\end{proof}

\begin{corollary}[type-row step at the stored row]\label{frame:cor:type}
Take the factor $L=E_v^{-1}S_v^{1/2}$ of $C^0_v$. Then $Z=S_v^{-1/2}W_vS_v^{1/2}$ and $\omega=S_v^{-1/2}\mu_v$, so the type solve runs in spatial coordinates at the local diagonal row with $R_{jk}=\sqrt{s_j/s_k}(G_W)_{jk}$ and $(W_v)_{jk}=\sqrt{s_j/s_k}Z_{jk}$, with no eigendecomposition. For the default completion, including the count factor,
\begin{equation}\label{frame:eq:type-step}
W_v=\frac{\Nsc}{n_v}S_vG_W\left(2S_v+\mu_v\mu_v^\top\right)^{-1}=\frac{\Nsc}{2n_v}\left(Q-\frac{(Q\mu_v)(S_v^{-1}\mu_v)^\top}{2+a_v}\right),\qquad Q_{jk}=\frac{s_j}{s_k}(G_W)_{jk},
\end{equation}
with $a_v=\mu_v^\top S_v^{-1}\mu_v$.
\end{corollary}
\begin{proof}
$L^{-1}YL=S_v^{-1/2}E_vYE_v^{-1}S_v^{1/2}=S_v^{-1/2}W_vS_v^{1/2}$ with $W_v=\Ad_{E_v}Y$, and $L^{-1}\xi^0_v=S_v^{-1/2}E_vE_v^{-1}\mu_v$. For the default completion, $\mathcal{M}_{0,2}$ in the spatial coordinates at $q_v$ reads $W\mapsto S_v^{-1}W(2S_v+\mu_v\mu_v^\top)$, so $W=S_vG_W(2S_v+\mu_v\mu_v^\top)^{-1}$. The Sherman--Morrison formula \cite{frame:ShermanMorrison1950} gives $(2S+\mu\mu^\top)^{-1}=\frac12\left(S^{-1}-S^{-1}\mu\mu^\top S^{-1}/(2+a)\right)$, and $SGS^{-1}=Q$.
\end{proof}

\begin{remark}[the Stein reduction]\label{frame:rem:stein}
For the isotropic whitened ridge, put $A=I+\omega\omega^\top$. Transposing $Z(A+\lambda I)+Z^\top=R$ and eliminating $Z^\top$ gives the Stein equation $Z-PZP=R^\top-PR$ with $P=A+\lambda I$. For $\lambda>0$ it is correct and invertible, with smallest singular value $2\lambda+\lambda^2$ but largest about $(1+a+\lambda)^2$, so it is worse conditioned than the direct form. Undamped it is not equivalent: the elimination adds $d(d-1)/2$ spurious null directions (for nonzero $\omega$, the ranks are $2d-1$ and $(d^2+3d-2)/2$, respectively), so it must not be used for a pseudo-inverse.
\end{remark}

\begin{proposition}[pooled positional operator]\label{frame:prop:pooled}
In row-major vectorization the completed positional block is
\begin{equation}\label{frame:eq:pooled}
\mathcal{M}_p=\frac1B\sum_bC_b^{-1}\otimes N_b+c_2K_{\mathrm{comm}},\qquad N_b=c_1C_b+\xi_b\xi_b^\top,
\end{equation}
with $(\xi_b,C_b)$ the middle-fiber moments of $v_b$ and $K_{\mathrm{comm}}$ the commutation matrix, and the positional step is $W'_p=(\Nsc/B)\mathcal{M}_p^{-1}G_{W'}$. For distinct Gaussians $\mathcal{M}_p$ is in general not a Kronecker product.
\end{proposition}
\begin{proof}
The formula is the sum of \eqref{frame:eq:M} over the column, expressed in the spatial coordinates of Proposition~\ref{frame:prop:position}, with the count factor from \eqref{frame:eq:batch}. The Van Loan--Pitsianis rearrangement \cite{frame:VanLoanPitsianis1993} maps \(\sum_b C_b^{-1}\otimes N_b\) to \(\sum_b\vect(C_b^{-1})\vect(N_b)^\top\). A nonzero operator is a single Kronecker product exactly when its full rearrangement has rank one. Generic distinct Gaussians give larger rank; proportional factors are a sufficient special case, not a necessary characterization of every possible cancellation. For nonzero \(c_2\), the rearrangement must also include the commutation term.
\end{proof}

The implementation forms the positional operator densely and uses a float64 Cholesky factorization after Jacobi scaling. Assembly costs \(O(Bd^4)\) per positional head and the factorization costs \(O(d^6)\). It accepts a solve when the normwise backward error \(\|\mathcal Mx-g\|/(\|\mathcal M\|_F\|x\|+\|g\|)\) is at most \(10^{-10}\). Replacing the spectral norm by the Frobenius norm avoids an SVD and can enlarge the denominator by at most \(\sqrt{d^2}=d\). The Jacobi-scaled spectrum is diagnostic and is computed on logging steps. These choices do not establish that dense Cholesky is the fastest solver at every block size or condition number.

\subsection{From the natural direction to a \texorpdfstring{$\varphi$}{phi} update}\label{frame:sec:update}

\begin{proposition}[chart direction and descent]\label{frame:prop:descent}
On a regular exponential chart, the chart direction and the body generator for staging are
\begin{equation}\label{frame:eq:chartdir}
v_\varphi=\Psi_R(\ad_{X_v})^{-1}W_v,\qquad \mathsf{Y}_{\mathrm{body}}=E_v^{-1}W_vE_v=\Psi_L(\ad_{X_v})v_\varphi,
\end{equation}
and likewise with $(\Theta_p,F_p,W'_p)$ for positional rows. In chart coordinates $v_\varphi=\mathcal{F}_\varphi^{-1}\nabla_\varphi\ell$ with $\mathcal{F}_\varphi=\Psi_R(\ad_{X_v})^*\mathcal{H}\Psi_R(\ad_{X_v})$, where $\mathcal{H}$ is the positive definite completed and count-normalized metric operator in spatial coordinates. Hence $v_\varphi$ is a descent direction whenever $\nabla_\varphi\ell\neq0$.
\end{proposition}
\begin{proof}
A chart velocity $\delta\varphi$ produces $W=\Psi_R(\ad_X)\delta\varphi$ by Lemma~\ref{frame:lem:dexp}, so the metric in chart coordinates is $\Psi_R^*\mathcal{H}\Psi_R$ and, by \eqref{frame:eq:covectors}, $\nabla_\varphi\ell=\Psi_R(\ad_X)^*G_W$. Then $\mathcal{F}_\varphi^{-1}\nabla_\varphi\ell=\Psi_R^{-1}\mathcal{H}^{-1}G_W=\Psi_R^{-1}W$. The body relation is $\Psi_L(\ad_X)=e^{-\ad_X}\Psi_R(\ad_X)$. Positive definiteness of $\mathcal{F}_\varphi$ gives $\ip{\nabla_\varphi\ell}{v_\varphi}>0$.
\end{proof}

\paragraph{Alternative chart maps without materializing $\Psi$.} Since $\Psi_R(z)^{-1}=z/(e^z-1)=\sum_kB_kz^k/k!$ with Bernoulli numbers $B_k$ ($B_1=-\frac12$, odd $B_k=0$ for $k\ge3$), $\Psi_R(\ad_X)^{-1}Y$ is a series with one commutator per term and radius $2\pi$; the adjoint is the same series in $X^\top$.

\begin{lemma}[Bernoulli tail]\label{frame:lem:bernoulli}
Let $\rho=\|X\|_F/\pi$, so that $\|\ad_X\|/(2\pi)\le\rho$ in the Frobenius operator norm. If $\rho<1$, the tail of the series after the even term $2n$ is bounded by $4\rho^{2n+2}/(1-\rho^2)\|Y\|_F$.
\end{lemma}
\begin{proof}
$\|\ad_XY\|_F\le2\|X\|_2\|Y\|_F\le2\|X\|_F\|Y\|_F$. From $|B_{2m}|=2(2m)!\zeta(2m)/(2\pi)^{2m}$ and $\zeta(2m)\le\zeta(2)=\pi^2/6$, $|B_{2m}|/(2m)!\le(\pi^2/3)(2\pi)^{-2m}<4(2\pi)^{-2m}$. Summing the geometric tail gives the bound.
\end{proof}

When $X$ is diagonalizable and the chart differential is invertible, a Daleckii--Krein form \cite[Ch.~3]{frame:Higham2008} with $X=P\Lambda P^{-1}$ gives
\begin{equation}\label{frame:eq:dk}
\Psi_R(\ad_X)^{-1}Y=P\left[(P^{-1}YP)\oslash\Psi_R(\lambda_j-\lambda_k)\right]P^{-1},\qquad \frac{1}{\Psi_R(z)}=\frac{e^{-z/2}}{\operatorname{sinhc}(z/2)},
\end{equation}
because \(\ad_X\) acts on \(PE_{jk}P^{-1}\) by \(\lambda_j-\lambda_k\). The adjoint is evaluated by the same analytic inverse function at \(X^\top\); for complex eigendecompositions its implementation must preserve the appropriate conjugate pairing and return the real operator. The Bernoulli and diagonalization constructions are mathematical alternatives. The current Fisher route instead uses the Fr\'echet derivative of the principal logarithm and rejects rows outside its certified strip.

\paragraph{Retraction and the certificate cap.} The chart retraction \(\varphi-\eta sv_\varphi\) agrees to first order with \(e^Xe^{-\eta s\mathsf{Y}_{\mathrm{body}}}=e^{-\eta sW}e^X\). A product-residual test nevertheless sees their second-order difference from Proposition~\ref{frame:prop:second}, so the accepted chart step can be smaller than the tangent trust radius permits. The group-logarithm candidate is
\begin{equation}\label{frame:eq:fisher-grouplog}
X_{\mathrm{new}}=\Log\left(e^Xe^{-\eta s\mathsf{Y}_{\mathrm{body}}}\right)=\Log\left(e^{-\eta sW}e^X\right).
\end{equation}
It removes the chart truncation discrepancy, subject to branch existence, reconstruction, and the other acceptance gates.

The $U$-space pullback uses the same interface with $W^{\mathrm{pb}}=G_W(EE^\top)^{-1}$, Euclidean descent on $E$. In spatial coordinates the Frobenius metric depends on the current factor, whereas the type-row Fisher metric depends on the fixed local prior moments. Pulling either back into stored chart coordinates reintroduces the exponential differential and therefore dependence on $\varphi$.

\subsection{The output Fisher alternative}\label{frame:sec:output}

The Fisher information of the training objective itself is the output Fisher
\begin{equation}\label{frame:eq:output}
\mathcal{F}^{\mathrm{out}}=\frac1N\sum_jJ_j^\top\left(\operatorname{diag}p_j-p_jp_j^\top\right)J_j,
\end{equation}
restricted to one $d\times d$ block of one row, where $p_j$ is the predictive distribution at position $j$ and $J_j$ the Jacobian of its logits. It is the Fisher information of the model output distribution, in contrast to the prior-family metric. For fixed conditioning contexts, draw independent sampled targets $\hat y_j\sim p_j$ and form $\hat g=N^{-1/2}\sum_j\nabla_\theta\log p_j(\hat y_j)$. Then $\mathbb E[\hat g\hat g^\top]=\mathcal F^{\mathrm{out}}$: cross terms vanish by the zero-mean score identity. Sampling a fresh autoregressive sequence changes the conditioning contexts and defines a different expectation. The biased empirical Fisher that uses the observed labels is not a substitute \cite{frame:Kunstner2019}.

Because $U$ and $U^{-1}$ enter as separate tensors, with $A=\partial L/\partial U$ and $B=\partial L/\partial U^{-1}$ the body gradient is $G_Y=U^\top A-BU^{-\top}$ and the spatial gradient $AU^\top-U^{-\top}B$ (from $\delta U=UY$, $\delta U^{-1}=-YU^{-1}$); both partial derivatives are therefore required. The body gradient of the stored type factor per occurrence is $F_p^{-\top}G_{Y,i}F_p^\top$, which follows from conjugating the realized body variation by $F_p^{-1}$.

The output Fisher has no general exact Kronecker factorization: its scores combine mean, covariance, inverse-frame, and write-back contributions across layers. Kronecker approximations must therefore be justified as approximations \cite{frame:MartensGrosse2015}. For example, if score matrices transform as \(G\mapsto h^\top Gh^{-\top}\), contractions using the belief metric give covariant factors \(\mathcal A=\sum GC^{-1}G^\top\) and \(\mathcal B=\sum G^\top CG\); Frobenius contractions preserve only orthogonal changes of reference basis. Inverting metric factors uses power \(-1\); matrix adaptive methods using square-root or fourth-root powers define different preconditioners \cite{frame:GuptaKorenSinger2018}. Estimating the output Fisher requires output-distribution scores or an equivalent construction, and practical factor estimates usually require persistent state. The observed-label empirical Fisher is a different object \cite{frame:Kunstner2019}. By contrast, the completed prior-family operator \(\mathcal M_{0,2}=C^{-1}\otimes(2C+\xi\xi^\top)\) is exactly a Kronecker product for that declared family and completion, not an exact factorization of the output Fisher.

\subsection{Comparison with AdamW}\label{frame:sec:adamw}

AdamW \cite{frame:KingmaBa2015,frame:LoshchilovHutter2019} divides each chart coordinate by a running root-mean-square minibatch-gradient estimate. This diagonal adaptive update is not the inverse Fisher operator of the declared Gaussian orbit family and is not generally equivariant under changes of frame basis or the type--position split. Its decoupled frame-coordinate weight decay pulls chart rows toward the identity frame. The implemented Frobenius and completed Fisher policies instead keep no Adam moments for the frame rows and override their weight decay to zero. Whether either improves training relative to a matched AdamW run is an empirical question; no such superiority follows from the geometric identities.

\subsection{Computational scope and current implementation}\label{frame:sec:implementation}
The type solve costs \(O(d^2)\) once its spatial covector is available; the pooled positional solve requires \(d^2\)-dimensional operators. Matrix exponentials, logarithms, and Fr\'echet logarithms operate on individual \(d\times d\) frame blocks. These costs are separate from the production attention computation. The earlier September 28 sizing exercise estimated roughly \(3.1\times10^9\) float64 operations per update at \(K=20\), two heads, 2,858 active type rows, and 128 positional rows using Bernoulli chart conversions. That is an arithmetic estimate for that algorithm and workload, not a timing estimate for the current Fr\'echet-logarithm implementation. No current GPU benchmark is supplied here.

The code exposes \texttt{m\_phi\_update\_mode} values \texttt{adamw}, \texttt{pullback\_group}, and \texttt{fisher\_group}. For both registered geometric policies the configuration loader currently requires untied \texttt{block\_glk}; the Frobenius policy additionally restricts head dimension to at most $12$. The two geometric policies use \texttt{GaugeManifoldAdamW} to stage only active frame rows; other parameter groups remain on AdamW. They set frame weight decay to zero and do not allocate Adam moment state for those rows. The separate selector \texttt{m\_phi\_group\_retraction} chooses \texttt{chart} or \texttt{group\_log}. The latter and the Fisher route require the built-in elementary block chart; the implemented Fisher configuration uses \texttt{block\_glk}, per-token encoding, full frame coordinates, and local diagonal layer-zero prior covariance. The general algebra results above do not imply every listed group has this optimized implementation.

\texttt{PhiFisherStatisticsProvider} reads the stored local prior moments and the current input tokens. Type rows use occurrence counts from those input tokens; positional rows pool the middle-fiber moments of each input column. Both are normalized by the scored-token count of the outer loss. \texttt{belief\_fisher\_group\_direction} recovers the spatial covector through a Fr\'echet logarithm, applies the completed type or positional solve, and converts the result to a body generator. Every frame block must lie within the principal-strip margin \(\pi-0.25\); off-strip rows are rejected rather than processed by the theoretical diagonalization alternative. Positional solve failure, zero occurrence count, or nonfinite output also rejects the row. The \texttt{group\_log} stager retains rejected rows exactly and tests reconstruction after casting the proposed coordinates to their storage precision.

The October 3, 2026 launcher selects \texttt{adamw}, \texttt{chart}, \texttt{block\_glk}, diagonal Gaussian beliefs, \texttt{projected\_predictive} attention, and no learned positional frame, with fixed right-inserted RoPE enabled. Thus the completed Fisher and group-logarithm routes exist and are selectable, but their formulas do not describe the active frame optimizer in that launcher. The proof-level distinctions between Frobenius, belief-family Fisher, and output Fisher remain valid independently of that selection.

The consolidation ran 26 existing focused checks on one CPU thread: Gaussian-KL Hessian and Fisher-pullback agreement; the completion spectrum; exact type and pooled positional solves; left, right, and split transformation laws; chart conversion and off-strip rejection; and the direct optimizer under both retractions. A separate construction check loaded the current launcher values without importing its training entry point, then constructed single-digit CPU models for AdamW, Frobenius/group-logarithm, and Fisher with both retractions. It confirmed the optimizer dispatch, statistical provider, and frame weight-decay override. These tiny checks validate formulas and code paths at the tested dimensions. They neither measure large-scale training quality nor justify claims about GPU throughput. The effect of the completion weight, occurrence normalization, noninvariant trust bounds, and trained positional conditioning remains a question for controlled experiments.

\section{Future integration and remaining obligations}
\label{unified:integration}

The intended later integration begins with Eq.~\eqref{unified:primary_link} as the primary interaction object. The low-rank exponential identities explain how selective channels arise within it; the current matrix links then provide an implementation bridge. Compression of a primitive link and the invertible routed head sum remain comparative constructions. This ordering is a theoretical choice by the author, not a conclusion drawn from small-network tests or the archived model measurements.

\subsection{Changes required in the GL(K) manuscript family}

In the main attention manuscript, the framed-state and reference-moment definitions can be retained. The matrix-link section should be extended to introduce the exponential interaction field before the selective content maps, with separate statements for invertibility, real-logarithm domains, and low-rank extraction. It should distinguish a primitive link from the weighted message operator assembled from several sources. The exact complement split becomes a receiver residual only after one specifies local complement preservation and prohibits sending the unsent component. The current full additive residual remains a separately declared architectural operation. Neither covariance nor exponentiation fixes that choice.

The main probabilistic-status theorem already supplies the predictive moment calculation. Its architecture discussion should identify covariance passthrough as the historical evaluated configuration and describe moment matching as a different supported realization. Gaussian moment projection, diagonal storage, covariance floors, and local-coordinate nonlinear blocks must each retain their own approximation or symmetry boundary. A correctly transformed full covariance obeys the passive GL(K) law, while discarding its off-diagonal entries does not commute with arbitrary frame changes.

The discussion of causal loops should retain the proof that directed causal walks cannot return through distinct positions and add the inverse-edge triangle construction for invertible links. The resulting conjugacy invariants need not vanish for the factorized content field. Their existence does not show that learned frames carry an independent curved connection, and does not choose an invariant dynamics for the links. When the field is parametrized by commuting head-diagonal factors, its nonzero defect need not exhibit noncommuting edge generators. A future nonabelian edge model should state which degrees of freedom actually vary with the edge.

The gradient companion should connect its general derivative formulas to the specific frozen receiver objective and to the current frame optimizers. Ambient alignment and channel likelihood count receiver entropy differently; exact natural-parameter evolution belongs to a frozen quadratic Gaussian objective. The pooled frame Fisher is defined by a specified statistical orbit and a declared sampling design, with stabilizer completion distinguished from the Fisher itself. A frame-table implementation does not establish that the content-map optimizer or the entire language-model training algorithm is Fisher-natural.

The supplement's fixed-joint evidence identity remains the general probabilistic reference. Its population cross-covariance formula is also the appropriate bridge to regression links and correlated receivers. The nonlinear extension already constructs source and reliability inference with information control under one objective. The present local mechanisms should be related to that construction by explicit model and recognition restrictions, rather than by substituting marginal KL scores for expected population record costs. Template mixture inference is one specified route to a nonlinear block, not a theorem that every feed-forward network optimizes that objective.

\subsection{Choices that remain open}

For Option A, the main unresolved specification is the law or objective that determines $X_{ij}$. A shared generator, a position-dependent generator, a routed generator, and a regression-derived map have different path laws and different expressive restrictions. The exact identities here support these alternatives without equating them. Regression supplies a statistical criterion when cross-agent covariances are available, but its map can be singular and may not admit a real logarithm. Turning it into an exponential group link therefore requires additional assumptions or a separately justified approximation.

Routing outside the exponent gives the stated categorical mixture over source content; routing inside it changes the single-edge operator and still requires a normalized aggregation rule. An exponential self-link can express a positive activation gate, but that rewriting does not determine the gate. A template mixture supplies a specific statistical gate and, under a different evidence interpretation, a structured Gaussian-mixture posterior. The full multimodal state, its projection, and its use by the next attention layer remain architectural and inference choices beyond the implemented local template readout.

Empirical comparisons remain separate from these derivations. A useful comparison must identify the same tokenizer, data, parameter and compute budgets, seed pairing, and evaluation protocol, and must distinguish optimizer changes from retraction, damping, and regularization changes. No training result produced in this consolidation establishes a benefit of Option A, moment-matched covariance, template routing, a new receiver objective, or frame-Fisher optimization. Tests at small dimensions ground the formulas and execution paths only. They do not establish or exclude large-population mechanisms, noncompact-group behavior, or training advantages.

The normalized autoregressive predictor and the local inference results remain valid within their stated domains. A single population ELBO for the entire composed stack still requires an explicit joint, an observation model, a recognition family, and an account of each finite update. This obligation survives the unification of the notation and should remain visible in the eventual GL(K) revision.
\appendix
\section{Source coverage and implementation provenance}
\label{unified:provenance}
The six named notes were read in full and their substantive constructions were retained. Historical workflow narration was consolidated into evidence qualifications. The source-to-section map is given below; the original files remain available beside this standalone manuscript.
\begin{longtable}{@{}p{0.29\textwidth}p{0.64\textwidth}@{}}
\toprule
Source note & Content retained and integrated \\
\midrule
\endhead
\path{link_origin.tex} & Exponential links, low-rank identities, real-logarithm chart domains, compression, head sums, complement aggregation, self-gates, and archived operator measurements (Section~\ref{unified:links}). \\
\path{link_curvature.tex} & Inverse-edge holonomy, recurrence and pure gauge, generator defects, conditional regression, reduced rank, Fisher pullback, singular-boundary scope, and archived checks (Section~\ref{unified:links}). \\
\path{projected_predictive.tex} & Encoding, comparison laws, routing, predictive mean and full covariance, diagonal projection, layer seams, decoder, and configuration interpretation (Section~\ref{unified:inference}). \\
\path{projected_gradient.tex} & Ambient and channel receiver energies, the two entropies, frozen versus joint source inference, stable natural-parameter updates, outer derivatives, and profiled self-coupling (Section~\ref{unified:inference}). \\
\path{template_routing.tex} & Mixture-prior responsibilities, MLP relation, tied and untied updates, covariance dependence, reference covariance, worked example, structured posterior, and scoped social interpretation (Section~\ref{unified:inference}). \\
\path{2026-09-28-phi-natgrad.tex} & Frobenius projection and group cases, chart recovery and retraction, statistical Fisher geometry and symmetries, batch blocks, completion, exact solvers, output Fisher, and current optimizer scope (Section~\ref{frame:sec:overview}). \\
\bottomrule
\end{longtable}

The GL(K) continuity readings covered the live \path{GL(K)_attention_v2.tex}, \path{GL(K)_supplementary_v2.tex}, \path{GL(K)_gradients.tex}, and \path{GL(K)_extension.tex}, concentrating on the constructions and qualifications identified in Sections~\ref{unified:main_bridge}, \ref{unified:companions}, and \ref{unified:integration}. These documents were used as research context, not as authority over their derivations or the author's direction.

The code snapshot began at VFE 3.0 commit \texttt{f35c573} with intentional working-tree changes. During the consolidation, an independent publication advanced that repository to \texttt{0f05ec92176391949034112a9ffcc578f88f065c}, merging the already-present routing-holonomy work. All initially recorded working-file hashes remained unchanged. Implementation descriptions therefore bind to inspected file bytes, not to a Git identifier alone. The accompanying consolidation record contains the file hashes, current configuration trace, source-section coverage, and machine-readable test summaries. This manuscript task made no VFE source or configuration changes.

The selected launcher uses projected predictive attention, moment-matched covariance, diagonal Gaussian storage, local nonlinear processing, and AdamW frame updates. Alternative receiver and frame routes are marked as such where introduced. The checks use tiny CPU models and one thread. Their dimensions were chosen to validate algebra and plumbing; no GPU benchmark, training comparison, or large-population study was performed.

\begin{thebibliography}{99}
\bibitem{unified:attention}
R. C. Dennis. \emph{Gauge-Covariant Gaussian Attention through Projected Matrix Links}. Working manuscript, 2026. \path{GL(K)_attention_v2.tex}, Research manuscript vault.

\bibitem{unified:levine2018}
S. Levine. Reinforcement learning and control as probabilistic inference: Tutorial and review. \emph{arXiv:1805.00909}, 2018. \url{https://arxiv.org/abs/1805.00909}.

\bibitem{unified:gradients}
R. C. Dennis. \emph{Variational Belief Gradients for Gauge-Covariant Gaussian Attention}. Working manuscript, 2026. \path{GL(K)_gradients.tex}.

\bibitem{unified:supplement}
R. C. Dennis. \emph{Supplementary Material for Gauge-Covariant Gaussian Attention through Projected Matrix Links}. Working manuscript, 2026. \path{GL(K)_supplementary_v2.tex}.

\bibitem{unified:extension}
R. C. Dennis. \emph{Nonlinear Inference and Information Control in Gauge-VFE Agent Networks}. Working manuscript, 2026. \path{GL(K)_extension.tex}.

\bibitem{unified:blei2017}
D. M. Blei, A. Kucukelbir, and J. D. McAuliffe. Variational inference: A review for statisticians. \emph{Journal of the American Statistical Association}, 112(518):859--877, 2017. \url{https://doi.org/10.1080/01621459.2017.1285773}.

\bibitem{links:dennis2026projpred} R.~C.~Dennis, \emph{The projected predictive route: theory and computational path}, technical note, manuscripts/Attention/projected\_predictive.tex, 2026.

\bibitem{links:wilson1974} K.~G.~Wilson, Confinement of quarks, \emph{Physical Review D} 10, 2445--2459, 1974.

\bibitem{links:kogut1979} J.~B.~Kogut, An introduction to lattice gauge theory and spin systems, \emph{Reviews of Modern Physics} 51, 659--713, 1979.

\bibitem{frame:Culver1966}
W. J. Culver, On the existence and uniqueness of the real logarithm of a matrix, \emph{Proceedings of the American Mathematical Society} \textbf{17}(5) (1966), 1146--1151.

\bibitem{links:dennis2026projgrad} R.~C.~Dennis, \emph{The projected gradient route: receiver energies, entropy, and stable updates}, technical note, manuscripts/Attention/projected\_gradient.tex, 2026.

\bibitem{links:he2016resnet} K.~He, X.~Zhang, S.~Ren, and J.~Sun, Deep residual learning for image recognition, \emph{Proceedings of the IEEE Conference on Computer Vision and Pattern Recognition}, 770--778, 2016.

\bibitem{links:elfwing2018silu} S.~Elfwing, E.~Uchibe, and K.~Doya, Sigmoid-weighted linear units for neural network function approximation in reinforcement learning, \emph{Neural Networks} 107, 3--11, 2018.

\bibitem{links:hendrycks2016gelu} D.~Hendrycks and K.~Gimpel, Gaussian error linear units (GELUs), arXiv:1606.08415, 2016.

\bibitem{links:dennis2026template} R.~C.~Dennis, \emph{Template routing}, technical note, manuscripts/Attention/template\_routing.tex, 2026.

\bibitem{links:humphreys1975} J.~E.~Humphreys, \emph{Linear Algebraic Groups}, Graduate Texts in Mathematics 21, Springer, 1975.

\bibitem{links:zhang2005schur} F.~Zhang (ed.), \emph{The Schur Complement and Its Applications}, Springer, 2005.

\bibitem{links:su2021roformer} J.~Su, Y.~Lu, S.~Pan, A.~Murtadha, B.~Wen, and Y.~Liu, RoFormer: enhanced transformer with rotary position embedding, arXiv:2104.09864, 2021.

\bibitem{links:vaswani2017} A.~Vaswani, N.~Shazeer, N.~Parmar, J.~Uszkoreit, L.~Jones, A.~N.~Gomez, {\L}.~Kaiser, and I.~Polosukhin, Attention is all you need, \emph{Advances in Neural Information Processing Systems} 30, 2017.

\bibitem{links:gu2022s4} A.~Gu, K.~Goel, and C.~R{\'e}, Efficiently modeling long sequences with structured state spaces, \emph{International Conference on Learning Representations}, 2022.

\bibitem{links:gu2023mamba} A.~Gu and T.~Dao, Mamba: linear-time sequence modeling with selective state spaces, arXiv:2312.00752, 2023.

\bibitem{frame:Hall2015}
B. C. Hall, \emph{Lie Groups, Lie Algebras, and Representations: An Elementary Introduction}, 2nd ed., Graduate Texts in Mathematics, vol.~222, Springer, Cham, 2015.

\bibitem{links:lauritzen1996} S.~L.~Lauritzen, \emph{Graphical Models}, Oxford University Press, 1996.

\bibitem{links:izenman1975} A.~J.~Izenman, Reduced-rank regression for the multivariate linear model, \emph{Journal of Multivariate Analysis} 5, 248--264, 1975.

\bibitem{links:reinsel1998} G.~C.~Reinsel and R.~P.~Velu, \emph{Multivariate Reduced-Rank Regression: Theory and Applications}, Lecture Notes in Statistics 136, Springer, 1998.

\bibitem{frame:AmariNagaoka2000}
S. Amari and H. Nagaoka, \emph{Methods of Information Geometry}, Translations of Mathematical Monographs, vol.~191, American Mathematical Society and Oxford University Press, 2000.

\bibitem[Vilnis and McCallum(2015)]{inference:vilnis2015}
L.~Vilnis and A.~McCallum.
\newblock Word representations via {G}aussian embedding.
\newblock In \emph{International Conference on Learning Representations}, 2015.

\bibitem{frame:Amari1998}
S. Amari, Natural gradient works efficiently in learning, \emph{Neural Computation} \textbf{10}(2) (1998), 251--276.

\bibitem[Plummer(2015)]{inference:plummer2015cuts}
M.~Plummer.
\newblock Cuts in {B}ayesian graphical models.
\newblock \emph{Statistics and Computing}, 25(1):37--43, 2015.

\bibitem[Jacob et~al.(2017)]{inference:jacob2017modules}
P.~E.~Jacob, L.~M.~Murray, C.~C.~Holmes, and C.~P.~Robert.
\newblock Better together? {S}tatistical learning in models made of modules.
\newblock arXiv:1708.08719, 2017.

\bibitem[Bishop(2006)]{inference:bishop2006}
C.~M.~Bishop.
\newblock \emph{Pattern Recognition and Machine Learning}.
\newblock Springer, New York, 2006.

\bibitem[Opper and Archambeau(2009)]{inference:opper2009variational}
M.~Opper and C.~Archambeau.
\newblock The variational {G}aussian approximation revisited.
\newblock \emph{Neural Computation}, 21(3):786--792, 2009.

\bibitem[Hunter and Lange(2004)]{inference:hunter2004mm}
D.~R.~Hunter and K.~Lange.
\newblock A tutorial on {MM} algorithms.
\newblock \emph{The American Statistician}, 58(1):30--37, 2004.

\bibitem[Geva et~al.(2021)]{inference:geva2021}
M.~Geva, R.~Schuster, J.~Berant, and O.~Levy.
\newblock Transformer feed-forward layers are key-value memories.
\newblock In \emph{Proceedings of the 2021 Conference on Empirical Methods in Natural Language Processing}, pages 5484--5495, 2021.

\bibitem[Sukhbaatar et~al.(2019)]{inference:sukhbaatar2019}
S.~Sukhbaatar, E.~Grave, G.~Lample, H.~J{\'e}gou, and A.~Joulin.
\newblock Augmenting self-attention with persistent memory.
\newblock arXiv:1907.01470, 2019.

\bibitem[Ramsauer et~al.(2021)]{inference:ramsauer2021}
H.~Ramsauer, B.~Sch{\"a}fl, J.~Lehner, P.~Seidl, M.~Widrich, T.~Adler, L.~Gruber, M.~Holzleitner, M.~Pavlovi{\'c}, G.~K.~Sandve, V.~Greiff, D.~Kreil, M.~Kopp, G.~Klambauer, J.~Brandstetter, and S.~Hochreiter.
\newblock Hopfield networks is all you need.
\newblock In \emph{International Conference on Learning Representations}, 2021.

\bibitem[Friston(2008)]{inference:friston2008hierarchical}
K.~Friston.
\newblock Hierarchical models in the brain.
\newblock \emph{PLoS Computational Biology}, 4(11):e1000211, 2008.

\bibitem[Hershey and Olsen(2007)]{inference:hershey2007}
J.~R.~Hershey and P.~A.~Olsen.
\newblock Approximating the {K}ullback {L}eibler divergence between {G}aussian mixture models.
\newblock In \emph{IEEE International Conference on Acoustics, Speech and Signal Processing}, pages IV-317--IV-320, 2007.

\bibitem[Ramachandran et~al.(2017)]{inference:ramachandran2017}
P.~Ramachandran, B.~Zoph, and Q.~V.~Le.
\newblock Searching for activation functions.
\newblock arXiv:1710.05941, 2017.

\bibitem[Ghahramani and Jordan(1994)]{inference:ghahramani1994}
Z.~Ghahramani and M.~I.~Jordan.
\newblock Supervised learning from incomplete data via an {EM} approach.
\newblock In \emph{Advances in Neural Information Processing Systems 6}, pages 120--127, 1994.

\bibitem[Jacobs et~al.(1991)]{inference:jacobs1991}
R.~A.~Jacobs, M.~I.~Jordan, S.~J.~Nowlan, and G.~E.~Hinton.
\newblock Adaptive mixtures of local experts.
\newblock \emph{Neural Computation}, 3(1):79--87, 1991.

\bibitem[Boyd and Richerson(1985)]{inference:boyd1985}
R.~Boyd and P.~J.~Richerson.
\newblock \emph{Culture and the Evolutionary Process}.
\newblock University of Chicago Press, Chicago, 1985.

\bibitem[Deutsch and Gerard(1955)]{inference:deutsch1955}
M.~Deutsch and H.~B.~Gerard.
\newblock A study of normative and informational social influences upon individual judgment.
\newblock \emph{Journal of Abnormal and Social Psychology}, 51(3):629--636, 1955.

\bibitem[Sherif(1936)]{inference:sherif1936}
M.~Sherif.
\newblock \emph{The Psychology of Social Norms}.
\newblock Harper \& Brothers, New York, 1936.

\bibitem[Schank and Abelson(1977)]{inference:schank1977}
R.~C.~Schank and R.~P.~Abelson.
\newblock \emph{Scripts, Plans, Goals and Understanding: An Inquiry into Human Knowledge Structures}.
\newblock Lawrence Erlbaum Associates, Hillsdale, NJ, 1977.

\bibitem[Dempster et~al.(1977)]{inference:dempster1977}
A.~P.~Dempster, N.~M.~Laird, and D.~B.~Rubin.
\newblock Maximum likelihood from incomplete data via the {EM} algorithm.
\newblock \emph{Journal of the Royal Statistical Society, Series B}, 39(1):1--38, 1977.

\bibitem[Henrich and Boyd(1998)]{inference:henrich1998}
J.~Henrich and R.~Boyd.
\newblock The evolution of conformist transmission and the emergence of between-group differences.
\newblock \emph{Evolution and Human Behavior}, 19(4):215--241, 1998.

\bibitem{frame:Absil2008}
P.-A. Absil, R. Mahony, and R. Sepulchre, \emph{Optimization Algorithms on Matrix Manifolds}, Princeton University Press, Princeton, 2008.

\bibitem{frame:Higham2008}
N. J. Higham, \emph{Functions of Matrices: Theory and Computation}, SIAM, Philadelphia, 2008.

\bibitem{frame:Mathias1996}
R. Mathias, A chain rule for matrix functions and applications, \emph{SIAM Journal on Matrix Analysis and Applications} \textbf{17}(3) (1996), 610--620.

\bibitem{frame:DenmanBeavers1976}
E. D. Denman and A. N. Beavers, The matrix sign function and computations in systems, \emph{Applied Mathematics and Computation} \textbf{2}(1) (1976), 63--94.

\bibitem{frame:Milnor1976}
J. Milnor, Curvatures of left invariant metrics on Lie groups, \emph{Advances in Mathematics} \textbf{21}(3) (1976), 293--329.

\bibitem{frame:NealHinton1998}
R. M. Neal and G. E. Hinton, A view of the EM algorithm that justifies incremental, sparse, and other variants, in M. I. Jordan (ed.), \emph{Learning in Graphical Models}, Kluwer, Dordrecht, 1998, 355--371.

\bibitem{frame:Skovgaard1984}
L. T. Skovgaard, A Riemannian geometry of the multivariate normal model, \emph{Scandinavian Journal of Statistics} \textbf{11}(4) (1984), 211--223.

\bibitem{frame:Williamson1936}
J. Williamson, On the algebraic problem concerning the normal forms of linear dynamical systems, \emph{American Journal of Mathematics} \textbf{58}(1) (1936), 141--163.

\bibitem{frame:Cencov1982}
N. N. \v{C}encov, \emph{Statistical Decision Rules and Optimal Inference}, Translations of Mathematical Monographs, vol.~53, American Mathematical Society, Providence, 1982.

\bibitem{frame:ShermanMorrison1950}
J. Sherman and W. J. Morrison, Adjustment of an inverse matrix corresponding to a change in one element of a given matrix, \emph{Annals of Mathematical Statistics} \textbf{21}(1) (1950), 124--127.

\bibitem{frame:VanLoanPitsianis1993}
C. F. Van Loan and N. Pitsianis, Approximation with Kronecker products, in M. S. Moonen, G. H. Golub, and B. L. R. De Moor (eds.), \emph{Linear Algebra for Large Scale and Real-Time Applications}, Kluwer, Dordrecht, 1993, 293--314.

\bibitem{frame:Kunstner2019}
F. Kunstner, P. Hennig, and L. Balles, Limitations of the empirical Fisher approximation for natural gradient descent, in \emph{Advances in Neural Information Processing Systems} \textbf{32} (2019).

\bibitem{frame:MartensGrosse2015}
J. Martens and R. Grosse, Optimizing neural networks with Kronecker-factored approximate curvature, in \emph{Proceedings of the 32nd International Conference on Machine Learning}, PMLR \textbf{37} (2015), 2408--2417.

\bibitem{frame:GuptaKorenSinger2018}
V. Gupta, T. Koren, and Y. Singer, Shampoo: preconditioned stochastic tensor optimization, in \emph{Proceedings of the 35th International Conference on Machine Learning}, PMLR \textbf{80} (2018), 1842--1850.

\bibitem{frame:KingmaBa2015}
D. P. Kingma and J. Ba, Adam: a method for stochastic optimization, in \emph{International Conference on Learning Representations}, 2015.

\bibitem{frame:LoshchilovHutter2019}
I. Loshchilov and F. Hutter, Decoupled weight decay regularization, in \emph{International Conference on Learning Representations}, 2019.
\end{thebibliography}
\end{document}
